%% Direct comparison of SL(2) with its metaplectic cover -- David Joyner and Claude
%% (formerly: Multiplicity one for SL(2))
%% Draft 7 (2026) of the 1995 Scientific Word manuscript aabook.tex and its
%% \input files, assembled into a single file.  See ../revision_log.md.
\documentclass[12pt,letterpaper,openany]{book}
\usepackage[american]{babel}
\usepackage[T1]{fontenc}
\usepackage{amsmath,amssymb,amsthm}
\usepackage{mathptmx}          % Times text and math
\usepackage{courier}           % Courier for typewriter text
\usepackage[letterpaper,margin=1in]{geometry}
\usepackage{microtype}

% --- theorem environments (formerly from the thmsa option) ---
\theoremstyle{plain}
\newtheorem{theorem}{Theorem}[chapter]
\newtheorem{lemma}[theorem]{Lemma}
\newtheorem{proposition}[theorem]{Proposition}
\newtheorem{corollary}[theorem]{Corollary}
\newtheorem{conjecture}[theorem]{Conjecture}
\newtheorem{hypothesis}[theorem]{Hypothesis}
\theoremstyle{definition}
\newtheorem{definition}[theorem]{Definition}
\newtheorem{example}[theorem]{Example}
\theoremstyle{remark}
\newtheorem{remark}[theorem]{Remark}
\theoremstyle{remark}
\newtheorem{remyy}{Remark}
\renewcommand{\theremyy}{2026.\arabic{remyy}}

\newcommand{\tr}{\operatorname{tr}}
\DeclareMathOperator{\sgn}{sgn}

% --- the book class has no abstract environment ---
\newenvironment{abstract}%
  {\begin{center}\bfseries Abstract\end{center}\begin{quote}}%
  {\end{quote}}

\usepackage[hidelinks]{hyperref}

\begin{document}

\author{David Joyner\\[0.5ex] \normalsize\texttt{wdjoyner@gmail.com}\thanks{This is a corrected and expanded
version of the first author's unpublished manuscript \emph{Multiplicity one for
$SL(2)$}, written in Scientific Word (files \texttt{aabook.tex},
\texttt{aabook1}--\texttt{aabook6}, \texttt{aabookrf.tex}, dated April 9,
1995). The title was changed in Draft~5.
This revision is Draft~7, November 2026.}
\and Claude\\[0.5ex] \normalsize Anthropic}
\title{Direct comparison of $SL(2)$ with its metaplectic cover\\[2ex]
\large Draft 7\\[0.5ex] Corrected and expanded version}
\date{Original draft: 1995. Corrected and expanded draft: 2026}
\maketitle

\begin{abstract}
This revision was prepared with the help of AI systems, whose roles are
described in the Acknowledgments at the end of the text.
By the work of Weil, Shimura, Waldspurger and others, genuine cuspidal
automorphic representations of the 2-fold metaplectic cover
$\overline{SL(2,\mathbb{A})}$ correspond to cuspidal automorphic
representations of $PGL(2,\mathbb{A})$, where $\mathbb{A}$ is the ring of
adeles of a number field. The cover lies over $SL(2)$, not over $PGL(2)$,
and its orbital integrals can be compared directly with those of $SL(2)$.
We ask whether this direct comparison gives a correspondence with
representations of $SL(2,\mathbb{A})$ itself, and, if not, why not and how
far it is from the known one. Locally the direct comparison works on the
split torus: for odd residual characteristic and for the real field we
describe the genuine characters of the cover of the torus, match functions
so that their Fourier transforms agree, and deduce a correspondence of
principal series and an isomorphism of spherical Hecke algebras. Globally
it does not: for every cuspidal representation $\pi$ of $SL(2,\mathbb{A})$,
no genuine cuspidal representation of the cover corresponds to $\pi$ under
this isomorphism at almost all places (Theorem~\ref{thm:nonexist}). The
direct comparison matches unramified parameters up to sign, whereas the
correspondence of Shimura and Waldspurger, followed by restriction from
$PGL(2)$ to $SL(2)$, squares them. So the direct comparison is off by a
square root. As a consequence, a comparison of trace formulas proposed in
the 1995 draft of this text, which would have deduced multiplicity one for
$SL(2)$ from Waldspurger's multiplicity one theorem for the cover, cannot
hold (Proposition~\ref{thm:main}, Corollary~\ref{cor:inconsistent}). We
also construct three matchings on $SL(2,\mathbb{R})$ that transfer genuine
discrete series to discrete series or limits of discrete series; these are
of the squaring type. Tracing the proposed comparison through the trace
formula, we find that it is consistent for the split torus, and fails for
the terms that involve the root: the intertwining operators, the residual
spectrum and the cuspidal spectrum. A last chapter reads the comparison
through the trace formula of the cover, which a companion text writes
out. The factor in the matching on the torus is the function that
defines stable conjugacy on the cover, the direct comparison is the
known correspondence with $PGL(2)$ with the tori identified in the wrong
way, and on the elliptic set it fails at the elements close to $-1$,
where the unit elements of the Hecke algebras have their stable and
unstable orbital integrals interchanged. Multiplicity one for $SL(2)$ is a theorem of
Ramakrishnan. Statements that have not been fully verified are marked as
``Remark 2026.$n$''.
\end{abstract}

\tableofcontents

\chapter{Introduction}\label{ch:intro}

The goal of this text is to use the trace formula to study the
multiplicity one problem for cuspidal automorphic representations of
$SL(2)$, by deducing it from the multiplicity one result for
the 2-fold metaplectic cover of $SL(2)$, where it is a theorem
of Waldspurger. The motivation for this problem
has its origins in several stimulating conversations
with R. Langlands in the late 1980s.

\section{Motivation}\label{sec:motivation}

This section explains the question behind the text.

The genuine automorphic forms on the 2-fold cover of $SL(2)$ are the
modular forms of half-integral weight. Weil's representation \cite{We3},
and the theta series built from it, connect them with quadratic forms.
Shimura \cite{Shi} attached to a cusp form of weight $k/2$, $k$ odd, a
modular form of even weight $k-1$; his proof used Rankin's method and
Weil's converse theorem. Waldspurger \cite{Wa1}, \cite{Wa2} put this
correspondence in representation-theoretic form and determined its
fibers: the genuine cuspidal automorphic representations of the cover
that are orthogonal to the elementary theta series correspond to cuspidal
automorphic representations of $PGL(2)$. So the linear group known to be
paired with the cover of $SL(2)$ is $PGL(2)$. In one sense $PGL(2)$ is
smaller than $SL(2)$: its representations restrict to $SL(2)$, and the
restriction forgets the difference between a representation and its
quadratic twists. On dual groups, the correspondence with $PGL(2)$ is the
identity of $SL(2,{\mathbb C})$, and restriction to $SL(2)$ is the
projection $SL(2,{\mathbb C})\rightarrow PGL(2,{\mathbb C})$.

The cover is, however, a cover of $SL(2)$ and not of $PGL(2)$. Its tori,
its conjugacy classes and its orbital integrals lie over those of
$SL(2)$, so functions and characters on the two groups can be compared
directly, element by element. This suggests the following question.

\begin{quote}
\emph{Does the direct comparison give a correspondence between genuine
representations of the cover and representations of $SL(2)$ itself? If
not, why is it wrong, and how far is it from the correspondence with
$PGL(2)$?}
\end{quote}

The 1995 draft of this text took the answer to the first question to be
yes, and tried to use the resulting correspondence, together with
Waldspurger's multiplicity one theorem, to prove multiplicity one for
$SL(2)$ by a comparison of trace formulas. This revision gives the
following partial answer.

\begin{itemize}
\item \emph{Locally, on the split torus, the direct comparison works.}
There is a bijection between characters of the torus and genuine
characters of its cover, and a matching of functions that preserves
Fourier transforms. It gives an identity of traces of principal series and
an isomorphism of spherical Hecke algebras (\S\ref{sec:charsA},
\S\ref{sec:spherical}).
\item \emph{Globally it is wrong.} No genuine cuspidal automorphic
representation is related to a cuspidal automorphic representation of
$SL(2)$ by this isomorphism at almost all places
(Theorem~\ref{thm:nonexist}).
\item \emph{It is off by a square root.} Let $t'$ be the unramified
parameter of a genuine representation. The correspondence of Shimura and
Waldspurger, followed by restriction to $SL(2)$, assigns to it the
representation of $SL(2,F)$ with parameter $t=t'^{\,2}$. The direct
comparison assigns the one with parameter $t=\pm t'$
(Lemma~\ref{lem:square}), a square root of the right parameter up to
sign. In terms of dual groups, the direct comparison would need a
homomorphism from the torus of $SL(2,{\mathbb C})$ to that of
$PGL(2,{\mathbb C})$ which is an isomorphism, and the natural map is of
degree two. The discrepancy is already visible locally: the direct
comparison does not preserve unitarity or reducibility
(Remark~\ref{rem:choices}), and at the real place the matchings that do
transfer discrete series double the infinitesimal character
(Lemma~\ref{lem:realds}).
\end{itemize}

We do not know whether the direct comparison has a
representation-theoretic meaning beyond these identities, nor what a
comparison of trace formulas built on the squaring correspondence would
say about $SL(2)$. These questions are listed in \S\ref{sec:summary} and
discussed in Remark~\ref{rem:even}. \S\ref{sec:where-fails} traces the
proposed comparison through the trace formula and shows which terms match
and which do not: the direct comparison is globally consistent on the
torus, and fails for the intertwining operators, the residual spectrum and
the cuspidal spectrum. Chapter~\ref{ch:cover} takes both points further,
using the trace formula of the cover as written out in the companion
text \cite{JC}. It gives the direct comparison a meaning in terms of
stable conjugacy, and it shows that the comparison fails on the elliptic
terms of the geometric side as well.

\begin{example}
The cusp form $\Delta$ of weight $12$ corresponds, under Shimura's
correspondence, to a cusp form of weight $13/2$. If
$\tau(p)=p^{11/2}(\alpha_p+\alpha_p^{-1})$, the parameter of the genuine
form at an odd prime $p$ is $t'_p=\pm\alpha_p$, and the parameter of the
restriction of $\Delta$ to $SL(2)$ is $t_p=\alpha_p^2=t_p'^{\,2}$. The
direct comparison would instead pair the restriction of $\Delta$ with a
genuine form of parameters $\pm\alpha_p^2$, and
Theorem~\ref{thm:nonexist} says that there is none.
\end{example}

\section{Status of the main result}

This section states what is known about multiplicity one for $SL(2)$ and
what this text proves.

Multiplicity one for the cuspidal spectrum of $SL(2)$ was conjectured
by Labesse and Langlands \cite{LL}, \cite{Lab}, and was proved by
Ramakrishnan \cite{Ra}. D. Blasius has shown that the analog of the
multiplicity one result is false for $SL(n)$, $n>2$ \cite{Blas}.
Strong multiplicity one fails for $SL(2)$ \cite{LL}; Chai and Zhang
\cite{CZ} (see also \cite{Zh}) proved a qualified form of it, for pairs
of cuspidal representations which also agree at the archimedean places
and at the places above $2$, and are generic with respect to the same
additive character. We stress that the following four
properties are different and are not used interchangeably here:
multiplicity one (each cuspidal representation occurs at most once),
strong multiplicity one (a cuspidal representation is determined by
almost all of its local components), the local theta correspondence,
and the description of global near-equivalence classes (``Waldspurger
packets'') on the metaplectic cover \cite{Wa1}, \cite{Wa2}, \cite{PS}.

The 1995 draft of this text aimed to deduce multiplicity one for $SL(2)$
from Waldspurger's theorem, by comparing the trace formula of $SL(2)$ with
that of its cover through the direct comparison of \S\ref{sec:motivation}.
This revision shows that the comparison cannot hold in
the form that was proposed, and explains why. The history of the argument, and the errors
found in earlier drafts, are recorded in the revision log (see the
Acknowledgments).

The results are of three kinds.

\emph{Local results} (Chapter~\ref{ch:local}).
\begin{itemize}
\item For $p$ odd and for ${\mathbb R}$, we describe the genuine characters
of the cover of the split torus and fix a bijection
$\chi\mapsto\overline\chi$ with the characters of the torus
(Lemma~\ref{lem:genuine-chars}, Definition~\ref{def:corrchar}). It commutes
with the Weyl group (Lemma~\ref{lem:Wequiv}) and depends on a choice
(Remark~\ref{rem:choices}).
\item We match functions on the torus and on its cover so that their
Fourier transforms agree (Lemmas~\ref{lem:FT-match}
and~\ref{lem:realA-match}). This gives an identity of traces of principal
series (Proposition~\ref{prop:hyp-ps}) and an isomorphism of spherical
Hecke algebras, the \emph{Satake correspondence} of
\S\ref{sec:spherical}. These matchings compare orbital integrals at an
element of the torus and at a lift of the same element. We call them
\emph{odd}.
\item On $SL(2,{\mathbb R})$ we construct three matchings of functions
with vanishing hyperbolic orbital integrals, which transfer genuine
discrete series to discrete series or limits of discrete series
(Lemma~\ref{lem:realds}). They compare a rotation $r(\theta)$ with the
element of the cover over $r(2\theta)$. We call matchings of this kind
\emph{even}.
\end{itemize}
No matching is established at the dyadic places.

\emph{A reading through the trace formula of the cover}
(Chapter~\ref{ch:cover}). The factor in the odd matching is the function
$\Delta_\psi$ that defines stable conjugacy on the cover
(Lemma~\ref{lem:Delta}). On the split torus the direct comparison is the
lifting of characters between the cover and $PGL(2)$, with the tori of
$PGL(2)$ and $SL(2)$ identified by a map that is not a homomorphism of
groups (Proposition~\ref{prop:wrong-orbit}). On the elliptic set the unit
elements of the Hecke algebras do not match
(Proposition~\ref{prop:unit-swap}, Corollary~\ref{cor:no-fl}). Most of
these statements are proved in \cite{JC}.

\emph{A conditional statement} (Proposition~\ref{thm:main}).
Fix a cuspidal automorphic representation $\pi=\otimes_v\pi_v$ of
$SL(2,\mathbb{A})$, and let $S=S(\pi)$ be the finite set consisting of the
archimedean places, the dyadic places, and the odd places $v$ at which
$\pi_v$ is not spherical with respect to $K_{0,v}=SL(2,{\cal O}_v)$.
Assume:
\begin{itemize}
\item (Hypothesis~\ref{hyp:GH}) for global functions $f$ on
$SL(2,\mathbb{A})$ and genuine global functions $\overline{f}$ on the cover
which correspond at every place (by the Satake correspondence outside
$S$, and through prescribed local classes of pairs at the places of
$S$), the traces of the cuspidal spectra agree after restriction to
each archimedean infinitesimal character $\nu$;
\item (Hypothesis~\ref{hyp:LH}) at each $v\in S$, the local class of
pairs separates characters, detects $\pi_v$, and transfers each relevant
genuine irreducible representation to at most one irreducible
representation of $SL(2,k_v)$.
\end{itemize}
Then $\pi$ occurs with multiplicity one, and $\pi$ corresponds to a
unique genuine cuspidal automorphic representation $\overline{\pi}$ of
the cover. The local hypothesis holds at odd places where
$\pi_v$ is spherical with respect to a hyperspecial maximal compact
subgroup (Proposition~\ref{prop:LH-spherical}).

\emph{The main theorem} (Theorem~\ref{thm:nonexist}). For every cuspidal
$\pi$ and every finite $S\supseteq S(\pi)$, there is no genuine cuspidal
automorphic representation of the cover whose components outside $S$
correspond to those of $\pi$ under the Satake correspondence. Consequently
Hypotheses~\ref{hyp:LH} and~\ref{hyp:GH} cannot both hold, for any $\pi$
(Corollary~\ref{cor:inconsistent}). The conditional statement is
therefore vacuous, and the proposed trace identity is false whenever the
local hypothesis can be satisfied.

The reason is the following. Let $t$ and $t'$ be the unramified parameters
of a representation of $SL(2,F)$ and of a genuine representation of the
cover. The Satake correspondence is $t=\pm t'$ (Lemma~\ref{lem:square}).
The correspondence realized by automorphic forms, through Waldspurger's
theorem and restriction from $PGL(2)$ to $SL(2)$, is $t=t'^{\,2}$. The
proof of Theorem~\ref{thm:nonexist} turns this into a contradiction between
the poles of two $L$-functions. It uses the Labesse--Langlands description
of the cuspidal spectrum of $SL(2)$, the Gelbart--Jacquet adjoint lift and
bound toward the Ramanujan conjecture, Waldspurger's correspondence, and
the Rankin--Selberg theory of Jacquet and Shalika.

A correct comparison would have to use even matchings at every place.
Remark~\ref{rem:even} explains why, even then, we do not see how to deduce
multiplicity one for $SL(2)$ from Waldspurger's multiplicity one theorem:
such a comparison passes through $PGL(2)$ and controls sums over
$L$-packets, not the multiplicity of a single representation.

\begin{example}
Let $k={\mathbb Q}$ and let $\pi$ be a cuspidal automorphic representation of
$SL(2,{\mathbb A})$ in the restriction of the representation $\Pi$ of
$GL(2,{\mathbb A})$ generated by the cusp form $\Delta$ of weight $12$. Since
$\Delta$ has level one, $\Pi_p$ is unramified at every prime $p$. If
$\tau(p)\neq0$, the restriction of $\Pi_p$ to $SL(2,{\mathbb Q}_p)$ is
irreducible, so $\pi_p$ has a vector fixed by $SL(2,{\mathbb Z}_p)$. Hence
$S(\pi)$ consists of the real place, the place $2$, and possibly some odd
primes $p$ with $\tau(p)=0$. By Lehmer's conjecture there are no such
primes. The genuine cusp form attached to $\Delta$ by the Shimura
correspondence has weight $13/2$, and its parameters are related to those
of $\pi$ by $t_p=t_p'^{\,2}$. We use this $\pi$ as a running example.
\end{example}

\section{Literature review}\label{sec:litreview}

This section summarizes the work that this text builds on, and the later
work that bears on its main result. It is divided at 1995, the date of the
original manuscript.

\subsection*{Work before 1995}

Kubota \cite{K1}, \cite{K2} constructed the 2-fold cover of $SL(2)$ over
local fields and over the adeles by an explicit cocycle, and showed that
the adelic cover splits over $SL(2,k)$. Moore \cite{M} determined the
topological central extensions of p-adic linear groups. These are the
groups studied here.

Gelbart \cite{Ge2} and Gelbart and Sally \cite{GS} described the Weil
representation and the genuine spectrum of the metaplectic group,
including the complementary series and the reducible principal series.
Moen \cite{Mo1}, \cite{Mo2} proved that genuine unitary principal series
are irreducible for $p>2$ and computed the intertwining operators.
Manderscheid \cite{M1}, \cite{M2} parametrized the supercuspidal
representations of $SL(2)$ and of its cover by compact induction. Appendix~\ref{ch:sc}
uses his parametrization.

Waldspurger \cite{Wa1} related genuine cusp forms on the cover to cusp
forms on $PGL(2)$ through the theta correspondence. He proved that the
space ${\cal A}_{00}$ of genuine cusp forms orthogonal to the elementary
theta series has multiplicity one. In \cite{Wa2} he described the
near-equivalence classes of genuine cuspidal representations using
quaternion algebras. Piatetski-Shapiro's survey \cite{PS} summarizes this
work. Gelbart and Piatetski-Shapiro studied the elementary theta series
and the distinguished representations \cite{GePS}, and compared the
correspondences of Shimura and Waldspurger \cite{GPS}.

Labesse and Langlands \cite{LL}, \cite{Lab} analyzed the stable trace
formula for $SL(2)$, introduced $L$-packets, and conjectured that the
cuspidal spectrum of $SL(2)$ has multiplicity one. Blasius \cite{Blas}
showed that the analogous statement fails for $SL(n)$, $n>2$.

The method of this text is a comparison of trace formulas. Its models are
the comparison of $GL(2)$ with a quaternion algebra by Jacquet and
Langlands \cite{JL} and by Gelbart and Jacquet \cite{GJ}, Langlands' base
change \cite{L4}, and above all Flicker's comparison of $GL(2)$ with its
covers \cite{Fl}, whose decomposition-of-measures argument was the intended
route to a trace identity such as Hypothesis~\ref{hyp:GH}. The trace formula used is Arthur's
\cite{A2}, \cite{A3}, \cite{A6}, \cite{A7}. Arthur's Paley--Wiener theorem
\cite{APW} provides the multipliers mentioned in
Remark~\ref{rem:why-refine}. Gelbart and Jacquet \cite{GJ78} constructed
the symmetric square lift from $GL(2)$ to $GL(3)$, which gives the bound
toward the Ramanujan conjecture used in Lemma~\ref{lem:A00}. Savin
\cite{Sav} studied the genuine Iwahori--Hecke algebras of metaplectic
covers.

\subsection*{Work after 1995}

Adams \cite{Ad2} defined a lifting of characters between the real
metaplectic groups $Mp(2n,{\mathbb R})$ and odd orthogonal groups. It is
the published form of the 1995 preprint used for the real character
identities of Chapter~\ref{ch:local}.

Ramakrishnan \cite{Ra} proved multiplicity one for the cuspidal spectrum
of $SL(2)$. He deduced it from his proof that the Rankin--Selberg product
of two cuspidal representations of $GL(2)$ is automorphic on $GL(4)$,
together with the results of Labesse and Langlands. This settles the
conjecture that the present text approaches by a different method.

Chai and Zhang \cite{CZ} proved a qualified strong multiplicity one
theorem for $SL(2)$. Two cuspidal representations that agree at almost all
places, and also at the archimedean places and the places above $2$, and
that are generic with respect to the same additive character, are equal. Their proof uses a local converse theorem. Zhang \cite{Zh} removed
the restriction on the residue characteristic in that local converse
theorem.

Takeda and Wood \cite{TW} established Hecke algebra correspondences at
Iwahori level for the metaplectic group, including the case $p=2$. They
pair the cover of $SL(2)=Sp(2)$ with $SO(3)\cong PGL(2)$ rather than with
$SL(2)$. This bears on Remark~\ref{claim:savin} and on the dyadic places
(\S\ref{sec:dyadic}).

Gan and Savin \cite{GS12} established, for odd residual characteristic,
the Hecke algebra correspondences that Takeda and Wood extend.

Li proved the basic results of local harmonic analysis on covering groups,
including the existence of characters as locally integrable functions
\cite{Li12}, established local character relations for the metaplectic
group \cite{Li19}, and stabilized its trace formula \cite{Li21}. The
elliptic endoscopic groups of the cover of $Sp(2n)$ are products of split
odd orthogonal groups. For $n=1$ the comparison is with
$SO(3)\cong PGL(2)$. This agrees with Waldspurger's description and with
Theorem~\ref{thm:nonexist}: the linear group naturally paired with the
cover of $SL(2)$ is $PGL(2)$, not $SL(2)$.

The companion text \cite{JC} is the 2026 revision of the first author's
notes of 1995 on the trace formula of the cover of $SL(2)$. It writes
out Arthur's trace formula for the 2-fold cover, in non-invariant and
invariant form, with the constants determined, and computes the action
of $GL(2)$ on the cover and on each term. It proves several statements
that the present text quotes (\S\ref{sec:cover-supplies}), and its last
chapter treats the direct comparison of the present text.
Chapter~\ref{ch:cover} is based on it.

\section{Conventions on the status of statements}

This section explains how the status of each statement is shown. Every
result in the text is of one of four kinds: proved here, quoted from the
literature, conditional on stated hypotheses, or conjectural.

Statements labeled Theorem, Proposition, Lemma, or Corollary are either
proved here or quoted, with a reference, from the literature. Anything that
has not been fully verified is labeled \textbf{Remark 2026.$n$}. This
includes statements of the 1995 draft that have not been checked or are
known to need correction, and the qualification is attached to the
statement itself. Statements labeled Conjecture or Hypothesis are
assumptions, and are named as such wherever they are used.

Several explicit local computations (Lemmas~\ref{lem:Acocycle},
\ref{lem:genuine-chars}, \ref{lem:realA-match} and~\ref{lem:realds}, and
some of the examples) are checked by exact computation. The script is
\[
\texttt{mult-one/scripts/local\_formulas\_check.py}
\]
in the repository \texttt{github.com/wdjoyner/trace-formulas}, and it
lists what it checks. A second script, \texttt{referee\_check.py} in the same directory, checks Lemma~\ref{lem:facts}(h), the splitting
(\ref{eq:sv}), and the example in Remark~\ref{rem:sc-caveat}. These checks supplement the proofs and do not replace them.

\begin{example}
The supercuspidal character relation of the 1995 draft appears as
Remark~\ref{prop:164}, with the qualification ``needs correction'' in its
heading. It is therefore not used in Chapter~\ref{ch:global}, whose results use
only proved or quoted statements (and, in Proposition~\ref{thm:main}, the
two Hypotheses~\ref{hyp:GH} and~\ref{hyp:LH}).
\end{example}

\section{Outline}

Chapter~\ref{ch:background} recalls the metaplectic covers of $SL(2)$,
locally and globally. Chapter~\ref{ch:reps} surveys the representation
theory of the local cover.
Chapter~\ref{ch:local} presents local correspondences between $SL(2,F)$
and its 2-fold cover. They are divided into three types according to the
residual characteristic of the field: $p=\infty$ (real), $p=2$,
and $p>2$. For $p>2$ and $p=\infty$, the idea is to show that if there
is a certain type of identity between the orbital integral of a smooth
function $f$ on $SL(2,F)$ and the orbital integral of a smooth function
$\overline{f}$ on the two-fold cover $\overline{SL(2,F)}$ (call this a
``correspondence'' or ``matching'' of functions), then there is an
analogous identity between the trace maps
$f^\wedge:\pi\longmapsto \tr(\pi(f))$
and $\overline{f}^\wedge:\overline{\pi}\longmapsto
\tr(\overline{\pi}(\overline{f}))$
on the unitary duals (call this a ``lifting'' or ``transfer''
of genuine virtual representations of $\overline{SL(2,F)}$ to
virtual representations of $SL(2,F)$).
In \cite{L4} and \cite{Fl}, only two types
of correspondences or matchings were needed: the
``Satake correspondence'' of spherical functions and the
``smooth matching'' of Schwartz-Bruhat functions.
In the case of $SL(2)$, some new
types of ``matchings'' are introduced. No matching is established at
$p=2$ (\S\ref{sec:dyadic}).

Chapter~\ref{ch:global} contains the global results: the two hypotheses,
their formal consequence (Proposition~\ref{thm:main}), the main
theorem (Theorem~\ref{thm:nonexist}), and an account of which terms of the
trace formulas match under the direct comparison
(\S\ref{sec:where-fails}).

Chapter~\ref{ch:cover} reads the comparison through the trace formula of
the cover of \cite{JC}: stable conjugacy, the elliptic terms, and a table
of the terms that match and the terms that do not. It is not used in the
earlier chapters.

Appendix~\ref{ch:sc} discusses a proposed correspondence between
supercuspidal representations for $p>2$. Its main statement is false when
$q\equiv3\pmod4$. Appendix~\ref{ch:tf} summarizes the trace formulas of
$SL(2)$ and of its cover. Neither appendix is used in the proofs.

Chapter~\ref{ch:global} depends on Chapter~\ref{ch:background}, on
\S\ref{sec:charsA}--\S\ref{sec:spherical}, and on formula (iii) of
\S\ref{sec:weil}. The lemmas on non-tempered representations in
\S\ref{sec:weil} and Lemma~\ref{lem:realds} are used only in
Remark~\ref{rem:LH-families} and Remark~\ref{rem:even}.

\begin{example}
The Satake correspondence of \S\ref{sec:spherical} is an odd matching: the
unit elements $1_K$ and $e_K$ of the two spherical Hecke algebras
correspond, and their traces agree on corresponding unramified principal
series. The matching of Lemma~\ref{lem:realds}(a) is even: it relates the
orbital integral of $f$ at $r(\theta)$ to that of $\overline f$ at the
element over $r(2\theta)$.
\end{example}

\section{Notation}\label{sec:notation}

The following symbols are used throughout. Local notation refers to a
local field $F$; global notation refers to a number field $k$.

\medskip
\begin{center}
\renewcommand{\arraystretch}{1.15}
\begin{tabular}{p{0.27\textwidth}p{0.66\textwidth}}
\hline
$G=SL(2)$, $G(F)$ & the special linear group and its $F$-points\\
$\overline{G(F)}$, $\overline{H}$ & the 2-fold metaplectic cover, and the
inverse image of a subgroup $H\subset G(F)$\\
$(g,\varsigma)$ & an element of $\overline{G(F)}$, $g\in G(F)$,
$\varsigma\in\mu_2=\{\pm1\}$\\
$\beta$, $(x,y)_2$ & the Kubota cocycle, and the quadratic Hilbert
symbol\\
$p$, $q$, ${\cal O}$, $\varpi$ & residual characteristic ($p=\infty$ if
$F$ is archimedean), order of the residue field, ring of integers, and a
uniformizer of $F$ \\
$\nu$, $|\cdot|$ & normalized valuation and absolute value,
$|\varpi|=q^{-1}$\\
$B=AN$ & the upper triangular Borel subgroup, its diagonal torus $A$ and
unipotent radical $N$\\
$a(x)$, $h(x)$ & the diagonal element $\operatorname{diag}(x,x^{-1})$\\
$T$, ${\cal T}$ & a Cartan subgroup of $G$; a set of representatives of
the $G(F)$-conjugacy classes of Cartan subgroups\\
$K_0$, $K_1$ & the hyperspecial maximal compact subgroups $G({\cal O})$
and $w_\varpi^{-1}G({\cal O})w_\varpi$ ($p$ odd)\\
$w$, $w_\varpi$ & the matrices $\left(\begin{smallmatrix}0&-1\\1&0\end{smallmatrix}\right)$ and
$\left(\begin{smallmatrix}0&\varpi\\1&0\end{smallmatrix}\right)$\\
$\pi$, $\overline{\pi}$ & a representation of $G$, and a genuine
representation of the cover\\
$\Theta_\pi$, $\tr\pi(f)$ & the character of $\pi$, and the trace of
$\pi(f)$\\
$\chi$, $\overline{\chi}$ & a character of $A(F)$, and a genuine
character of $\overline{A(F)}$\\
$\Pi(H)$, $\Pi_u(H)$ & irreducible admissible (genuine, for a cover)
representations of $H$, and the unitary ones (\S\ref{sec:chars})\\
$I_A(\chi)$ & the unitarily induced principal series\\
$O_f$, $F_f$, $\Phi_f$ & orbital integral, normalized hyperbolic orbital
integral, normalized elliptic orbital integral of $f$\\
$S$ & a finite set of places of $k$ (only in Chapter~\ref{ch:global}
and Appendix~\ref{ch:tf})\\
$S(\pi)$ & the exceptional set of places of
Definition~\ref{def:exceptional}\\
${\cal C}_v$ & a class of corresponding pairs at a place $v\in S$\\
$t_v$ & the local transfer of genuine representations in
Hypothesis~\ref{hyp:LH}\\
$\nu_\infty(\pi)$ & the archimedean infinitesimal character of $\pi$
(\S\ref{sec:hypotheses})\\
$\mathfrak{F}_v$, $\overline{\mathfrak{F}}_v$ & local components of cuspidal
representations, and of representations in ${\cal A}_{00}$
(\S\ref{sec:hypotheses})\\
$t$, $t'$ & unramified parameters on $SL(2,F)$ and on the cover
(\S\ref{sec:nonexist})\\
\hline
\end{tabular}
\end{center}


\section{A running example}\label{sec:running}

This section collects one example that is followed through the text:
the field ${\mathbb Q}$, the set of places $S=\{\infty,2\}$, and the cusp
form $\Delta$. It is the smallest case in which every feature of the
comparison appears. The same field and the same $S$ are the running
example of the companion text \cite[\S1.3]{JC}, where they are followed
through the trace formula of the cover. The items below are in the order
of the chapters.

\begin{itemize}
\item \emph{The data.} $k={\mathbb Q}$ and $S=\{\infty,2\}$. The additive
character $\psi$ of ${\mathbb Q}\backslash{\mathbb A}$ has
$\psi_\infty(x)=e^{-2\pi ix}$ and $\psi_p(x)=e^{2\pi i\{x\}_p}$, where
$\{x\}_p$ is the $p$-adic fractional part. The representation $\pi$ of
$SL(2,{\mathbb A})$ is a cuspidal constituent of the restriction of the
representation of $GL(2,{\mathbb A})$ generated by $\Delta$, the cusp
form of weight $12$ and level one, with
$\tau(p)=p^{11/2}(\alpha_p+\alpha_p^{-1})$. Its exceptional set
(Definition~\ref{def:exceptional}) is $S(\pi)=\{\infty,2\}$, provided
that $\tau(p)\neq0$ for every odd prime $p$. Test functions are
\[
f=f_\infty\otimes f_2\otimes\bigotimes_{p\ \mathrm{odd}}1_{K_p},\qquad
\overline f=\overline f_\infty\otimes\overline f_2\otimes
\bigotimes_{p\ \mathrm{odd}}e_{K_p},
\]
with $1_{K_p}$ and $e_{K_p}$ the unit elements of the two spherical Hecke
algebras.

\item \emph{The cover} (Chapter~\ref{ch:background}). At the real place,
$(-1,-1)_2=-1$, so $(-I,1)^2=(I,-1)$: the center of
$\overline{SL(2,{\mathbb R})}$ is cyclic of order $4$, and the cover does
not split over $SO(2)$ (Lemma~\ref{lem:facts}(d),~(e)). The same holds
for the center at the place $2$ (the example at the end of
\S\ref{sec:dyadic}). At an odd prime $p$ the cover splits over
$K_{0,p}=SL(2,{\mathbb Z}_p)$ by $\kappa_p(k)=(k,s_p(k))$, and the
splitting is not $k\mapsto(k,1)$. For
$\gamma=\left(\begin{smallmatrix}a&b\\c&d\end{smallmatrix}\right)\in
SL(2,{\mathbb Z})$ with $cd\neq0$, formula (\ref{eq:sv}) gives
$\prod_{p\ \mathrm{odd}}s_p(\gamma)=\left(\frac d{c'}\right)$, the Jacobi
symbol, where $c'$ is the odd part of $|c|$. This is the symbol in the
multiplier of the theta function $\theta(z)=\sum_ne^{2\pi in^2z}$.

\item \emph{Representations} (Chapter~\ref{ch:reps}). At an odd prime
$p$, $\pi_p$ is the spherical constituent of $I_A(\chi)$ with
$\chi(a(p))=t_p=\alpha_p^2$. At the real place $\pi_\infty$ is one of the
two discrete series $\pi_{\pm11}$ into which the discrete series of
weight $12$ of $GL(2,{\mathbb R})$ decomposes.

\item \emph{The torus} (\S\ref{sec:charsA}, \S\ref{sec:realA}). At $p=3$,
$(3,3)_2=-1$, so a genuine character of $\overline{A({\mathbb Q}_3)}$
takes a value $\pm i\vartheta$ at $(a(3),1)$ where the corresponding
character of $A({\mathbb Q}_3)$ takes the value $\vartheta$. At $p=5$,
$(5,5)_2=1$ and the values agree up to sign. At the real place the
factor in the matching of functions is $-i$ for $x<0$
(Lemma~\ref{lem:realA-match}). These constants are the values of Weil's
character $\chi_\psi$ for the character $\psi$ above, up to the choices
discussed in Example~\ref{ex:eta-psi}.

\item \emph{The odd places} (\S\ref{sec:spherical}). The unit elements
$1_{K_p}$ and $e_{K_p}$ correspond under the Satake correspondence, and
a genuine unramified representation with parameter $t'$ corresponds to
the representation of $SL(2,{\mathbb Q}_p)$ with $t=\pm t'$
(Lemma~\ref{lem:square}). So a direct partner of $\pi$ would have
$t'_p=\pm\alpha_p^2$.

\item \emph{The real place} (\S\ref{sec:realds}). The even elliptic
matching transfers the genuine discrete series
$\overline\pi{}^{\pm}_{13/2}$ to $\pi_{\pm11}$, by
Lemma~\ref{lem:realds}(a) with $k=13$. This is weight $13/2$ going to
weight $12$, the archimedean form of the Shimura correspondence.

\item \emph{The place $2$} (\S\ref{sec:dyadic}). No matching is
established. Everything that the comparison would need at $2$ is a
hypothesis.

\item \emph{The global statements} (Chapter~\ref{ch:global}). For this
$\pi$, Hypothesis~\ref{hyp:GH} is an identity of traces for all pairs
$(f,\overline f)$ as above with $(f_\infty,\overline f_\infty)$ and
$(f_2,\overline f_2)$ in prescribed classes, and Hypothesis~\ref{hyp:LH}
is a condition on those two classes. Proposition~\ref{thm:main} would
then give $m(\pi)=1$ and a genuine cuspidal partner with parameters
$t'_p=\pm\alpha_p^2$. Theorem~\ref{thm:nonexist} says that there is no
such genuine cusp form. The partner that does exist is the form of
weight $13/2$ attached to $\Delta$ by Shimura, with $t'_p=\pm\alpha_p$
and $t_p=t_p'^{\,2}$.

\item \emph{The trace formulas} (\S\ref{sec:where-fails},
Chapter~\ref{ch:cover}). The scalar of the intertwining operator is
$\xi(s)/\xi(1+s)$ for $SL(2)$, with a pole at $s=1$ that gives the
constants, and $\xi(2s)/\xi(1+2s)$ for the cover, with a pole at
$s=1/2$ that gives $\theta$. Four theta representations are unramified
outside $S$, those of the square classes of $\pm1$ and $\pm2$, and
$GL(2,{\mathbb Q})$ permutes them (\S\ref{sec:cover-spectral}). The
hyperbolic classes that contribute are those of $a(\pm2^j)$, $j\geq1$,
and their orbital integrals agree under the direct comparison. The
elliptic element
$\gamma=\left(\begin{smallmatrix}2&-1\\3&-1\end{smallmatrix}\right)$ of
order $6$ has $2+\tr\gamma=3$, and its term in the trace formula of the
cover is $-i$ times its term in the trace formula of $SL(2)$
(Example~\ref{ex:epsilon}).
\end{itemize}

%% ===== begin aabook1.tex =====
\chapter{Background on covers of \texorpdfstring{$SL(2)$}{SL(2)}}\label{ch:background}

\section{Local covers of \texorpdfstring{$SL(2)$}{SL(2)}}\label{sec:local}

This section constructs the 2-fold cover of $SL(2,F)$ over a local field $F$ from Kubota's cocycle. It also records the facts about the cover used later: which subgroups split, which are abelian, and how conjugation acts.


Let $F$ be a local field of characteristic zero and let $\mu_2=\{\pm 1\}$.
Let $p=\infty$ if $F$ is archimedean and otherwise
let $p$ denote the residual characteristic of $F$.
If $F$ is a p-adic field, let ${\cal O}$ denote its ring of integers,
$\varpi$ a uniformizer, $q$ the cardinality of the residue field, and
$\nu =\nu _F=\nu _v:F^{\times }\rightarrow {\mathbb{Z}}$ the
normalized valuation, so that $\nu (\varpi )=1$; we write
$|x|_F=q^{-\nu(x)}$, so that $|\varpi|_F=q^{-1}$.


For $g=\left(
\begin{array}{cc}
a & b \\
c & d
\end{array}
\right) \in SL_2(F)$, let
\[
x(g)=\left\{
\begin{array}{c}
c,\;\;\;c\neq 0, \\
d,\;\;\;c=0,
\end{array}
\right.
\]
and let $\beta =\beta _v:SL_2(F)\times SL_2(F)\rightarrow \mu _2$ be defined
by
\begin{align}
\beta (g_1,g_2) &=(x(g_1),x(g_2))_2\cdot (-x(g_1)^{-1}x(g_2),x(g_1g_2))_2 \label{eq:beta}\\
&=(\frac{x(g_1g_2)}{x(g_1)},\frac{x(g_1g_2)}{x(g_2)})_2,  \nonumber
\end{align}
where $(...,...)_2=(...,...)_{2,v}:F^{\times }
\times F^{\times }\rightarrow \mu _2$
denotes the Hilbert symbol \cite{W1}.
Note that this is trivial if $F={\mathbb {C}}$.

\begin{lemma}\label{lem:kubota}
(Kubota \cite{K1}, Moore \cite{M})
The map $\beta $ defines a continuous Borel
2-cocycle (giving $\mu _2$ the discrete topology of course). Its cohomology
class in the continuous cohomology group $H^2(SL_2(F),\mu _2)$ corresponds
to a 2-fold topological central extension $\overline{SL_2(F)}$ which fits
into the short exact sequence
\[
1\rightarrow \mu _2\rightarrow \overline{SL_2(F)}\stackrel{p}{\rightarrow }%
SL_2(F)\rightarrow 1.
\]
\end{lemma}

If $H<SL_2(F)$ is a subgroup then we define $\overline{H}=p^{-1}(H)\subset
\overline{SL_2(F)}$.

If $g\in GL_2(F)$, let
\begin{align*}
p(g) &=\left(
\begin{array}{cc}
1 & 0 \\
0 & \det g
\end{array}
\right) ^{-1}g, \\
d(a,b) &=\left(
\begin{array}{cc}
a & 0 \\
0 & b
\end{array}
\right) , \\
g^x &=x^{-1}gx,
\end{align*}
and define the extension $\beta _1$ of the above cocycle to
$G_1(F)=GL_2(F)$ by
\begin{equation}
\beta _1(g_1,g_2)=\beta (p(g_1)^{d(1,\det g_2)},p(g_2))\cdot v(\det
g_2,p(g_1)),\label{eq:beta1}
\end{equation}
where $v:F^{\times }\times SL_2(F)\rightarrow \mu_2$ is defined by
\[
v(z,g)=\left\{
\begin{array}{c}
1,\;\;\;c\neq 0, \\
(z,d)_2,\;\;\;\;\;\;\;c=0,
\end{array}
\right.
\]
for every local field $F$, including $F=\mathbb R$. This defines a
2-fold metaplectic cover of $GL_2(F)$; see \cite[\S2.1]{Ge2}.
The 1995 draft incorrectly set $v\equiv 1$ at the real place.
That exception violates the cocycle identity; for example, the identity
applied to $-I_2$, $w$, and $d(1,-1)$ forces the factor
$(-1,-1)_{\mathbb R}$ arising from $v(-1,-I_2)$.

Elements of $\overline{SL_2(F)}$ will be denoted by $(g,\varsigma )$,
where $g\in SL_2(F)$, $\varsigma \in \mu _2$.
Multiplication on $\overline{SL_2(F)}$ is given by
\begin{equation}
(g_1,\varsigma _1)\cdot (g_2,\varsigma _2)=(g_1g_2,\beta (g_1,g_2)\varsigma
_1\varsigma _2),\label{eq:mult}
\end{equation}
so inverses are given by
\begin{equation}
(g,\varsigma )^{-1}=(g^{-1},\beta (g,g^{-1})^{-1}\varsigma ^{-1})\label{eq:inv}
\end{equation}
and conjugation is given by
\begin{equation}
(g,\varsigma )\cdot (h,\varsigma ^{\prime })\cdot (g,\varsigma
)^{-1}=(ghg^{-1},\beta (gh,g^{-1})\beta (g,h)\beta (g,g^{-1})^{-1}\varsigma
^{\prime }),\label{eq:conj}
\end{equation}
for $g,h,g_1,g_2\in G(F)$.

These identities also hold for $\beta _1$ in place of $\beta$,
for $g,h,g_1,g_2\in G_1(F)$.

We let $B=\{\left(
\begin{array}{cc}
\ast  & * \\
0 & *
\end{array}
\right) \}\subset G=SL_2$ denote the standard Borel, $N$ its unipotent
radical, and $A$ its Levi component. Let
$\theta\in F^\times -F^{\times 2}$ and let
$T=T_\theta\subset G(F)$ denote the non-split Cartan subgroup
\[
T=\{
    \left(
    \begin{array}{cc}
    a&b\theta\\
    b&a
    \end{array}
    \right)
    \ |\ a^2-b^2\theta =1,\ a,b\in F
  \}.
\]
\begin{lemma}\label{lem:Ainv}
For all $g\in A(F)$, 
\[
g=\left(
\begin{array}{cc}
a & 0 \\
0 & a^{-1}
\end{array}
\right) ,
\]
$\varsigma \in \mu _2$, we have $(g,\varsigma
)^{-1}=(g^{-1},(a,a)_2\varsigma ^{-1})$.
\end{lemma}

\begin{lemma}\label{lem:facts}
(a) $\overline{N(F)}=N(F)\times \mu _2$ (direct product),

(b) if $p>2$ then
$\overline{SL_2({\cal O})}\cong SL_2({\cal O})\times
\mu _2$ (direct product; the splitting is unique, and it is not
$k\mapsto(k,1)$),

(c) if $K_0=SL_2({\cal{O}})$ and
\[
K_{m}=\{\left(
\begin{array}{cc}
a & b \\
c & d
\end{array}
\right) \in SL_2({\cal{O}})\;|\;a-1,d-1,b,c\in \varpi ^m{\cal{O}}\},
\;m>0,
\]
then $\overline{K_{m}}=K_{m}\times \mu _2$
(direct product) for all sufficiently large $m$,

(d) if $F={\mathbb {R}}$, and $K=SO(2,{\mathbb {R})}$ then
$\overline{K}$ does not split,

(e) The center of $\overline{SL_2(F)}$ is $Z(\overline{SL_2(F)})=\{\pm
1_2\}\times \mu _2$ (direct product if $(-1,-1)_2=1$).

(f) $\overline{A(F)}$ is abelian.

(g) $\overline{T}$ is abelian.

(h) If $g\in G(F)$ and
$n=
\left(
\begin{array}{cc}
1 & 1\\
0 & 1
\end{array}
\right)
$
then $(g,\zeta)(n,\zeta ')(g,\zeta)^{-1}
=(gng^{-1},\zeta')$.

\end{lemma}

\begin{remark}
These well-known facts
were proven in \cite{K2} or follow from the properties of the Hilbert
symbol \cite{W1}.
For example, (g) is a consequence of a calculation in \cite{Fl}, p. 128.
The splitting in (b) is unique because $SL_2({\cal O})$ has no quotient of
order $2$ when $p$ is odd. It is given by (\ref{eq:sv}) below.
For (h), let $c$ be the lower left entry of $g$. If $c\neq0$, then
$\beta(g,n)=(1,c)_2=1$, $\beta(g,g^{-1})=(c,-c)_2=1$, and
$\beta(gn,g^{-1})=(-c,c)_2=1$, because $gng^{-1}$ has lower left entry
$-c^2$. So the sign in (\ref{eq:conj}) is $1$. The case $c=0$ is similar.
\end{remark}

\subsection{Local algebras of functions}\label{sec:localalg}

Here we fix the spaces of test functions on the cover: the genuine compactly supported smooth functions, and the genuine spherical Hecke algebra.


If $H$ is the group of $F$-rational points of a reductive algebraic group
over an archimedean
or p-adic field $F$, let $C_c^\infty (H)$ denote the usual
algebra of compactly supported smooth (if $F$ is p-adic: locally constant)
functions under convolution, and, for an open compact subgroup $K$ of
$H$ ($F$ p-adic), let $C_c^\infty (H//K)$ denote the subalgebra of
functions in $C_c^\infty(H)$ which are bi-$K$-invariant. A
function $f:\overline{G(F)}\rightarrow {\mathbb{C}}$ such that
\begin{equation}
f(g,\varsigma )=\varsigma \cdot f(g,1)\label{eq:genuine}
\end{equation}
will be called \textbf{genuine}. For a function $f$ on
$\overline{G(F)}$ we write $f|_{G(F)}$ for the function $g\mapsto f(g,1)$
on $G(F)$. Let
\begin{align*}
C_c^\infty (\overline{G(F)})_{gen}
&=\{f:\overline{G(F)}\rightarrow {\mathbb{C}}\
|\;(i)\;f|_{G(F)}\in C_c^\infty (G(F)),
\\
&\quad (ii)\;f\;\mathrm{is\;genuine}\}
\end{align*}
denote the Schwartz-Bruhat algebra of smooth genuine functions.
If $F$ is p-adic, $K=G({\cal O})$, and $p>2$, let
$\kappa:K\rightarrow\overline{K}$, $\kappa(k)=(k,s_K(k))$, be the splitting
of $\overline{K}$ over $K$ (Lemma~\ref{lem:facts}(b)), and let
\begin{align*}
{\cal{H}}(\overline{G(F)},K)_{gen}
&=\{f\in C_c^\infty (\overline{G(F)})_{gen}
\ |\ f(\kappa(k)x\kappa(k'))=f(x)\ \text{for all } k,k'\in K,\
x\in\overline{G(F)}\}
\end{align*}
denote the Hecke algebra of genuine spherical functions. It is an
algebra under convolution because its elements have compact support.
We also write ${\cal H}(G(F),K)=C_c^\infty(G(F)//K)$.

\begin{example}
The unit element of ${\cal{H}}(\overline{G(F)},K)_{gen}$ (for the
measure giving $\overline{K}$ volume $1$, as in the next subsection) is
the function $e_K$ with $e_K(k,\varsigma)=\varsigma\, s_K(k)$ for $k\in K$ and
$e_K=0$ off $\overline{K}$. Thus $e_K|_{G(F)}=s_K\cdot 1_K$, which is
not the characteristic function of $K$, because $s_K$ is not identically
$1$. This is why bi-invariance is imposed with respect to
$\kappa(K)$, not by requiring $f|_{G(F)}$ to be bi-$K$-invariant.
\end{example}

\subsection{Local measures}

This subsection fixes Haar measures on $G(F)$, on its subgroups, and on their covers, so that orbital integrals on the two groups can be compared.


Let $F$ be a local field and let $\overline{F}$ denote a separable
algebraic closure.
Let $dk$ denote the Haar measure on $K=G({\cal O})$
normalized so that $K$ has
volume 1. Fix an additive non-trivial character $\psi$ of
$F$ and let $dn$ denote a Haar
measure on $N(F)$ which is induced by a self-dual measure on
the additive group $F$ (with respect to $\psi$).
Let $da$ be a Haar measure on $A(F)$. Every Cartan subgroup $T$ of
$G$ defined over $F$ is either $G(F)$-conjugate to $A$, or anisotropic
over $F$. In the second case $T(F)$ is compact, and we let $dt$ be the
Haar measure on $T(F)$ of total volume $1$. (An isomorphism
$T\rightarrow A$ defined over $\overline{F}$ does not carry $T(F)$ to
$A(F)$, so it does not by itself transport Haar measures between these
groups. Any fixed normalization will do for the local comparisons below, provided
the same one is used on $T(F)$ and on $\overline{T(F)}$.)
Let $dg$ be the Haar measure
on $G(F)$ such that $dg=\delta _B(a)\,dk\,da\,dn$, $g=kan$, where
\begin{equation}
\delta _B(a)=|t|^2,\;\;a=\left(
\begin{array}{cc}
t & 0 \\
0 & t^{-1}
\end{array}
\right)
\end{equation}
is the modulus of $B$.

We extend these measures to the 2-fold cover by imposing the discrete
measure $d\varsigma $ on $\mu _2$ which gives $\mu _2$ total volume 1. For
example, the measure on $\overline{G(F)}$ will be denoted
$d\overline{g}=dg\cdot d\varsigma $.

\section{Global aspects}\label{sec:global}

This section assembles the local covers into the adelic cover and explains how $G(k)$ embeds in it.


Let $k$ denote a number field. Let $\Sigma _k$ denote the set of places of
$k$, $\Sigma _k^\infty $ denote the set of archimedean places of
$k$ (which we may identify with the set of inequivalent
complex embeddings $k\hookrightarrow {\mathbb{C}}$), and let
$\Sigma _k^f$ denote the set of non-archimedean places of
$k$ (which we may identify with the set of non-zero prime ideals of
${\cal{O}}_k$). Let ${\mathbb{A}}={\mathbb{A}}_k$ denote the adele
ring of $k$ and, for $v\in \Sigma _k$, let $k_v$ denote
the completion of $k$ with respect to $v$.

Let $G=SL(2)$.

The global cover is constructed from the local ones as follows
\cite{K2}. For each odd finite place $v$ let
$\kappa_v:K_{0,v}\rightarrow\overline{SL_2(k_v)}$ be the splitting of
Lemma~\ref{lem:facts}(b). It is $\kappa_v(k)=(k,s_v(k))$, with
$s_v$ as in (\ref{eq:sv}) below \cite{K2}, \cite{Ge2}. Let
$\prod'_v\overline{SL_2(k_v)}$ denote the restricted direct product of the
local covers with respect to the subgroups $\kappa_v(K_{0,v})$, and let
\[
Z_0=\Bigl\{(\varsigma_v)_v\in\bigoplus_{v\in\Sigma_k}\mu_2\ \Bigm|\
\prod_v\varsigma_v=1\Bigr\},
\]
a discrete central subgroup (only finitely many $\varsigma_v$ are
$\neq 1$). Then
\[
\overline{G({\mathbb{A}})}=\Bigl({\prod_v}'\,\overline{SL_2(k_v)}\Bigr)\Big/Z_0 ,
\]
and the product map $\mu_2^{(\Sigma_k)}\rightarrow\mu_2$ identifies the
image of $\bigoplus_v\mu_2$ with a central subgroup $\mu_2$, giving the
exact sequence
\begin{equation}
1\rightarrow \mu _2\rightarrow
\overline{G({\mathbb{A}})}\stackrel{p}{\rightarrow }G({\mathbb{A}})\rightarrow 1.
\end{equation}
Thus $\overline{G({\mathbb{A}})}$ is a \emph{quotient} of the restricted
product of the local covers, not a subgroup of it. In these original
restricted-product coordinates the diagonal lift of
$\gamma\in G(k)$ is represented by the family $((\gamma,1))_v$.
Hilbert reciprocity makes this a homomorphism after passing to $Z_0$.
At almost every place this family belongs to the prescribed compact
subgroup, since $s_v(\gamma)=1$ there.

After changing the local coordinates by the cochains $s_v$, the same
extension is represented on $G(\mathbb A)\times\mu_2$.
In these \emph{modified-cocycle coordinates}, the splittings of
$K_{0,v}$ are represented by $k\mapsto(k,1)$.
Equivalently, $\overline{G({\mathbb{A}})}$ is the set $G({\mathbb{A}})\times\mu_2$
with the multiplication (\ref{eq:mult}) for the 2-cocycle
\[
\beta_{\mathbb{A}}(g_1,g_2)=\prod_{v\in \Sigma_k}\beta'_v(g_{1,v},g_{2,v}),\qquad
\beta'_v(g_1,g_2)=\beta_v(g_1,g_2)\,s_v(g_1)\,s_v(g_2)\,s_v(g_1g_2)^{-1}
\]
(with $s_v$ as in (\ref{eq:sv}), and $s_v=1$ at the archimedean places). The cocycle $\beta'_v$
is cohomologous to $\beta_v$ and is trivial on $K_{0,v}\times K_{0,v}$ for
almost all $v$ \cite{K2}, so for $g_1,g_2\in G({\mathbb{A}})$ the product
has only finitely many factors different from $1$. (The product of the
unmodified $\beta_v$ is not a finite product in general.) The rules of multiplication, inverses, and conjugation are as in
(\ref{eq:mult})--(\ref{eq:conj}), with $\beta_{\mathbb{A}}$ in place of
$\beta$. If $H<G({\mathbb{A}})$ is a subgroup then we define
$\overline{H}=p^{-1}(H)<\overline{G({\mathbb{A}})}$.

\begin{example}
Let $v_1\neq v_2$ be places of $k$. The element of the restricted
product with $\varsigma_{v_1}=\varsigma_{v_2}=-1$ and all other
components trivial lies in $Z_0$, so it is the identity of
$\overline{G({\mathbb{A}})}$. The element with $-1$ at $v_1$ only is the
nontrivial central element of $\mu_2\subset\overline{G({\mathbb{A}})}$.
\end{example}

\begin{lemma}\label{lem:Gk-split}
$\overline{G(k)}\cong G(k)\times \mu _2$,
where the splitting isomorphism in $\cong$ will be
described below.
\end{lemma}

\begin{proof} This is well-known \cite{K2}.
\end{proof}

In the modified-cocycle coordinates of the preceding paragraph, the
splitting of the above exact sequence is
$\gamma\longmapsto(\gamma,s_{\mathbb A}(\gamma)^{-1})$,
rather than $\gamma\longmapsto(\gamma,1)$.
Indeed, Hilbert reciprocity gives
$\beta_{\mathbb A}(\gamma_1,\gamma_2)
=s_{\mathbb A}(\gamma_1)s_{\mathbb A}(\gamma_2)
s_{\mathbb A}(\gamma_1\gamma_2)^{-1}$ on $G(k)$.
This expression should not be confused with the diagonal lift
$((\gamma,1))_v$ in the original restricted-product coordinates.
Here
\[
s_{{\mathbb{A}}}(\gamma )=
\prod_{v\in \Sigma _k}
s_v(\gamma),
\]
and
\begin{equation}\label{eq:sv}
s_v(\gamma)=\left\{
\begin{array}{c}
(c,d)_{v,2}^{-1},
\;\;\;cd\neq 0{\rm \ and\ }
2\;\mathrm{does\;not\;divide}\;v(c)
{\rm \ and\ }v\notin \Sigma_k^\infty, \\
1,\;\;\;\;\;\;\;\;\;\;\;\;\;\;\;\;
\mathrm{\;otherwise}.\;\;\;\;\;\;\;\;
\end{array}
\right.
\end{equation}
Here $\gamma =\left[
\begin{array}{cc}
a & b \\
c & d
\end{array}
\right] \in G(k)$.

{\em From now on, we identify $G(k)$ with its image
in $\overline{G({\mathbb{A}})}$ under this splitting (and the diagonal
embedding).} We identify $N({\mathbb{A}})$ with the
subgroup $N({\mathbb{A}})\times 1\subset \overline{G({\mathbb{A}})}$.

Let $K_0=\Pi_{v\in \Sigma_k}K_{0,v}$
denote the ``usual'' maximal compact subgroup of $G({\mathbb{A}})$.
Here, if $v\in \Sigma_k^f $ then $K_{0,v}=G({\cal O}_v)$ and
otherwise
$K_{0,\infty}=\prod_{\substack{ v|\infty\\ v\in\Sigma_k}}
K_v$, where $K_v=SO(2,{\mathbb{R}})$ if $v$ is real and
$K_v=SU(2,{\mathbb{C}})$ if $v$ is complex.

\subsection{Global algebras of functions}\label{sec:globalalg}

The global test functions are restricted tensor products of local ones.


A function $f:\overline{G({\mathbb{A}})}\rightarrow {\mathbb{C}}$
such that (\ref{eq:genuine})
holds will be called \textbf{genuine}. Let
${\cal{H}}(\overline{G({\mathbb{A}})})$ denote the algebra of
functions generated by the elements of the
restricted tensor product
$\bigotimes_{v\in \Sigma_k }
C_c^\infty (\overline{G(k_v)})_{gen}$,
where almost all $f_v$ in the product
$f=\otimes _{v\in \Sigma_k}f_v\in
\bigotimes_{v\in \Sigma_k }C_c^\infty (\overline{G(k_v)})_{gen}$ are
spherical in the sense of \S\ref{sec:localalg}.

\subsection{Global measures}

Globally we use Tamagawa measures, extended to the cover as in the local case.


Let $dk$ denote the Haar measure on $K=\prod_{v\in \Sigma_k }K_v$
defined by $dk=\prod_{v\in \Sigma_k }dk_v$.
Let $dg$ be the Tamagawa measure on
$G({\mathbb{A}})$ \cite{W2}. For any Cartan subgroup
$T$ of $G$ defined over $k$, let $dt$
denote the global Tamagawa measure on
$T({\mathbb{A}})$. If $T=A$ then this will
be denoted by $da$. Let $dn$ denote the product measure on $N({\mathbb{A}})$.

We extend these measures to the 2-fold cover by imposing the discrete
measure $d\varsigma $ on $\mu _2$ which gives $\mu _2$ total volume 1. For
example, the measure on $\overline{G({\mathbb{A}})}$ will be denoted
$d\overline{g}=dg\cdot d\varsigma $.
%% ===== end aabook1.tex =====



%% ===== begin aabook3.tex =====
\chapter{Representations of \texorpdfstring{$\overline{SL(2,F)}$}{SL(2,F)-bar}}\label{ch:reps}

This chapter presents a survey of certain aspects of the
representation theory of central covering groups of p-adic reductive groups
in the case of $SL(2)$.

\section{Characters}\label{sec:chars}

Characters occur naturally in the trace formula, so we recall
some of their properties.

For a group $H$ as in \S\ref{sec:localalg}, or its cover, $\Pi(H)$ denotes
the set of equivalence classes of irreducible admissible representations
of $H$ (genuine ones, for a cover). For $H=A(F)$ or $\overline{A(F)}$ these
are the characters, not necessarily unitary. $\Pi_u(H)$ denotes the subset
of unitary classes. For $\chi\in\Pi(A(F))$ or $\chi\in\Pi(\overline{A(F)})$,
$I_A(\chi)$ denotes the representation of $G(F)$, respectively of
$\overline{G(F)}$, unitarily induced from the character of $B(F)=A(F)N(F)$,
respectively of $\overline{A(F)}N(F)$, which is $\chi$ on the torus and
trivial on $N(F)$.

\subsection{The character as a function}

Harish-Chandra's theorem that characters are given by locally integrable functions extends to the cover. We state the form needed.


The existence of the character as a function on covering groups is proved
in \cite{Li12} (a reference not consulted directly).
The argument in \cite{S} extends to the metaplectic case with only
minor modification.

Let $\overline{G(F)}^{\prime }=\{(g,\varsigma )\in \overline{G(F)}%
\;|\;g\;{\rm has\;distinct\;eigenvalues}\},$ called the {\bf regular set} of
$\overline{G(F)}.$ In particular, every element in the regular set is
semi-simple (i.e., diagonalizable). The following result is a corollary of a
result of Howe and Harish-Chandra \cite{S}:

\begin{lemma}
Let $F$ be p-adic and $\pi\in\Pi(\overline{G(F)})$.
There exists a genuine locally constant function denoted $\Theta _\pi $ on
the regular set $\overline{G(F)}^{\prime }$ which represents the trace of $%
\pi :$%
\[
\tr(\pi (f))=\int_{\overline{G(F)}}\Theta _\pi (x)f(x)dx,
\]
for all $f\in C_c^\infty (\overline{G(F)})_{gen}.$
\end{lemma}

This trace, which we sometimes denote simply by $\Theta _\pi (f)$, is called
the {\bf Fourier transform} of $f$ at $\pi $. The function $\Theta _\pi (x)$
is called the {\bf character} of $\pi $.

\subsection{Principal series representations}\label{sec:ps-char}

The character of a principal series representation is supported on the conjugates of the split torus, where it is given by an explicit formula. We use this to write the trace as a Fourier transform of an orbital integral.


Let ${\cal T}$ denote a complete set of representatives of conjugacy classes of
Cartan subgroups of $G(F)$. We may and do assume that $A\in {\cal T}$. For each
$T\in {\cal T}$, let $W_T=N_{G(F)}(T)/T$ denote the Weyl
group of $T$. If $t\in T$ is regular then
\[
T=G(F)_t=\operatorname{Cent}(t,G(F)).
\]

\begin{lemma}
Let $\chi\in \Pi (\overline{A(F)})$. In the formula below,
$a$ lies over $\operatorname{diag}(a_1,a_2)\in A(F)$,
so $a_1a_2=1$ and
$|D(a)|^{1/2}=|a_1-a_2|$.
The character of the
(possibly reducible) representation $\pi =I_{A}(\chi )$\ is given on the regular
set by
\[
\Theta _\pi (x)=\left\{
\begin{array}{c}
\frac{\chi (a)+\chi (a^w)}{|a_1-a_2|},
\;\;x=a^g,\;{\rm for\;some}\;g\in \overline{G(F)},
\;a\in \overline{A(F)}, \\
0,\;\;\;\;\;\;\;\;\;\;\;\;\;\;\;\;\;\;\;\;\;\;\;\;\;\;\;\;\;
{\rm \;if\;}\;\;\;
x\in \overline{G(F)}^{\prime }-\overline{A(F)}^{\overline{G(F)}},
\end{array}
\right.
\]
where
\[
\overline{A(F)}^{\overline{G(F)}}=\{a^g\;|\;a\in \overline{A(F)},\;g\in
\overline{G(F)}\}.\
\]
\end{lemma}

The proof, which is omitted, follows the proof of \cite{Ge1}, Theorem 5.19.
The formula does not require $I_A(\chi)$ to be irreducible.

\begin{remark} As noted above, on the regular set,
this character vanishes off the conjugacy class of $\overline{A(F)}$.
\end{remark}

Now we describe the character of a principal series representation
as the Fourier transform of the function $a\longmapsto J_G(a,f)$.
Here $J_G(a,f)=F_f^{\overline{A(F)}}(a)=|D(a)|^{1/2}\int_{\overline{A(F)}\backslash\overline{G(F)}}f(x^{-1}ax)\,dx$
is the normalized orbital integral of \S\ref{sec:weighted-defs}, and
$|D(a)|^{1/2}=|a_1-a_2|$ for $a$ over $\operatorname{diag}(a_1,a_2)$.

Let $\chi \in \Pi (\overline{A(F)})$,
$\pi =I_{A}(\chi )\in \Pi (\overline{G(F)})$.
This lemma and the Weyl integration formula
imply that, for $f\in C_c^\infty
(\overline{G(F)})$,
\begin{equation}
\begin{array}{c}
\Theta _\pi (f)=\int_{\overline{G(F)}}\Theta _\pi (x)f(x)dx \\
=\frac 12\int_{\overline{A(F)}}(\chi (a)+
\chi(a^w))F_f^{\overline{A(F)}}(a)da \\
=(F_f^{\overline{A(F)}})^{\wedge }(\chi ) ,
\end{array}
\end{equation}
where the Fourier transform is defined by
\begin{equation}
\begin{array}{c}
\varphi ^{\wedge }(\chi )=\int_{\overline{A(F)}}\chi (x)\varphi
(x)da \\
\ \ \ \ \ \ \chi \in \Pi (\overline{A(F)}).
\end{array}
\end{equation}
The last equality in the first display holds because $F_f^{\overline{A(F)}}$ is
invariant under $a\mapsto a^w$, since $a$ and $a^w$ are conjugate in
$\overline{G(F)}$. The same formulas, without bars, hold for $G(F)$.

\begin{lemma}
Let $\varphi(\chi)=\tr I_A(\chi)(f)$, and let $\varphi^\vee$ be its inverse
Fourier transform on $\overline{A(F)}$. Then
\[
\varphi^\vee(\gamma)
=J_G(\gamma,f),
\]
for $\gamma =(a,\zeta)\in \overline{A(F)}$, $a\not= \pm 1$.
\end{lemma}

\begin{proof} This follows from our results above.
\end{proof}

\begin{example}
Let $\chi\neq\chi^w$ and let $x=(t,\zeta)$ with $t$ regular elliptic. Then
$x$ is not conjugate to an element of $\overline{A(F)}$, so the lemma above
gives $\Theta_{I_A(\chi)}(x)=0$. Thus principal series characters vanish on
the regular elliptic set. This fact is used in
\S\ref{sec:weil}, in the proof of the matching for non-tempered
representations.
\end{example}

\section{Description of the dual spaces}

The dual of $\overline{G(F)}$ is classified in \cite{Wa2}.
We shall recall some details from that paper and the
papers of Moen and Ariturk.

\subsection{Principal series, reducible principal series}

We recall when genuine principal series representations are irreducible.


Let $\chi \in \Pi (\overline{A(F)})$. The induced representation
$I_{A}(\chi )$ is in general not irreducible. However, we do have the
following

\begin{theorem}
(Moen \cite{Mo2}) If $p>2$ then $I_{A}(\chi )$ is
irreducible for all $\chi \in \Pi _u(\overline{A(F)})$.
\end{theorem}

In fact, C. Moen \cite{Mo1}
explicitly computes the intertwining operators as
matrices using the Kirillov model.

\subsection{Complementary series}

Some principal series induced from non-unitary characters are unitarizable. These form the complementary series.


Let $\chi\in\Pi(\overline{A(F)})$ and write
$\chi(a(x),1)=\chi_\psi(x)\mu(x)$ with $\chi_\psi$ as in
\S\ref{sec:nonexist} and $\mu$ a character of $F^\times$. We write
$I_A(\mu)$ for $I_A(\chi)$. If $\mu (x)=\mu _0(x)|x|^s$, for some
character $\mu _0$ of finite order and some $s\in {\mathbb C}$, then
we write $s=s(\mu )=s(\chi )$.

\begin{lemma} If $0<s<1/2$ and $\mu (x)=|x|^s$ then $I_A(\mu )$ is unitary.
\end{lemma}

This follows from \cite{GS}, Proposition 1 following
the argument of \cite{Ar}. See also
\cite{Mo1} if $p>2$ and \cite{Ge1}, \S 5.4.

\subsection{Special representation, non-tempered representation}

At the end of the complementary series the principal series reduces. Its two pieces play the roles that the Steinberg and trivial representations play for $SL(2)$.


The following result is due to Gelbart and Sally \cite{GS}:

\begin{lemma}
If $s=1/2$ and $\mu (x)=|x|^s$ then there is an irreducible
constituent of $I_A(\mu)$, ``the special
representation'' $\pi _{sp}$, which is elliptic (in the sense that its
character has support on the elliptic set) and square-integrable.
Its complement in $I_A(\mu)$ is an infinite-dimensional,
elliptic, non-tempered representation $\pi _{nt}$.
\end{lemma}

Furthermore, in this case

(1) the representation $\pi _{sp}$ is the only square-integrable
representation in $\Pi _u(\overline{G(F)})$ which is not cuspidal.

(2) $\pi _{sp}$ is tempered.

(3) The representation $\pi _{nt}$ and the complementary series comprise the
only non-tempered representations in $\Pi _u(\overline{G(F)})$.

\subsection{Cuspidal representations}

Thanks to ideas originally introduced by Weil and Shale, much is known about
the cuspidal representations in $\Pi _u(\overline{G(F)})$ in the case
where $p>2$. See for example \cite{M1}, \cite{M2}, \cite{M3}.
The information we shall need will be recalled in Appendix~\ref{ch:sc}.
%% ===== end aabook3.tex =====



%% ===== begin aabook4.tex =====
\begin{example}
Take $\mu(x)=|x|^{1/4}$. By the lemma on complementary series, $I_A(\mu)$ is
unitary, although $\mu$ is not unitary. At $s=1/2$, that is for
$\mu(x)=|x|^{1/2}$, the lemma of Gelbart and Sally shows that $I_A(\mu)$ is
reducible, with the special representation $\pi_{sp}$ as one constituent.
\end{example}

\chapter{Local comparisons on the groups}\label{ch:local}

\section{Introduction}
We study local identities of the kind that a comparison of the
trace formula of $SL(2)$ with its metaplectic analog would need.

We shall describe, in the real case, several
correspondences between functions in
$C_c^\infty(G({\mathbb{R}}))$ and functions in
$C_c^\infty(\overline{G({\mathbb{R}})})$, each of which is associated
to a matching of characters. See \S\ref{sec:realA} and
\S\ref{sec:realds}.

We shall describe, in the p-adic case with $p\not= 2$,
several correspondences between functions in
$C_c^\infty(G(F))$ and functions in
$C_c^\infty(\overline{G(F)})$, each of which is associated
to a matching of characters. The case $p=2$ is discussed in
\S\ref{sec:dyadic}, where no matching is established in this version.

\begin{example}
The simplest corresponding pair consists of the unit elements of the two
spherical Hecke algebras at an odd place (Example~\ref{ex:units}). Their
traces agree on corresponding unramified principal series.
\end{example}

\section{Orbital integrals - basics}

This section defines the orbital integrals and weighted orbital integrals on the cover that appear in the trace formula, and identifies the centralizers involved.


Let $G(F)_\gamma$ denote the connected component of the centralizer of
$\gamma$ in $G(F)$.

\subsection{Orbital integrals: definitions}\label{sec:weighted-defs}

Certain weighted orbital integrals occur naturally in the invariant trace
formula. We will need some of their properties to put the trace formula in
stable form as well as to compare trace formulas.

If $M\in \{A,G\}$, $f\in C_c^\infty (\overline{G(F)})_{gen}$, and $g\in
\overline{M(F)}$ is regular, let
\[
J_M(g,f)=|D(g)|^{1/2}
\int_{\overline{G(F)}_g\backslash \overline{G(F)}}
f(x^{-1}gx)v_M(x)dx.
\]
If $F$ is a local field then
\[
v_M(x)=\left\{
\begin{array}{ll}
1, & M=G, \\
\log (1+|n|^2)^{1/2}, & M=A,\ F\ {\rm archimedean},\\
\log \max(1,|n|_F), & M=A,\ F\ {\rm non\hbox{-}archimedean},
\end{array}
\right.
\]
where $x=a\left(
\begin{array}{cc}
1 & n \\
0 & 1
\end{array}
\right) k$.

Note that the centralizer in $\overline{G(F)}$ of a regular element
$a\in A(F)$ is $\overline{A(F)}$. We may therefore write,
for a regular $a\in \overline{A(F)}$,
\begin{equation}
J_M(a,f)=
|D(a)|^{1/2}\int_{\overline{A(F)}\backslash \overline{G(F)}}
f(x^{-1}ax)v_M(x)dx.
\end{equation}

\subsection{Elliptic orbital integrals}

We now discuss the elliptic orbital integrals $J_G(\gamma)$,
$\gamma\in \overline{T}_{reg}$, where $\overline{T}_{reg}$
denotes the subset of elements of the form $(x,\zeta)$
where $x\in T$ has distinct eigenvalues and $T$ is an elliptic
Cartan of $SL(2,F)$.

By Lemma~\ref{lem:facts}(g), we have
\[
\overline{G(F)}_\gamma
=\overline{T},\ \ \ \ \ \
\gamma\in \overline{T}_{reg}.
\]
We have
\begin{equation}
\int_{\overline{G(F)}_\gamma\backslash \overline{G(F)}}
f(x^{-1}\gamma x)\frac{dx}{dt}
=\int_{\overline{T}\backslash \overline{G(F)}}
f(x^{-1}\gamma x)\frac{dx}{dt},
\end{equation}
where $\gamma =(g,1)\in \overline{T}_{reg}$.

\subsection{Hyperbolic orbital integrals}

For split regular elements, the centralizer in the cover is the cover of $A$.


By Lemma~\ref{lem:facts}(f), we have
\[
\overline{G(F)}_\gamma
=\overline{A(F)},\ \ \ \ \ \
\gamma\in \overline{A(F)}_{reg}.
\]
We have
\begin{equation}
\int_{\overline{G(F)}_\gamma\backslash \overline{G(F)}}
f(x^{-1}\gamma x)\frac{dx}{da}
=\int_{\overline{A(F)}\backslash \overline{G(F)}}
f(x^{-1}\gamma x)\frac{dx}{da},
\end{equation}
where $\gamma =(g,1)\in \overline{A(F)}_{reg}$.

\subsection{Unipotent orbital integrals}

The unipotent orbital integrals are taken over the cover of $N$.


Using (\ref{eq:conj}) and Lemma~\ref{lem:facts}(h), a straightforward calculation involving only the definitions
yields the following result:
\[
\overline{G(F)}_\gamma
=\overline{N(F)},\ \ \ \ \ \
\gamma=
(
\pm \left(
\begin{array}{cc}
1 & 1\\
0 & 1
\end{array}
\right)
,1)\in \overline{N(F)}.
\]
We have
\begin{equation}
\int_{\overline{G(F)}_\gamma\backslash \overline{G(F)}}
f(x^{-1}\gamma x)\frac{dx}{dn}
=\int_{\overline{N(F)}\backslash \overline{G(F)}}
f(x^{-1}\gamma x)\frac{dx}{dn},
\end{equation}
where $\gamma $ is as above.

\subsection{Weighted orbital integrals}\label{sec:weighted-central}

We write out the weighted orbital integrals at regular and at central elements of the cover of $A$.


If $\gamma =(g,1)\in \overline{A(F)}_{reg}$ and if
\[
g=
\pm \left(
\begin{array}{cc}
a & 0\\
0 & a^{-1}
\end{array}
\right)
\in A(F),
\]
then
\begin{equation}
\int_{\overline{G(F)}_\gamma\backslash \overline{G(F)}}
f(x^{-1}\gamma x)v_A(x)\frac{dx}{da}
=\int_{\overline{A(F)}\backslash \overline{G(F)}}
f(x^{-1}\gamma x)v_A(x)\frac{dx}{da}.
\end{equation}

If
\[
\gamma=
(
\pm \left(
\begin{array}{cc}
1 & 0\\
0 & 1
\end{array}
\right)
,1)\in Z_{\overline{G}}(F),
\]
the integrals $J_A(\gamma a,f)$ do not converge as $a\rightarrow1$. For
$a=a(t)$ and $x=n(u)k$ one has
$x^{-1}ax=k^{-1}a\,n\bigl((1-t^{-2})u\bigr)k$, so for $F$ p-adic
\[
J_A(\gamma a,f)=-\log|1-t^{-2}|\cdot J_G(\gamma a,f)+O(1)
\qquad(t\rightarrow1).
\]
Following \cite[(8.1)]{A7}, one removes the divergent term:
\begin{equation}
J_A(\gamma,f)=
\lim_{t\rightarrow 1}
\Bigl[J_A(\gamma a(t),f)+\log|1-t^{-2}|\,J_G(\gamma a(t),f)\Bigr]
=\int_K\int_F f\bigl(k^{-1}\gamma\,n(u)k\bigr)\log|u|\,du\,dk
\end{equation}
for $F$ p-adic, where the conjugation is computed in the cover. In the
archimedean case $\log|u|$ is replaced by the corresponding archimedean
weight. (The normalization of the correction term should be compared with
\cite{A7}.)

\begin{example}
Let $p$ be odd, $K=G({\cal O})$, and $f=1_K$. For regular $a=a(x)$, the
normalized hyperbolic orbital integral of $f$ equals its Satake transform
\cite{C}, so $F_f(a(x))=1$ if $|x|=1$ and $F_f(a(x))=0$ otherwise (with
$dn$ giving $N\cap K$ volume $1$).
\end{example}

\section{Corresponding functions}

We define the notion of a ``smooth matching''. We shall use separate
definitions for $p=\infty$ and $p>2$. For $p=2$ see
\S\ref{sec:dyadic}.
In these cases we shall base our ``matching'' of functions (or, more
properly, of orbital integrals) on our knowledge of the characters of the
representations. For the purpose of avoiding confusion in later
chapters, we shall assume $p>2$ in some sections even though
that hypothesis will not be used until later.

\subsection{Corresponding functions on \texorpdfstring{$A$}{A}}

We first match characters of the split torus and its cover, and then match functions so that their Fourier transforms agree.


We concentrate first on the set $A(F)_{reg}$.

\subsubsection{Characters of \texorpdfstring{$\overline{A(F)}$}{A(F)-bar}}\label{sec:charsA}

This subsection describes the genuine characters of $\overline{A(F)}$ and fixes a bijection between them and the characters of $A(F)$.


We assume $p>2$ throughout this subsection. (The case $p=2$ is
not covered. See \S\ref{sec:dyadic}.) We write $a(x)=\operatorname{diag}(x,x^{-1})$
and, for $a\in A(F)$, $x_a$ for the element with $a=a(x_a)$. We write
$(x,y)$ for the Hilbert symbol $(x,y)_2$. Its values are $\pm1$, so
$(x,y)^{-1}=(x,y)=(y,x)$.

\begin{lemma}\label{lem:Acocycle}
(a) For $x,y\in F^\times$, $\beta(a(x),a(y))=(x,y)$. Hence
\[
(a(x),\zeta)(a(y),\zeta')=(a(xy),(x,y)\zeta\zeta').
\]

(b) Let $w=\left(\begin{smallmatrix}0&-1\\1&0\end{smallmatrix}\right)$.
For every $a\in A(F)$, $(w,1)(a,1)(w,1)^{-1}=(a^{-1},1)$.
(This holds for every local field $F$.)
\end{lemma}

\begin{proof}
(a) Since $c=0$, $x(a(x))=x^{-1}$, so by (\ref{eq:beta})
$\beta(a(x),a(y))=((xy)^{-1}/x^{-1},(xy)^{-1}/y^{-1})=(y^{-1},x^{-1})=(x,y)$.

(b) Let $a=a(x)$. We have $x(w)=1$, $x(w^{-1})=-1$, $x(a)=x^{-1}$,
$wa=\left(\begin{smallmatrix}0&-x^{-1}\\x&0\end{smallmatrix}\right)$ so
$x(wa)=x$, and $waw^{-1}=a(x^{-1})$ so $x(waw^{-1})=x$. Hence
$\beta(w,a)=(x,x^2)=1$, $\beta(w,w^{-1})=(1,-1)=1$ and
$\beta(wa,w^{-1})=(1,-x)=1$. By (\ref{eq:mult}) and (\ref{eq:inv}),
$(w,1)^{-1}=(w^{-1},1)$ and
$(w,1)(a,1)(w^{-1},1)=(wa,1)(w^{-1},1)=(a^{-1},1)$.
\end{proof}

Let
\[
A_2(F)=\{a(x)\ |\ \nu (x)\ {\rm \ even}\}.
\]
If $\nu(x)$ and $\nu(y)$ are even then $(x,y)=1$ (as $p$ is odd), so by
Lemma~\ref{lem:Acocycle}(a), $\overline{A_2(F)}=A_2(F)\times \mu_2$
(direct product), and every genuine character of $\overline{A_2(F)}$
has the form $(a(x),\zeta)\mapsto\zeta\,\chi_1(x)$ for a character
$\chi_1$ of $A_2(F)$ (which we regard as a character of the group of
$x$ with $\nu(x)$ even). Every element of $A(F)-A_2(F)$ is $a(x)$ with
$x/\varpi$ of even valuation.

\begin{lemma}\label{lem:genuine-chars}
Let $\chi_1\in\Pi(A_2(F))$ and let $\epsilon\in{\mathbb C}$ satisfy
\begin{equation}\label{eq:eps}
\epsilon^2=(\varpi,\varpi)^{-1}\chi_1(\varpi^2).
\end{equation}
Then
\begin{equation}\label{eq:chi-eps}
\overline{\chi}_\epsilon (a(x),\zeta)=
\left\{
\begin{array}{ll}
\zeta\chi_1(x), & \nu(x)\ \text{even},\\
\epsilon\,\zeta(x,\varpi)^{-1}\chi_1(x/\varpi),
& \nu(x)\ \text{odd},
\end{array}
\right.
\end{equation}
is a genuine character of $\overline{A(F)}$, and every genuine character
of $\overline{A(F)}$ is of this form for exactly one pair
$(\chi_1,\epsilon)$. Thus for each $\chi_1$ there are exactly two genuine
characters extending $\zeta\chi_1$.
\end{lemma}

\begin{proof}
By Lemma~\ref{lem:Acocycle}(a), a genuine function
$(a(x),\zeta)\mapsto\zeta\psi(x)$ is a character if and only if
$\psi(xy)=(x,y)\psi(x)\psi(y)$ for all $x,y$. If $\psi$ is such a
function, put $\chi_1=\psi|_{A_2}$ and
$\epsilon=(\varpi,\varpi)\psi(\varpi)$. For $\nu(x)$ odd write
$x=\varpi x_0$. Then
$\psi(x)=(\varpi,x_0)\psi(\varpi)\chi_1(x_0)=\epsilon(x,\varpi)\chi_1(x/\varpi)$,
because $(x,\varpi)=(\varpi,\varpi)(x_0,\varpi)$, and
$\chi_1(\varpi^2)=\psi(\varpi^2)=(\varpi,\varpi)\psi(\varpi)^2
=(\varpi,\varpi)\epsilon^2$, which is (\ref{eq:eps}). Conversely, let
$\psi$ be defined by (\ref{eq:chi-eps}) with $\zeta=1$. If $\nu(x)$,
$\nu(y)$ are both even the relation is clear. If $\nu(x)$ is odd and
$\nu(y)$ even, it reads $(xy,\varpi)=(x,y)(x,\varpi)$, i.e.\
$(y,x\varpi)=1$, which holds because $y$ and $x\varpi$ have even
valuation. If both are odd, it reads
$\chi_1(xy)=(x,y)\epsilon^2(x,\varpi)(y,\varpi)\chi_1(xy/\varpi^2)$. By
(\ref{eq:eps}) the right side is
$(x,y)(x,\varpi)(y,\varpi)(\varpi,\varpi)\chi_1(xy)
=(x\varpi,y\varpi)\chi_1(xy)=\chi_1(xy)$, since $x\varpi$ and $y\varpi$
have even valuation.
\end{proof}

Since $A(F)/A_2(F)$ has order 2, the characters of $A(F)$ are
also in 2-1 correspondence with the characters of $A_2(F)$:
given $\chi_1\in \Pi(A_2(F))$, its two extensions to $A(F)$ are
\begin{equation}\label{eq:chi-theta}
\chi_\vartheta (a(x))=
\left\{
\begin{array}{ll}
\chi_1(x), & \nu(x)\ \text{even},\\
\vartheta\, \chi_1(x/\varpi), & \nu(x)\ \text{odd},
\end{array}
\right.
\qquad \vartheta^2 = \chi_1(\varpi^2).
\end{equation}

Comparing (\ref{eq:eps}) with $\vartheta^2=\chi_1(\varpi^2)$, we see that
$\epsilon=\vartheta$ is impossible when $(\varpi,\varpi)=-1$. Since
$(\varpi,\varpi)=(\varpi,-1)=(-1)^{(q-1)/2}$, this happens exactly when
$q\equiv 3\pmod 4$. Fix once and for all $\eta_\varpi\in{\mathbb C}$ with
\[
\eta_\varpi^2=(\varpi,\varpi)^{-1},
\]
and take $\eta_\varpi=1$ when $(\varpi,\varpi)=1$.

\begin{definition} Let $\chi_1\in \Pi(A_2(F))$ be given.
We say that
$\overline{\chi}_\epsilon\in \Pi(\overline{A(F)})$ {\bf corresponds}
to $\chi_\vartheta \in \Pi(A(F))$ if $\epsilon =\eta_\varpi\vartheta$.
If $\chi \in \Pi(A(F))$ is given we let
$\overline{\chi}\in \Pi(\overline{A(F)})$ denote
its corresponding genuine character.
\label{def:corrchar}
\end{definition}

If $\vartheta^2=\chi_1(\varpi^2)$ then $\epsilon=\eta_\varpi\vartheta$
satisfies (\ref{eq:eps}), so Definition~\ref{def:corrchar} gives a
bijection $\chi\mapsto\overline{\chi}$ between $\Pi(A(F))$ and the genuine
characters of $\overline{A(F)}$. It depends on the choices of $\varpi$
and $\eta_\varpi$ (Remark~\ref{rem:choices}).

\begin{lemma}\label{lem:Wequiv}
If $\overline\chi$ corresponds to $\chi$, then $\overline\chi^{\,w}$
corresponds to $\chi^w=\chi^{-1}$, where
$\overline\chi^{\,w}(x)=\overline\chi(x^w)$.
\end{lemma}

\begin{proof}
Let $\overline\chi=\overline\chi_\epsilon$ and $\chi=\chi_\vartheta$, with
$\epsilon=\eta_\varpi\vartheta$. By Lemma~\ref{lem:Acocycle}(b),
$\overline\chi^{\,w}(a(x),1)=\overline\chi(a(x^{-1}),1)$. This restricts to
$\chi_1^{-1}$ on $A_2(F)$, and by the proof of
Lemma~\ref{lem:genuine-chars} its parameter is
\[
\epsilon''=(\varpi,\varpi)\,\overline\chi(a(\varpi^{-1}),1)
=\epsilon\,\chi_1(\varpi^2)^{-1}=(\varpi,\varpi)\epsilon^{-1}
=(\varpi,\varpi)\eta_\varpi^{-1}\vartheta^{-1},
\]
where we used (\ref{eq:chi-eps}) and (\ref{eq:eps}). Since
$\eta_\varpi^2=(\varpi,\varpi)$, this is $\eta_\varpi\vartheta^{-1}$. The
character $\chi^{-1}$ extends $\chi_1^{-1}$ with parameter
$\vartheta^{-1}$, so its corresponding genuine character has parameter
$\eta_\varpi\vartheta^{-1}$.
\end{proof}

\begin{remark}\label{rem:choices}
(1) Let $\varepsilon_0$ be the unramified quadratic character
$a(x)\mapsto(-1)^{\nu(x)}$ of $A(F)$. Replacing $\eta_\varpi$ by
$-\eta_\varpi$ (possible when $q\equiv3\pmod4$) replaces the bijection
$\chi\mapsto\overline\chi$ by $\chi\mapsto\overline{\chi\varepsilon_0}$.
Replacing $\varpi$ by $u\varpi$, with $u$ a unit which is not a square,
has the same effect: one finds $\vartheta'=\chi_1(u)\vartheta$ and
$\epsilon'=-\chi_1(u)\epsilon$ for the new parameters. So there are exactly
two bijections of this kind, they differ by the twist by $\varepsilon_0$,
and neither is preferred.

(2) The bijection does not respect unitarity or reducibility of the
induced representations. If $|\chi(a(x))|=|\overline\chi(a(x),1)|=|x|^s$
with $s>0$, the genuine principal series is unitary for $s<1/2$ and
reducible at $s=1/2$ (\S\ref{sec:chars} and the lemmas on complementary
series in Chapter~\ref{ch:reps}), while $I_A(\chi)$ is unitary for $s<1$
and reducible at $s=1$. The global meaning of this is discussed in
Remark~\ref{rem:even}.
\end{remark}

\begin{example}
Let $F={\mathbb Q}_3$, $\varpi=3$, and $\chi_1=1$. Then
$(3,3)_3=(3,-1)_3=-1$, so $\vartheta=\pm1$ while $\epsilon=\pm i$. With
$\eta_3=i$, the trivial character $\chi_{1}$ of $A({\mathbb Q}_3)$
corresponds to $\overline{\chi}_i$, and the character $\chi_{-1}$
($a(x)\mapsto(-1)^{\nu(x)}$) corresponds to $\overline{\chi}_{-i}$. For
$F={\mathbb Q}_5$, $(5,5)_5=1$, $\eta_5=1$, and the definition reduces to
that of the 1995 draft.
\end{example}

\subsubsection{Corresponding functions on \texorpdfstring{$\overline{A(F)}_{reg}$}{A(F)-bar_reg}}

We now determine the conditions on
$\phi \in C_c^\infty(A(F))$,
$\overline{\phi}
\in C_c^\infty(\overline{A(F)})_{gen}$
under which
$\phi^{\wedge}(\chi)=
\overline{\phi}^{\wedge}(\overline{\chi})$ for all $\chi \in \Pi(A(F))$.
Let $\chi=\chi_\vartheta$ extend $\chi_1$. Then
\begin{equation}\label{eq:FT-lin}
\phi^{\wedge}(\chi)=
\int_{A(F)}\phi(a)\chi(a)\, da
=\int_{A_2(F)}\phi(a)\chi_1(x_a)\, da +
\vartheta \int_{A(F)-A_2(F)}\phi(a)\chi_1(x_a/\varpi)\, da .
\end{equation}
Since $\overline{\phi}\,\overline{\chi}$ is invariant under
$\mu_2$ and $d\varsigma$ gives $\mu_2$ volume $1$, and since
$\overline{\chi}=\overline{\chi}_{\epsilon}$ with
$\epsilon=\eta_\varpi\vartheta$,
\begin{equation}\label{eq:FT-gen}
\begin{aligned}
\overline{\phi}^{\wedge}(\overline{\chi})&=
\int_{A(F)}\overline{\phi}(a,1)\overline{\chi}(a,1)\, da\\
&=\int_{A_2(F)}\overline{\phi}(a,1)\chi_1(x_a)\, da +
\eta_\varpi\vartheta \int_{A(F)-A_2(F)}\overline{\phi}(a,1)\,(x_a,\varpi)^{-1}
\chi_1(x_a/\varpi)\, da.
\end{aligned}
\end{equation}

\begin{definition}\label{def:corrA} We say that
$\phi \in C_c^\infty(A(F))$ {\bf corresponds} to
$\overline{\phi}
\in C_c^\infty(\overline{A(F)})_{gen}$
if
\[
\overline{\phi}(a,1)=
\left\{
\begin{array}{ll}
\phi(a), & a\in A_2(F),\\
\eta_\varpi^{-1}(x_a,\varpi)\,\phi(a), & a\in A(F)-A_2(F).
\end{array}
\right.
\]
\end{definition}

\begin{lemma}\label{lem:FT-match}
(a) If $\phi \in C_c^\infty(A(F))$ corresponds to
$\overline{\phi}\in C_c^\infty(\overline{A(F)})_{gen}$
then
$\phi^{\wedge}(\chi)=
\overline{\phi}^{\wedge}(\overline{\chi})$
for all $\chi \in \Pi(A(F))$. Conversely, if this holds for all $\chi$,
then $\phi$ corresponds to $\overline{\phi}$.

(b) If $\phi$ corresponds to $\overline{\phi}$ and $\phi(a^{-1})=\phi(a)$
for all $a$, then $\overline{\phi}((a,1)^w)=\overline{\phi}(a,1)$ for all
$a$.
\end{lemma}

\begin{proof}
(a) The first terms of (\ref{eq:FT-lin}) and (\ref{eq:FT-gen}) agree
when $\overline{\phi}(a,1)=\phi(a)$ on $A_2(F)$. In the second term of
(\ref{eq:FT-gen}), substituting
$\overline{\phi}(a,1)=\eta_\varpi^{-1}(x_a,\varpi)\phi(a)$ gives
$\vartheta\int\phi(a)(x_a,\varpi)^2\chi_1(x_a/\varpi)\,da$, which is the
second term of (\ref{eq:FT-lin}). Conversely, as $\chi_1$ ranges over
$\Pi(A_2(F))$ and $\vartheta$ over both square roots of
$\chi_1(\varpi^2)$, Fourier inversion on $A_2(F)$ and on its coset
$A(F)-A_2(F)$ forces both pointwise identities.

(b) By Lemma~\ref{lem:Acocycle}(b), $(a,1)^w=(a^{-1},1)$, and
$(x_a^{-1},\varpi)=(x_a,\varpi)$.
\end{proof}

\begin{definition} We say that
$f \in C_c^\infty(G(F))$ {\bf corresponds oddly
on the hyperbolic set} to
$\overline{f}\in C_c^\infty(\overline{G(F)})_{gen}$
if the normalized orbital integral $F_f$ on $A(F)_{reg}$ corresponds to
the normalized orbital integral $F_{\overline{f}}$ on
$\overline{A(F)}_{reg}$ in the sense of Definition~\ref{def:corrA}.
\end{definition}

By Lemma~\ref{lem:FT-match}(b), this definition is compatible with the
invariance of $F_f$ and $F_{\overline f}$ under the Weyl group.

\begin{definition}
We say $f,\overline{f}$ {\bf correspond evenly on the hyperbolic set}
if $\Delta(g)O_{\overline{f}}(g)=\Delta(h)O_f(h)$, when
$g=t(h)$ is regular, and $\Delta(g)O_{\overline{f}}(g,1)=0$,
if $g\notin {\rm image}(t)$. Here $O_f$ is the orbital integral
\[
O_f(x)=\int_{G_x(F)\backslash G(F)} f(g^{-1}xg)dg,
\]
and $\Delta$, $t$ are as in \cite{Fl}. (This is the definition of the
1995 draft. The corresponding matching of representations, as in
\cite{Fl}, is not verified here.)
\end{definition}

\begin{remyy}[1995 draft; false as stated]\label{lem:weighted-small}
For $f$ of sufficiently small support, if
$f \in C_c^\infty(G(F))$ corresponds oddly on $A(F)_{reg}$ to
$\overline{f}\in C_c^\infty(\overline{G(F)})_{gen}$
then the weighted orbital integral of
$f$ corresponds to the weighted orbital integral of
$\overline{f}$ in the sense above (as functions on $A(F)$).
\end{remyy}

\begin{remark}\label{rem:weighted-noninv}
The statement cannot hold for \emph{all} corresponding pairs. Indeed, if
$(f,\overline{f})$ correspond, so does $(f,\overline{f}+\overline{h})$
for any genuine $\overline{h}$ whose orbital integrals vanish, for
example $\overline{h}=\overline{g}-\overline{g}^{\,y}$ with
$\overline{g}^{\,y}(x)=\overline{g}(yxy^{-1})$ and $y\in\overline{G(F)}$.
It would then force $J_A(a,\overline{h})=0$, i.e.,
\[
|D(a)|^{1/2}\int_{\overline{A(F)}\backslash\overline{G(F)}}
\overline{g}(x^{-1}ax)\,\bigl(v_A(x)-v_A(xy)\bigr)\,dx=0
\]
for all small-support $\overline{g}$ and all $y$, which fails because
weighted orbital integrals are not invariant distributions
\cite{A7}. A correct statement must either restrict to pairs constructed
in a specific way from one another, or replace the weighted orbital
integrals $J_A$ by the invariant distributions $I_A$ of the invariant
trace formula \cite{A6}. The same objection applies to the remark on
real weighted orbital integrals in \S\ref{sec:realA}. Neither statement
is used in Chapter~\ref{ch:global}.
\end{remark}

\begin{proposition}\label{prop:hyp-ps}
Suppose that $\overline{f}\in C^\infty_c(\overline{G(F)})_{gen}$
corresponds oddly on $A(F)_{reg}$ to
$f\in C_c^\infty(G(F))$. Let $\chi\in\Pi(A(F))$ and let $\overline{\chi}$
be the corresponding genuine character
(Definition~\ref{def:corrchar}). Then
\[
\tr I_A(\overline{\chi})(\overline{f})=\tr I_A(\chi)(f).
\]
\end{proposition}

\begin{proof}
By \S\ref{sec:ps-char},
\[
\tr I_A(\chi)(f)=(F_f)^\wedge(\chi),\qquad
\tr I_A(\overline{\chi})(\overline{f})=(F_{\overline{f}})^\wedge(\overline{\chi}).
\]
These agree by Lemma~\ref{lem:FT-match}(a). (The metaplectic character
formula of \S\ref{sec:ps-char} is quoted there with its proof omitted,
so the proposition depends on it.)
\end{proof}

\begin{example}
Let $F={\mathbb Q}_3$, $\varpi=3$, $\eta_3=i$, and let $\phi$ be the
characteristic function of $\{a(x)\ |\ x\in 3{\mathbb Z}_3^\times\}$.
Then the corresponding $\overline\phi$ has
$\overline{\phi}(a(3u),1)=-i\,(3u,3)_3=-i\,(3,3)_3(u,3)_3=i\left(\frac{u}{3}\right)$
for $u\in{\mathbb Z}_3^\times$, where $(\frac{\cdot}{3})$ is the Legendre
symbol, and $\overline\phi(a(x),1)=0$ otherwise. For the trivial character
$\chi=\chi_1$ (with $\chi_1=1$, $\vartheta=1$), $\phi^\wedge(\chi)={\rm vol}(3{\mathbb Z}_3^\times)$,
and $\overline{\phi}^\wedge(\overline\chi)
=i\int i\,(\tfrac{u}{3})(3u,3)_3\,da=-\int(\tfrac{u}{3})\cdot(-(\tfrac{u}{3}))\,da
={\rm vol}(3{\mathbb Z}_3^\times)$, as it should be.
\end{example}

A similar result holds in the real case, which is discussed in the
next subsection.

\subsection{Corresponding functions on \texorpdfstring{$\overline{A({\mathbb{R}})}$}{A(R)-bar}}\label{sec:realA}

This subsection does the same for $F={\mathbb R}$. There a genuine character of the torus is determined by a character of ${\mathbb R}_+^\times$ and a choice of $\eta=\pm i$.


Let $F={\mathbb{R}}$ in this subsection, and let $a(x)$ denote the
diagonal matrix with entries $x,x^{-1}$. The real Hilbert symbol is
\[
(x,y)_2=
\left\{
\begin{array}{ll}
-1, & x<0\ \text{and}\ y<0,\\
1, & \text{otherwise}.
\end{array}
\right.
\]
By
Lemma~\ref{lem:Acocycle}(a), which holds for every local field,
$(a(x),\zeta)(a(y),\zeta')=(a(xy),(x,y)_2\zeta\zeta')$.

Let $\chi_1$ be a character of ${\mathbb{R}}_+^\times$. Define
$\chi$ on $A({\mathbb{R}})$ by
\[
\chi(a(x))=\chi_\epsilon(a(x))
=
\left\{
\begin{array}{ll}
\chi_1(x), & x>0,\\
\epsilon \chi_1(-x), & x<0,
\end{array}
\right.
\]
where $\epsilon =\pm 1$. Define
$\overline{\chi}$ on $\overline{A({\mathbb{R}})}$ by
\[
\overline{\chi}(a(x),\zeta)=\overline{\chi}_\eta(a(x),\zeta)
=
\left\{
\begin{array}{ll}
\zeta\chi_1(x), & x>0,\\
\zeta\eta \chi_1(-x), & x<0,
\end{array}
\right.
\]
where $\eta=\pm i$. This is a genuine character: the only nontrivial
condition, at $x=y=-1$, reads $1=(-1,-1)_2\,\eta^2=-\eta^2$. Every
genuine character of $\overline{A({\mathbb R})}$ is of this form (compare
\cite{GS}). We say that
$\chi_\epsilon$ {\bf corresponds} to $\overline{\chi}_\eta$ if
$\eta =\epsilon i$. This is a bijection
$\chi\mapsto\overline{\chi}$ from $\Pi(A({\mathbb R}))$ onto the genuine
characters of $\overline{A({\mathbb R})}$.

\begin{definition}\label{def:corrA-real}
We say that $\phi \in C_c^\infty(A({\mathbb R}))$ {\bf corresponds} to
$\overline{\phi}\in C_c^\infty(\overline{A({\mathbb R})})_{gen}$
if
$\phi^{\wedge}(\chi)=
\overline{\phi}^{\wedge}(\overline{\chi})$
for all $\chi \in \Pi(A({\mathbb R}))$.
\end{definition}

\begin{lemma}\label{lem:realA-match}
$\phi$ corresponds to $\overline{\phi}$ if and only if
\begin{equation}\label{eq:realA-match}
\overline{\phi}(a(x),1)=
\left\{
\begin{array}{ll}
\phi(a(x)), & x>0,\\
-i\,\phi(a(x)), & x<0.
\end{array}
\right.
\end{equation}
\end{lemma}

\begin{proof}
As in the p-adic case,
$\phi^\wedge(\chi_\epsilon)=\int_{x>0}\phi(a(x))\chi_1(x)\,d^\times x
+\epsilon\int_{x<0}\phi(a(x))\chi_1(-x)\,d^\times x$ and
$\overline{\phi}^\wedge(\overline{\chi}_{\epsilon i})
=\int_{x>0}\overline{\phi}(a(x),1)\chi_1(x)\,d^\times x
+\epsilon i\int_{x<0}\overline{\phi}(a(x),1)\chi_1(-x)\,d^\times x$.
These agree for all $\chi_1$ and both values of $\epsilon$ if and only if
$\overline{\phi}=\phi$ on $x>0$ and $i\,\overline{\phi}(a(x),1)=\phi(a(x))$
on $x<0$, i.e.\ $\overline{\phi}(a(x),1)=-i\,\phi(a(x))$.
\end{proof}

(The factor in (\ref{eq:realA-match}) cannot depend on $\epsilon$, since
$\epsilon$ depends on the character.) By Lemma~\ref{lem:Acocycle}(b), $(a,1)^w=(a^{-1},1)$, so the
condition (\ref{eq:realA-match}) is compatible with Weyl group
invariance.

\begin{definition} We say that
$f \in C_c^\infty(G({\mathbb R}))$ {\bf corresponds oddly on the
hyperbolic set} to
$\overline{f}\in C_c^\infty(\overline{G({\mathbb R})})_{gen}$
if $F_f$ corresponds to
$F_{\overline{f}}$
in the sense of Definition~\ref{def:corrA-real}.
\end{definition}

With this definition, the proof of Proposition~\ref{prop:hyp-ps} gives
$\tr I_A(\overline{\chi})(\overline{f})=\tr I_A(\chi)(f)$ for
$f$, $\overline{f}$ corresponding oddly on the hyperbolic set, given the
character formula of \S\ref{sec:ps-char} for ${\mathbb R}$.

\begin{remark}
The 1995 draft added here that, by a direct calculation, weighted
orbital integrals of functions corresponding on
$A({\mathbb{R}})_{reg}$ also correspond. For arbitrary corresponding
pairs this cannot hold, for the reason given in
Remark~\ref{rem:weighted-noninv}.
\end{remark}

\subsection{Corresponding functions on the elliptic set}\label{sec:elliptic}

Finally, we concentrate on the elliptic set. Let ${\cal T}_{ell}$ be a
set of representatives of the $G(F)$-conjugacy classes of anisotropic
(elliptic) Cartan subgroups of $G$ over $F$. If $F={\mathbb R}$, there is
one class, represented by $SO(2)$. If $F$ is p-adic, there are several:
each anisotropic torus is isomorphic to the norm-one torus of a
quadratic extension $E/F$, and for $p$ odd there are three such $E$ and
two $G(F)$-conjugacy classes of tori for each $E$. Matching on a single
elliptic Cartan subgroup is therefore not enough in the p-adic case.

For $T$ a Cartan subgroup, $t\in T(F)$ regular, and
$f\in C_c^\infty(G(F))$, let $\Phi_f(t)=|D(t)|\,O_f(t)$, where
$O_f(t)=\int_{T(F)\backslash G(F)}f(x^{-1}tx)\,dx/dt$, and define
$\Phi_{\overline{f}}$ on $\overline{T(F)}_{reg}$ in the same way for
$\overline{f}\in C_c^\infty(\overline{G(F)})_{gen}$. Here the notation $\Phi_f$ is \emph{not} the usual
square-root-normalized orbital integral $F_f=|D|^{1/2}O_f$
introduced earlier. We deliberately use the full Weyl
discriminant $|D|$ in $\Phi_f$, so that the unnormalized
character $\Theta_\pi(t)$ appears in the Weyl integration formula.
In particular $|D(r(\theta))|=4\sin^2\theta$
when $F=\mathbb R$. With this
normalization (and the measures of \S\ref{sec:local}), the Weyl
integration formula reads
\[
\Theta _{\pi}(f)=\sum_{T\in{\cal T}}\frac 1{|W_T|}\int_{T(F)}\Theta _{\pi}(t)
\Phi_f(t)\,dt,\;\;\;f\in C_c^\infty (G(F)),
\]
and
\[
\Theta _{\overline{\pi } }(\overline{f})
=\sum_{T\in{\cal T}}\frac 1{|W_T|}\int_{\overline{T(F)}}
\Theta _{\overline{\pi } }(t)
\Phi_{\overline{f}}(t)\,dt\,d\varsigma,
\;\;\;\overline{f}\in C_c^\infty (\overline{G(F)})_{gen},
\]
where ${\cal T}$ consists of $A$ and the elements of ${\cal T}_{ell}$,
and $W_T=N_{G(F)}(T)/T(F)$. Since $\Theta_{\overline{\pi}}$ and
$\Phi_{\overline{f}}$ are both genuine, the integrand on
$\overline{T(F)}$ is invariant under $\mu_2$, and the integral over
$\overline{T(F)}$ equals $\int_{T(F)}\Theta_{\overline\pi}(t,1)\Phi_{\overline f}(t,1)\,dt$.

\begin{definition}\label{def:ell-corr}
We say that $f$ and $\overline{f}$ {\bf have vanishing hyperbolic
orbital integrals} if $\Phi_f=0$ on $A(F)_{reg}$ and
$\Phi_{\overline f}=0$ on $\overline{A(F)}_{reg}$.

(a) If $f,\overline{f}$ have vanishing hyperbolic orbital integrals and
satisfy
\[
\Phi_f(t)=\Phi_{\overline{f}}(t,1),\ \ \ \ t\in T(F)_{reg},\ \text{for every } T\in{\cal T}_{ell},
\]
then we say that they \textbf{correspond oddly on the elliptic set}.

(b) If $f,\overline{f}$ have vanishing hyperbolic orbital integrals and
satisfy, for every $T\in{\cal T}_{ell}$,
\[
\Phi_{\overline{f}}(g)
=
\left\{
\begin{array}{ll}
\Phi_f(h), & g=\tau(h),\ h\in T(F)_{reg},\\
0, & g\notin \tau(T(F)),
\end{array}
\right.
\]
then we say that they \textbf{correspond evenly on the elliptic set}.
Here $\tau:h\longmapsto (h,1)^2$ is the map denoted
$h\longmapsto \tilde{h}$ in \cite{Fl}.
\end{definition}

For instance, if $f$ and $\overline{f}|_{G(F)}$ are supported in the
(open) set of regular elliptic elements, they have vanishing hyperbolic
orbital integrals.

\begin{lemma}\label{lem:ell-odd}
Let $f,\overline{f}$ correspond oddly on the elliptic set. Let $\pi$ and
$\overline{\pi}$ be an admissible representation of $G(F)$ and a genuine
admissible representation of $\overline{G(F)}$ whose characters satisfy
\begin{equation}\label{eq:ell-char-rel}
\Theta_{\overline{\pi}}(t,\zeta)=\zeta\,\Theta_\pi(t),\qquad
t\in T(F)_{reg},\ \zeta\in\mu_2,\ \text{for every }T\in{\cal T}_{ell}.
\end{equation}
Then $\Theta _{\pi }(f)=\Theta _{\overline{\pi }}(\overline{f})$.
\end{lemma}

\begin{proof}
In both Weyl integration formulas the term for $T=A$ vanishes, because
the hyperbolic orbital integrals vanish. For $T\in{\cal T}_{ell}$, the
term for the cover equals
$|W_T|^{-1}\int_{T(F)}\Theta_{\overline\pi}(t,1)\Phi_{\overline f}(t,1)\,dt
=|W_T|^{-1}\int_{T(F)}\Theta_{\pi}(t)\Phi_{f}(t)\,dt$ by
(\ref{eq:ell-char-rel}) and Definition~\ref{def:ell-corr}(a), which is
the term for $G$. Summing over ${\cal T}_{ell}$ gives the result.
\end{proof}

The character relation (\ref{eq:ell-char-rel}) for supercuspidal
representations is Remark~\ref{prop:164}, which needs correction
(Remark~\ref{rem:sc-caveat}). Lemma~\ref{lem:ell-odd} is therefore
stated with (\ref{eq:ell-char-rel}) as a hypothesis. The condition
$\Phi_f(t)=\Phi_{\overline f}(t,1)$ and the relation
(\ref{eq:ell-char-rel}) both use the section $t\mapsto(t,1)$, which is
not compatible with conjugation. \S\ref{sec:cover-factor} gives a form of
both that is, with the function $\Delta_\psi$ in place of the section
(Definition~\ref{def:direct}), and \S\ref{sec:cover-elliptic} shows what
happens to it on the elliptic set.

\subsection{Corresponding spherical functions, \texorpdfstring{$p> 2$}{p> 2}}\label{sec:spherical}

We sketch the Satake correspondence for $\overline{G(F)}$
following \cite{J}, \S 1.3 and using its notation.

Recall that the genuine characters of
$\overline{A(F)}$ are in 1-1 correspondence with
the characters of $A(F)$. Moreover, from results recalled in
Chapter~\ref{ch:reps} we deduce that all
the unramified genuine representations of
$\overline{G(F)}$ belong to the unramified principal series.
Replacing $GL(r)$ by
$SL(2)$ in \cite{J}, \S 1.3, we first observe that the
set $\Pi_{unr.p.s.}(\overline{G(F)})_{gen}$ is
parametrized by the $W$-orbit of unramified
genuine characters of $\overline{A(F)}$. Therefore, we have the following

\begin{lemma} The set $\Pi_{unr.p.s.}(\overline{G(F)})_{gen}$ is
in 1-1 correspondence with
$\Pi_{unr.p.s.}(G(F))$, by $I_A(\overline\chi)\leftrightarrow I_A(\chi)$
for corresponding unramified characters. The correspondence is compatible
with the Weyl group (Lemma~\ref{lem:Wequiv}) and depends on the choices in
Remark~\ref{rem:choices}.
\end{lemma}

Following now \cite{J}, \S 1.4, we obtain a transform
\begin{equation}
\begin{array}{c}
{\cal{H}}(\overline{G(F)},K)_{gen}\rightarrow
{\cal{A}}_{cl.fcn}(PGL(2,{\mathbb{C}}))\\
f\longmapsto S_0 f(\chi)=\tr I_A(\chi)(f).
\end{array}
\end{equation}

The target $\mathcal A_{cl.fcn}(PGL(2,\mathbb C))$
in this \emph{direct-comparison} transform is obtained by
transporting the ordinary $SL(2)$ Satake parametrization
through the character bijection of Lemma~\ref{lem:genuine-chars}.
It must not be read as asserting that the intrinsic unramified
$L$-group of the metaplectic cover is $PGL(2,\mathbb C)$:
the genuine parameter in the Shimura--Waldspurger correspondence
belongs naturally to $SL(2,\mathbb C)$, as discussed in
\S\ref{sec:where-fails}.

This is essentially the inverse Fourier transform of the Satake
transform, which we define next. The Satake transform of a
genuine spherical function is defined as in the non-metaplectic case:
\begin{equation}
Sf(a)=\delta(a)^{1/2}\int_{N(F)}f(an)dn.
\end{equation}
Here $a\in\overline{A(F)}$ and $n\in N(F)$ is lifted by
$n\mapsto(n,1)$; the integral uses the fixed splitting of the
unipotent subgroup.
The relationship between this expression and an orbital integral
is well-known - see for example \cite{C}.

This defines an algebra isomorphism
\begin{equation}
\begin{array}{c}
{\cal{H}}(\overline{G(F)},K)_{gen}\rightarrow
{\cal{H}}(\overline{A(F)},\kappa(K\cap A(F)))^W_{gen}\\
f\longmapsto Sf(a).
\end{array}
\end{equation}

By Lemma~\ref{lem:FT-match}, the correspondence $\phi\mapsto\overline\phi$
of Definition~\ref{def:corrA} preserves Fourier transforms, hence
convolution. It maps $W$-invariant functions to $W$-invariant functions.
And, since the splitting $\kappa$ is $a(u)\mapsto(a(u),1)$ on $K\cap A(F)$
(by (\ref{eq:sv}), as $c=0$), it maps functions invariant under
$K\cap A(F)$ to genuine functions invariant under $\kappa(K\cap A(F))$
(the factor $(x_a,\varpi)$ transforms correctly under $a\mapsto a\,a(u)$
by Lemma~\ref{lem:Acocycle}(a)). It therefore gives an algebra isomorphism
${\cal{H}}(A(F),K\cap A(F))^W\cong
{\cal{H}}(\overline{A(F)},\kappa(K\cap A(F)))^W_{gen}$, and so
\[
{\cal{H}}(G(F),K)\cong{\cal{H}}(\overline{G(F)},K)_{gen}.
\]

\begin{definition}\label{def:satake-corr}
We say that $f\in{\cal H}(G(F),K)$ and
$\overline f\in{\cal H}(\overline{G(F)},K)_{gen}$ {\bf correspond under
the Satake isomorphism} if $Sf$ corresponds to $S\overline{f}$ in the
sense of Definition~\ref{def:corrA}.
\end{definition}

If $f$ and $\overline f$ correspond under the Satake isomorphism and $\chi$ is
an unramified character of $A(F)$, then
$\tr I_A(\overline\chi)(\overline f)=(S\overline f)^\wedge(\overline\chi)
=(Sf)^\wedge(\chi)=\tr I_A(\chi)(f)$, by Lemma~\ref{lem:FT-match}(a).
This is the \emph{spectral} matching supplied by the Satake isomorphism.

\begin{remyy}\label{rem:satake-oi}
The 1995 draft concluded: ``This implies the matching of orbital
integrals for spherical functions.'' That does not follow from the
existence of an algebra isomorphism alone. What does follow is the
following. For $f\in{\cal H}(G(F),K)$ and regular $a\in A(F)$, the
normalized hyperbolic orbital integral $F_f(a)$ equals the Satake
transform $Sf(a)$ \cite{C}. Granted the analogous identity on the cover
(the same computation, using the splitting $\kappa$ and
Lemma~\ref{lem:facts}(h). The cocycle bookkeeping has not been written out
here), functions corresponding under the Satake isomorphism correspond
oddly on the hyperbolic set. The \emph{elliptic} orbital integrals of
corresponding spherical functions are not determined by their Satake
transforms in any formal way, and their matching (a fundamental lemma
for this comparison) is not established here. It is one of the
ingredients assumed in Hypothesis~\ref{hyp:GH}.
\end{remyy}

\begin{example}\label{ex:units}
For the unit elements $f=1_K$ and $\overline f=e_K$ of the two Hecke
algebras, $Sf=1_{K\cap A(F)}$ and $S\overline f$ is supported on
$\overline{A_2(F)}$, where it equals $(a,\zeta)\mapsto\zeta\,1_{K\cap A(F)}(a)$.
These correspond (the factor $\eta_\varpi^{-1}(x_a,\varpi)$ only enters off
$A_2(F)$). Hence $\tr I_A(\overline\chi)(e_K)=\tr I_A(\chi)(1_K)=1$ for
unramified $\chi$.
\end{example}

\begin{remyy}\label{rem:weighted-units}
On the matching of
weighted orbital integrals of spherical functions:
every global smooth genuine function
$\overline{f}=\otimes_v\overline{f}_v$ is a restricted tensor product
of local smooth functions, almost all of which are units of the
local Hecke algebra of genuine spherical functions on
$\overline{G(k_v)}$. In order to match weighted orbital integrals
of such global functions, one must match them locally.
The 1995 draft states that a direct calculation of the weighted orbital
integrals shows that the weighted orbital integral of the unit element
$\overline{f}$ of the genuine spherical Hecke algebra on $\overline{A(F)}$
matches that of the unit element $f$ of the spherical Hecke algebra of
$G(F)$ on $A(F)$. That calculation is not reproduced here. (This is a
statement about a specific pair, so the objection of
Remark~\ref{rem:weighted-noninv} does not apply to it.)
\end{remyy}

\begin{lemma}
If $f,\overline{f}$ are the unit elements
of the respective Hecke algebras, $p\not= 2$, then
$\Theta_\pi(f)=\Theta_{\overline{\pi}}(\overline{f})$,
where
$\pi \leftrightarrow \overline{\pi}$ are corresponding
unramified principal series as above.
\end{lemma}

This is the last statement of Example~\ref{ex:units}.

\section{The dyadic places}\label{sec:dyadic}

This section explains why the dyadic places are excluded from the local results of this chapter.


In this section $F$ is a dyadic field, i.e., $p=2$, and we write
$G=SL(2,F)$ and $\overline{G}=\overline{SL(2,F)}$.

An earlier draft attempted to obtain a matching at dyadic places from a
correspondence of generalized Hecke algebras in the sense of Howe and Moy
\cite[ch.~3]{HM}. That argument is withdrawn. The reasons are recorded in
the revision log, and the Hecke algebra material is revised separately in
\cite{J3}. Two features of the dyadic case explain the difficulty. The
affine Weyl group of $SL(2)$ is infinite dihedral, so there is no braid
relation between its two simple reflections. And the set of elements over
which the cocycle $\beta$ is trivial is not a subgroup, so the genuine
functions supported over it do not form a convolution algebra.

In this version no local matching of functions or of characters is
established at the dyadic places. Instead, the dyadic places are always
put in the exceptional set $S$ of Chapter~\ref{ch:global}, and what is
needed there is assumed in Hypotheses~\ref{hyp:GH} and~\ref{hyp:LH}.

\begin{remark}
(1) The construction of \S\ref{sec:charsA} does not extend to $p=2$
without change: there $\overline{A_2(F)}$ need not be the direct product
$A_2(F)\times\mu_2$ (see the example at the end of this section).

(2) At Iwahori level, there are Hecke algebra correspondences at
$p=2$ for the metaplectic group (Takeda and Wood \cite{TW}, extending
work of Gan and Savin for $p\neq 2$), but they relate genuine
representations of $\overline{SL(2,F)}$ to representations of
$SO(3)\cong PGL(2)$ and its inner form, not of $SL(2)$. Whether such
results can be used to verify Hypothesis~\ref{hyp:LH} at dyadic places
is not addressed here.
\end{remark}

%% (aabook2b.tex withdrawn in Draft 1; see revision_log.md)



%% ===== begin aabook2c.tex =====
\begin{example}
Let $F={\mathbb Q}_2$. Since $x(-I)=-1$, the cocycle gives
$\beta(-I,-I)=(-1,-1)_2=-1$, so $(-I,1)^2=(I,-1)$. Thus the lifts of
$-I\in A_2(F)$ have order $4$, and $\overline{A_2(F)}$ is not the direct
product $A_2(F)\times\mu_2$, unlike the case $p>2$ of \S\ref{sec:charsA}.
\end{example}

\section{On the correspondence of limits of d.s. for
\texorpdfstring{$\overline{SL(2,\mathbb{R})}$}{SL(2,R)-bar}}\label{sec:realds}

This section constructs three matchings of functions on $SL(2,{\mathbb R})$ and on its cover whose hyperbolic orbital integrals vanish. It shows that each one transfers genuine discrete series to discrete series or limits of discrete series.


Let $T$ denote a Cartan of $G=SL(2)$.
%$F$ be a field and
%$\overline{F}$ a separable algebraic closure. By the density
%theorem (see for example, [Hum], \S 22.2), the conjugate $T^G$ contains
%a $\overline{F}$-Zariski-open dense subset of $G$.
%The $F$-rational
%points $T^G(F)$ remain $F$-Zariski open in $G(F)$ (but not necessarily dense,
%see for example, Theorem 1 of ch II, \S 4, in [Mum]). If $F$ is a
%complete local field then the points $T^G(F)$ are open (in the induced
%$F$-manifold topology) in $G(F)$. However, in general,
%$T(F)^{G(F)}\not= T^G(F)$.
If $F={\mathbb{R}}$ and $T$ is an
anisotropic Cartan of $G$ defined over $F$ then recall $T(F)^{G(F)}$ is the
elliptic set of $G$.

We shall describe 3 different correspondences
$c_i:\overline{\pi}\longmapsto \pi$, $i=1,2,3$, between representations of
$\overline{SL(2,\mathbb{R})}$ and those of $SL(2,\mathbb{R})$. They are
consistent with a ``matching'' $c_i^*:\overline{f}\longmapsto f$,
$i=1,2,3$, of functions whose hyperbolic orbital integrals vanish, in the
sense that the characters match:
$\Theta_{\overline{\pi}}(\overline{f})=\Theta_{\pi}(f)$, where
$\pi=c_i(\overline{\pi})$, $f=c_i^*(\overline{f})$, $i=1,2,3$
(Lemma~\ref{lem:realds}).
Furthermore, if $\pi$ is any element of the discrete series or limit of
discrete series of $G(F)$ then $\pi$ occurs in the image of at least one
of the $c_i$.

\subsection{Discrete series for \texorpdfstring{$\overline{SL(2,\mathbb{R})}$}{SL(2,R)-bar}}

We recall the genuine principal series and discrete series of the real cover.


Let
\begin{align*}
\overline{B} &=\overline{M}\cdot A(\mathbb{R})^o\cdot N(\mathbb{R}), \\
A(\mathbb{R})^o &=\{a(x)=\left(
\begin{array}{cc}
x & 0 \\
0 & x^{-1}
\end{array}
\right) \;|\;x>0\}, \\
N(\mathbb{R}) &=\{n=\left(
\begin{array}{cc}
1 & x \\
0 & 1
\end{array}
\right) \;|\;x\in \mathbb{R}\}, \\
\overline{M} &=\langle\gamma\rangle,\;\gamma =(-\left(
\begin{array}{cc}
1 & 0 \\
0 & 1
\end{array}
\right) ,1),
\end{align*}
and let
\[
\overline{r}(\theta )=\left(
\begin{array}{cc}
\cos (\theta ) & -\sin (\theta ) \\
\sin (\theta ) & \cos (\theta )
\end{array}
\right) ,\;\;0\leq \theta \leq 4\pi ,
\]
which we use to denote the element of $\overline{K}$ lying over the
displayed rotation, chosen so that $\theta\mapsto\overline r(\theta)$ is a
continuous homomorphism ${\mathbb R}\rightarrow\overline K$ with
$\overline r(0)$ the identity. Then $\overline r(2\pi)=(I,-1)$, since
$\overline K$ is connected (Lemma~\ref{lem:facts}(d)).
Let $\overline{\chi }=\overline{\chi }_{\lambda ,\tau }$ denote the
character of $\overline{B}$ given by
\[
\overline{\chi }(\gamma ^ja(x)n)=e^{ij\tau \pi }x^\lambda ,\;\;\tau \in
\{0,1,1/2,-1/2\},\ \ \ \ \ \ x>0.
\]
This character is genuine if and only if $\tau \in \{1/2,-1/2\}$. Every
genuine principal series is given by
\[
\rho (\lambda ,\tau )=\operatorname{Ind}(\overline{\chi },\overline{B},\overline{SL(2,\mathbb{R%
})}).
\]

Let $\Pi (\overline{SL(2,\mathbb{R})})_{ds}$ denote the set of equivalence
classes of genuine unitary discrete series representations of $\overline{%
SL(2,\mathbb{R})}$.

\subsection{\texorpdfstring{Gelbart's parameterization \cite[p.~75]{Ge2}}{Gelbart's parameterization}}

Gelbart's tables list the genuine discrete series by their parameters.


Let $k>1$ be an odd integer and let $\tau \in \{1/2,-1/2\}$. Let
\[
\phi _m(na(x)\overline{r}(\theta ))=x^\lambda e^{im\theta }=x^\lambda
e^{i(2s+\tau )\theta },\;\;\;\lambda \in \mathbb{C},\;\;m=2s+\tau ,\;\;s\in
\mathbb{Z}.
\]

These form a collection of functions which can be used to describe a basis
for the discrete series. The discrete series are tabulated as follows:
\[
\begin{array}{cccc}
\mathrm{d.s.\;repn} & k & \tau & \lambda \\
B_{-}^{1/2}(k/2) & k\equiv 1(4) & 1/2 & \frac k2-1 \\
B_{+}^{-1/2}(-k/2) & k\equiv 1(4) & -1/2 & \frac k2-1 \\
B_{+}^{1/2}(-k/2) & k\equiv 3(4) & 1/2 & \frac k2-1 \\
B_{-}^{-1/2}(k/2) & k\equiv 3(4) & -1/2 & \frac k2-1
\end{array}
\]
These representations exhaust $\Pi (\overline{SL(2,\mathbb{R})})_{ds}$.
Gelbart also gives a second realization, with the sign of
$\lambda$ reversed. We retain it for comparison of conventions;
the later character formulas use the corresponding
$\overline K$-type notation $\overline\pi_{k/2}^{\pm}$ instead
of the labels $B_{\pm}^{\pm1/2}$. The second table is:
\[
\begin{array}{cccc}
\mathrm{d.s.\;repn} & k & \tau & \lambda \\
B_{-}^{1/2}(-k/2+2) & k\equiv 3(4) & 1/2 & -\frac k2+1 \\
B_{+}^{-1/2}(k/2-2) & k\equiv 3(4) & -1/2 & -\frac k2+1 \\
B_{+}^{1/2}(k/2-2) & k\equiv 1(4) & 1/2 & -\frac k2+1 \\
B_{-}^{-1/2}(-k/2+2) & k\equiv 1(4) & -1/2 & -\frac k2+1
\end{array}
\]

\subsection{Discrete series for \texorpdfstring{$SL(2,\mathbb{R})$}{SL(2,R)}}

We recall Harish-Chandra's formulas for the characters of the discrete series of $SL(2,{\mathbb R})$.


Let
\begin{align*}
B &=M\cdot A(\mathbb{R})^o\cdot N(\mathbb{R}), \\
M &=\{\pm I_2\},\;I_2=\left(
\begin{array}{cc}
1 & 0 \\
0 & 1
\end{array}
\right) ,
\end{align*}
and let
\[
r(\theta )=\left(
\begin{array}{cc}
\cos (\theta ) & -\sin (\theta ) \\
\sin (\theta ) & \cos (\theta )
\end{array}
\right) ,\;\;0\leq \theta \leq 2\pi .
\]

For each non-zero $m\in \mathbb{Z}$ there is a $\pi _m\in \Pi (SL(2,\mathbb{R}%
))_{ds}$ whose character $\Theta _{\pi _m}$ satisfies
\[
\Theta _{\pi _m}(r(\theta ))=-\sgn(m)
\frac{e^{im\theta }}{e^{i\theta }-e^{-i\theta }},
\;\;\;0<\theta <2\pi ,\;\theta \neq \pi ,
\]
and
\[
\Theta _{\pi _m}((-I_2)^j a(e^t ))=(-1)^{(m-1)j}
\frac{e^{-|mt| }}{|e^t-e^{-t}|}.
\]

For $m$ odd, $\pi _m$ has trivial central character. For $m$ even, $m\neq 0$
($m=0$ corresponds to the two irreducible components of the unitary
reducible principal series), $\pi _m$ has non-trivial central character.
This may be found in \cite{HC2}, page 660.

We denote the two irreducible components of the unitary reducible principal
series by $\pi _0^{+}$ and $\pi _0^{-}$, labeled in such a way that their
character $\Theta _{\pi _0^{\pm }}$ satisfies
\begin{align*}
\Theta _{\pi _0^\varepsilon }(r(\theta )) &=-\varepsilon \frac 1{e^{i\theta
}-e^{-i\theta }},\;\;\;0<\theta <2\pi ,\;\theta \neq \pi , \\
\Theta _{\pi _0^\varepsilon }((-I_2)^ja(x)) &=(-1)^j\frac
1{|x-x^{-1}|},\;\;\;\varepsilon =\pm .
\end{align*}

\subsection{A correspondence between the discrete series}

We now compare the elliptic terms of the two Weyl integration formulas, and deduce the three matchings.


\emph{Labels and characters.} For $k\geq 3$ odd, let
$\overline{\pi}{}^{+}_{k/2}$ denote the genuine holomorphic discrete series
representation of $\overline{SL(2,{\mathbb R})}$ with lowest
$\overline{K}$-type $\overline{r}(\theta)\mapsto e^{ik\theta/2}$ (its
$\overline K$-types are $e^{i(k/2+2j)\theta}$, $j\geq0$), and
$\overline{\pi}{}^{-}_{k/2}$ the antiholomorphic one, with
$\overline K$-types $e^{-i(k/2+2j)\theta}$. On the regular elliptic set
their characters are given by the $\overline K$-characters, summed as
geometric series (Harish-Chandra; compare \cite{Ad2}):
\begin{equation}\label{eq:mds-char}
\Theta_{\overline{\pi}{}^{+}_{k/2}}(\overline r(\theta))
=-\frac{e^{i(k/2-1)\theta}}{e^{i\theta}-e^{-i\theta}},\qquad
\Theta_{\overline{\pi}{}^{-}_{k/2}}(\overline r(\theta))
=\frac{e^{-i(k/2-1)\theta}}{e^{i\theta}-e^{-i\theta}},\qquad
\theta\notin\pi{\mathbb Z}.
\end{equation}
The same computation for $SL(2,{\mathbb R})$ gives the formulas for
$\Theta_{\pi_m}$ and $\Theta_{\pi_0^\pm}$ above. In particular $\pi_m$,
$m>0$, is holomorphic with lowest $K$-type $m+1$, and $\pi_0^+$ is the
holomorphic limit of discrete series. Here $\varepsilon=+$ always means
holomorphic (see Remark~\ref{rem26:ds-labels}).

\emph{Weyl integration formula.} With $\Phi$ normalized as in
\S\ref{sec:elliptic}, $dk=d\theta/2\pi$ on $K=\{r(\theta)\}$ and
$d\overline k=d\theta/4\pi$ on $\overline K=\{\overline r(\theta)\}$,
the elliptic terms of the Weyl integration formulas for $\Theta_{\pi}(f)$
and $\Theta_{\overline\pi}(\overline f)$ are
\begin{gather*}
\int_0^{2\pi }\Theta _{\pi}(r(\theta ))\Phi_f(r(\theta ))\,\frac{d\theta}{2\pi}
\qquad\text{and}\\
\int_0^{4\pi }
\Theta _{\overline{\pi }}(\overline{r}(\varphi ))
\Phi_{\overline{f}}(\overline{r}(\varphi ))\,\frac{d\varphi}{4\pi}
=\int_0^{2\pi }
\Theta _{\overline{\pi }}(\overline{r}(2\theta ))
\Phi_{\overline{f}}(\overline{r}(2\theta ))\,\frac{d\theta}{2\pi}.
\end{gather*}
Here the Weyl group of $SO(2)$ in $SL(2,{\mathbb R})$ is trivial. Since
\[
e^{2i\theta}-e^{-2i\theta}=(e^{i\theta}-e^{-i\theta})(e^{i\theta}+e^{-i\theta}),
\]
the substitution $\varphi=2\theta$ in (\ref{eq:mds-char}) produces the
factor $e^{i\theta}+e^{-i\theta}=2\cos\theta$ in the denominator. This
factor vanishes at $\theta=\pi/2,3\pi/2$, where $r(\theta)$ is regular
but $\overline r(2\theta)$ lies over the central element $-I$. We
therefore state the matching conditions without division.

\begin{definition}\label{def:realds-corr}
Let $f,\overline{f}$ have vanishing hyperbolic orbital integrals, and let
$E=\{\pi/2,\pi,3\pi/2\}$. We say that $f,\overline f$
\textbf{correspond evenly}, \textbf{$+$oddly}, or \textbf{$-$oddly}
\textbf{on the elliptic set} if, for all $\theta\in(0,2\pi)-E$,
\begin{align*}
\Phi_{\overline{f}}(\overline{r}(2\theta ))
&=(e^{i\theta }+e^{-i\theta })\,\Phi_f(r(\theta )) &&(\mathrm{even}),\\
\Phi_{\overline{f}}(\overline{r}(2\theta ))
&=e^{-i\theta}(e^{i\theta }+e^{-i\theta })\,\Phi_f(r(\theta )) &&(+\mathrm{odd}),\\
\Phi_{\overline{f}}(\overline{r}(2\theta ))
&=e^{i\theta}(e^{i\theta }+e^{-i\theta })\,\Phi_f(r(\theta )) &&(-\mathrm{odd}),
\end{align*}
respectively.
\end{definition}

\begin{remark}\label{rem:realds-consistency}
(1) As $\theta$ runs over $(0,2\pi)-E$, $\overline r(2\theta)$ runs over
all regular elements of $\overline K$, so each condition determines
$\Phi_{\overline f}$ on the regular elliptic set from $\Phi_f$. Since
$\Phi_f$ is smooth near $r(\pi/2)$ and $r(3\pi/2)$, which are regular,
each condition forces $\Phi_{\overline f}(\overline r(\varphi))\rightarrow0$
as $\varphi\rightarrow\pi$ or $3\pi$. (The factor
$(e^{i\theta}+e^{-i\theta})^{-1}$ is undefined at these points, which is
why the conditions are written without division.)

(2) Since $\overline r(2\pi)=(I,-1)$ and $\Phi_{\overline f}$ is genuine,
replacing $\theta$ by $\theta+\pi$ shows that the even condition can hold
only if $\Phi_f(-t)=\Phi_f(t)$, and the odd conditions only if
$\Phi_f(-t)=-\Phi_f(t)$, for regular elliptic $t$. This is consistent
with Lemma~\ref{lem:realds}: the even matching sees the $\pi_m$ with
trivial central character ($m$ odd), the odd matchings those with
nontrivial central character ($m$ even).
\end{remark}

\begin{example}
Let $\Phi_f(r(\theta))=\sin(\theta)\,g(\theta)$ with $g$ smooth,
$2\pi$-periodic and supported in a small neighborhood of $\theta=\pi/2$
and of $\theta=3\pi/2=\pi/2+\pi$ (with $g(\theta+\pi)=-g(\theta)$, so that
$\Phi_f(-t)=\Phi_f(t)$). The even condition requires
$\Phi_{\overline f}(\overline r(\varphi))=2\cos(\varphi/2)\sin(\varphi/2)g(\varphi/2)
=\sin(\varphi)g(\varphi/2)$, which is smooth and vanishes at
$\varphi=\pi,3\pi$, as Remark~\ref{rem:realds-consistency}(1) requires.
\end{example}

\begin{lemma}\label{lem:realds}
Let $k\geq3$ be odd and $\varepsilon\in\{+,-\}$.

(a) If $f,\overline{f}$ correspond evenly, then
$\Theta _{\pi _m}(f)
=\Theta _{\overline{\pi }{}_{k/2}^\varepsilon }(\overline{f})$ for
$m=\varepsilon (k-2)$.

(b) If $f,\overline{f}$ correspond $+$oddly, then
$\Theta _{\pi _m}(f)
=\Theta _{\overline{\pi }{}_{k/2}^\varepsilon }(\overline{f})$ for
$m=\varepsilon (k-2)-1$, unless $k=3$ and
$\varepsilon =+$, and
$\Theta _{\pi _{0}^{+}}(f)=\Theta _{\overline{\pi }{}_{3/2}^{+}}(\overline{f})$.

(c) If $f,\overline{f}$ correspond $-$oddly, then
$\Theta _{\pi _m}(f)
=\Theta _{\overline{\pi }{}_{k/2}^\varepsilon }(\overline{f})$ for
$m=\varepsilon (k-2)+1$, unless $k=3$ and
$\varepsilon =-$, and
$\Theta _{\pi _{0}^{-}}(f)=\Theta _{\overline{\pi }{}_{3/2}^{-}}(\overline{f})$.
\end{lemma}

\begin{proof}
Since the hyperbolic orbital integrals vanish, only the elliptic terms
of the Weyl integration formulas contribute. By (\ref{eq:mds-char}) and
the factorization above, for $\theta\notin E$,
\[
\Theta_{\overline{\pi}{}^{+}_{k/2}}(\overline r(2\theta))
=-\frac{e^{i(k-2)\theta}}{(e^{i\theta}-e^{-i\theta})(e^{i\theta}+e^{-i\theta})},
\qquad
\Theta_{\overline{\pi}{}^{-}_{k/2}}(\overline r(2\theta))
=\frac{e^{-i(k-2)\theta}}{(e^{i\theta}-e^{-i\theta})(e^{i\theta}+e^{-i\theta})}.
\]
Multiplying by $\Phi_{\overline f}(\overline r(2\theta))$ as given by
Definition~\ref{def:realds-corr}, the factor $e^{i\theta}+e^{-i\theta}$
cancels, and the elliptic integrand of $\Theta_{\overline\pi}(\overline f)$
becomes $\Phi_f(r(\theta))$ times:
\[
\renewcommand{\arraystretch}{1.6}
\begin{array}{c|cc}
 & \overline{\pi}{}^{+}_{k/2} & \overline{\pi}{}^{-}_{k/2}\\ \hline
\text{even} & -\dfrac{e^{i(k-2)\theta}}{e^{i\theta}-e^{-i\theta}}
 & \dfrac{e^{-i(k-2)\theta}}{e^{i\theta}-e^{-i\theta}}\\
+\text{odd} & -\dfrac{e^{i(k-3)\theta}}{e^{i\theta}-e^{-i\theta}}
 & \dfrac{e^{-i(k-1)\theta}}{e^{i\theta}-e^{-i\theta}}\\
-\text{odd} & -\dfrac{e^{i(k-1)\theta}}{e^{i\theta}-e^{-i\theta}}
 & \dfrac{e^{-i(k-3)\theta}}{e^{i\theta}-e^{-i\theta}}
\end{array}
\]
Comparing with
$\Theta_{\pi_m}(r(\theta))=-\sgn(m)e^{im\theta}/(e^{i\theta}-e^{-i\theta})$
and $\Theta_{\pi_0^{\pm}}(r(\theta))=\mp1/(e^{i\theta}-e^{-i\theta})$, the
entries are $\Theta_{\pi_m}(r(\theta))$ with $m=\pm(k-2)$ (even),
$m=k-3$ or $-(k-1)$ ($+$odd), $m=k-1$ or $-(k-3)$ ($-$odd). When $k=3$
the entries $k-3$ and $-(k-3)$ are replaced by $\pi_0^+$ and $\pi_0^-$.
These agree with $\Theta_{\pi_m}(r(\theta))\Phi_f(r(\theta))$ for almost
all $\theta$, so the integrals agree.
\end{proof}

\begin{remyy}\label{rem26:ds-labels}
The labels $\overline\pi{}^{\pm}_{k/2}$ above are fixed by
(\ref{eq:mds-char}). The 1995 draft used Gelbart's labels \cite{Ge2},
and its version of Lemma~\ref{lem:realds} differs from the one above by
$\varepsilon\mapsto-\varepsilon$ together with an interchange of $+$odd and
$-$odd. That is the effect of reversing the orientation
$\theta\mapsto-\theta$, so the content is the same. The labels of
\cite{Ge2} and \cite{Ad2} have not been checked against
(\ref{eq:mds-char}).
\end{remyy}

In every case the infinitesimal character of the representation of
$SL(2,{\mathbb R})$ is $|m|\in\{2\nu-1,2\nu,2\nu+1\}$, where
$\nu=k/2-1$ is the infinitesimal character of $\overline\pi{}^\pm_{k/2}$
(both normalized so that the trivial representation of $SL(2,{\mathbb R})$
has infinitesimal character $1$). Every discrete series or limit of
discrete series representation of $SL(2,{\mathbb R})$ occurs in at least
one of (a)--(c).

\begin{remark}
First, we recall from \cite{HC2}, page 660, that
\[
\Theta _{\pi _m}((-I_2)^ja(x))
=(-1)^{(m-1)j}\frac{e^{-|m\ln(x)|}}{|x-x^{-1}|}%
,\;\;m\neq 0,\;x>0,
\]
and from \cite{Ad2}, (15.7)(b), that
\[
\Theta _{\overline{\pi }_{k/2}^\varepsilon }(a(x))=\frac{%
c_{\varepsilon,k}(x)\,e^{-|(\frac k2-1)\ln|x||}}{|x-x^{-1}|},
\]
where $c_{\varepsilon,k}(x)=1$ for $x>0$ and $c_{\varepsilon,k}(x)$ is a fourth
root of unity for $x<0$, written $\sgn(x)^{-\varepsilon k/2}$ in \cite{Ad2}
(with the labels of \cite{Ad2}, which have not been compared with those
of Lemma~\ref{lem:realds}).
By definition, if $f,\overline{f}$ are corresponding on the hyperbolic set
then $(F_f)^{\wedge }(\chi )
=(F_{\overline{f}})^{\wedge }(\overline{\chi })$,
where $\chi ,\overline{\chi }$ are corresponding characters and
$F_\phi$ is the normalized orbital integral. Combining these
together, we see that without the assumption that
the hyperbolic orbital integrals vanish, the result is
inconsistent with the matching of principal series. Indeed the
exponents $|m|=k-2$ and $k/2-1$ differ by a factor of $2$: the elliptic
matchings double the infinitesimal character, and the hyperbolic matching
of \S\ref{sec:realA} preserves it (compare Remark~\ref{rem:even}).
\end{remark}

%\begin{center}
%\textbf{REFERENCES}
%\end{center}

%[A] J. Adams, Lifting of characters on orthogonal and metaplectic groups,
%preprint, 1995

%[Fl] Y. Flicker, Automorphic forms on covering groups of GL(2), Inv. Math.,
%1980

%[G] S. Gelbart, \underline{Weil's representation and the spectrum of the
%metaplectic group}, SLNM 530, 1976

%[HC] Harish-Chandra, \underline{Collected Works, vol III}, Springer-Verlag

%[Hum] J. Humpheries, \underline{Linear algebraic groups}, Springer-Verlag Grad Texts

%[Mum] D. Mumford, \underline{The red book}, Springer LNM, 1358, 1988

%[V] V. S. Varadarajan, An introduction to harmonic analysis on semi-simple
%Lie groups, Cambridge U. Press, 1989
%% ===== end aabook2c.tex =====


%% ===== begin aabook2d.tex (revised) =====
\section{On the correspondence of d. s. for
\texorpdfstring{$\overline{SL(2,F)}$}{SL(2,F)-bar}, \texorpdfstring{$p>2$}{p>2}}\label{sec:padic-ds}

We have described several correspondences
$c_i:\overline{\pi}\longmapsto \pi$
between representations of
$\overline{SL(2,F)}$ and those of $SL(2,F)$. They are meant to be
consistent with a ``matching''
$c_i^*:\overline{f}\longmapsto f$
of functions, in the sense that the characters match:
$\Theta_{\overline{\pi}}(\overline{f})=\Theta_{\pi}(f)$, where
$\pi=c_i(\overline{\pi})$ and $f=c_i^*(\overline{f})$.
In the 1995 draft it was asserted that every element $\pi$ of the
unitary dual of $G(F)$ occurs in the image of at
least one of the $c_i$. In view of Remarks~\ref{rem:sc-caveat}
and~\ref{rem:savin-caveat}, this is not established here.

Let $\tau:G(F)\rightarrow \overline{G(F)}$ denote the map
$g\longmapsto \tilde{g}$ of \cite{Fl}.

\subsection{The Iwahori algebra}

This subsection records how the 1995 draft used Savin's results on Iwahori--Hecke algebras, with the caveats they need.


Let ${\cal I}$ denote the Iwahori subgroup of $G(F)$ and assume $p$ is odd.
Savin \cite{Sav} studied the genuine Iwahori--Hecke algebras of
metaplectic covers of split groups. The 1995 draft used his work in the
following form.

\begin{remyy}[1995 draft, attributed to Savin \cite{Sav}; unverified, see
Remark~\ref{rem:savin-caveat}]\label{claim:savin} There is an isomorphism
of algebras
${\cal H}(\overline{G}//{\cal I})\rightarrow {\cal H}(G//{\cal I})$.
\end{remyy}

\begin{remyy}[1995 draft; depends on Remark~\ref{claim:savin}]
The isomorphism of Remark~\ref{claim:savin}
induces a set-theoretic bijection
\[
{\cal H}(\overline{G}//{\cal I})^\wedge\rightarrow {\cal H}(G//{\cal I})^\wedge
\]
on the sets of equivalence
classes of irreducible (genuine, in the case of $\overline{G}$)
representations with an ${\cal I}$-fixed vector.
Furthermore, if
\[
\begin{array}{ccc}
{\cal H}(G//{\cal I})&\rightarrow &{\cal H}(\overline{G}//{\cal I})\\
f & \longmapsto & \overline{f}
\end{array}
\]
and
\[
\begin{array}{ccc}
{\cal H}(\overline{G}//{\cal I})^\wedge &\rightarrow &{\cal H}(G//{\cal I})^\wedge\\
\overline{\pi}&\longmapsto &\pi
\end{array}
\]
then $\tr(\overline{\pi}(\overline{f}))=\tr(\pi(f))$.
Furthermore, if $\overline{\pi}$ is a discrete series representation
then so is $\pi$.
\end{remyy}

\begin{remark}\label{rem:savin-caveat}
In the later literature on Hecke algebra correspondences
for metaplectic groups (see \cite{TW} and the references there), the
linear group paired with the metaplectic cover of $SL(2)=Sp(2)$ is
$SO(3)\cong PGL(2)$, not $SL(2)$, and the resulting bijection of
Iwahori-spherical representations is with representations of $PGL(2,F)$.
The two statements above, with $G=SL(2)$, should therefore be
checked against \cite{Sav} before being used. In particular, the 1995
draft's assertion that this isomorphism is associated with an identity
of orbital integrals rested on the argument of the 1995 version of
\S\ref{sec:dyadic}, which has been withdrawn.
Neither statement is used in Chapter~\ref{ch:global}.
\end{remark}

The 1995 draft intended this corollary to be used in the comparison of
the special representations at non-archimedean places $v$ not dividing
$2$.

%% ===== end aabook2d.tex =====


%% ===== begin aabook2e.tex =====
\begin{example}
The Steinberg representation of $SL(2,F)$ has a nonzero vector fixed by
the Iwahori subgroup ${\cal I}$, but none fixed by $K_0$ or $K_1$. It is therefore
not covered by Proposition~\ref{prop:LH-spherical}. A correspondence of
Iwahori--Hecke algebras, as in Remark~\ref{claim:savin}, is what would be
needed at such a place.
\end{example}

\section{On the character of the Weil representation}\label{sec:weil}

We formally compute the character of the ``even'' and ``odd'' parts
of the Weil representation on the ``big cell'' of the Bruhat decomposition
by integrating the kernel along the diagonal. This diagonal integral is
not absolutely convergent, and identifying its regularized value with
the distribution character requires an argument beyond the formal
calculation. We then discuss possible consequences for local matching
of representations. Only formula (iii) of \S\ref{sec:weil-rep} (in
Lemma~\ref{lem:A00}) and the lemmas on non-tempered representations (in
Remark~\ref{rem:LH-families}) are used later. The kernel of the Weil
representation and the character of $r_+-r_-$ on the regular set are
computed, for Kubota's cocycle and with a proof that the diagonal
integral is the character, in the companion text \cite{JC}; see item~6
of \S\ref{sec:cover-supplies}.

Here $F$ may be either p-adic or real. If $F$ is p-adic we shall
sometimes impose the assumption that $p\not= 2$.

\subsection{A Bruhat decomposition}

We write elements of the big cell of the cover as products of simple elements, keeping track of the cocycle.


Let $G=SL(2)$, let
\[
h(x)=\left(
\begin{array}{cc}
x & 0 \\
0 & x^{-1}
\end{array}
\right)
,\ \ \ \
n(x)=\left(
\begin{array}{cc}
1 & x \\
0 & 1
\end{array}
\right)
\]
and let
\[
w:=\left(
\begin{array}{cc}
0 & -1 \\
1 & 0
\end{array}
\right),\ \ \
r(\theta )=\left(
\begin{array}{cc}
\cos(\theta) & -\sin(\theta) \\
\sin(\theta) & \cos(\theta)
\end{array}
\right) .
\]

Let $F$ be either ${\mathbb{R}}$ or a p-adic field.
If $F$ is p-adic, let $|\ |$ denote the norm which is normalized
so that $|\varpi|=q^{-1}$, where $q$ denotes the number of elements of the
residue field, and let ${\cal O}_F$ denote the ring of integers.
Let $(\ ,\ )$ denote
the quadratic Hilbert symbol and let $\psi$ denote a non-trivial
additive character of $F$. If $F$ is real, take $\psi(x)=\exp(\pi ix)$
and if $F$ is p-adic then assume $\psi$ has conductor ${\cal O}_F$.

Let $\beta: G(F)\times G(F)\rightarrow \mu_2=\{\pm 1\}$
denote the cocycle
\[
\beta(g,g')=({\frac{x(gg')}{x(g)}}, {\frac{x(gg')}{x(g')}} ),
\]
where $g,g'\in G(F)$ and
\[
x\left(
\begin{array}{cc}
a & b \\
c & d
\end{array}
\right)
=
\left\{
    \begin{array}{c}
    c,\ \ \ c\not= 0,\\
    d,\ \ \ {\rm otherwise}.
    \end{array}
\right.
\]
We define a group structure $\overline{G(F)}$ on
$G(F)\times \mu_2$ by
\[
(g,\zeta)\cdot (g',\zeta')=(gg',\beta(g,g')\zeta\zeta'),
\]
$g,g'\in G(F),\ \zeta,\zeta'\in \mu_2$.
The following simple fact has appeared in several places:

\begin{lemma} If
\[
g=\left(
\begin{array}{cc}
a & b \\
c & d
\end{array}
\right)
\]
has determinant 1 then
\[
(g,(-c,-1))=(n(a/c),1)\cdot (w,1)\cdot (h(-c),1)\cdot (n(d/c),1),
\]
provided $c\not= 0$.
\end{lemma}

\begin{proof} We have
\[
n(a/c)\cdot w =
\left(
\begin{array}{cc}
-a/c & 1 \\
-1 & 0
\end{array}
\right)
\]
so $\beta(n(a/c),w)
=({\frac{x(n(a/c)w)}{x(n(a/c))}},{\frac{x(n(a/c)w)}{x(w)}})
=(-1,1)=1$.
Also,
\[
n(a/c)\cdot w h(-c) =
\left(
\begin{array}{cc}
a & -1/c \\
c & 0
\end{array}
\right)
\]
so $\beta(n(a/c)w,h(-c))
=({\frac{x(n(a/c)wh(-c))}{x(n(a/c)w)}},{\frac{x(n(a/c)wh(-c))}{x(h(-c))}})
=(-c,-c^2)=(-c,-1)$.
Finally,
\[
n(a/c)\cdot w h(-c)n(d/c) =
\left(
\begin{array}{cc}
a & b \\
c & d
\end{array}
\right)=g
\]
so $\beta(n(a/c)wh(-c),n(d/c))
=({\frac{x(g)}{x(n(a/c)wh(-c))}},{\frac{x(g)}{x(n(d/c))}})
=(1,c)=1$.
\end{proof}

\subsection{The Weil representation}\label{sec:weil-rep}

We recall the Weil representation of the cover and its even and odd parts.


Let $\omega$, $q(x)$, and $\gamma(q)$ be as in \cite{Fl}, \S 2.2. See also
\cite{Ge2}, page 36, and \cite{Rao}, Appendix.

For $\phi \in L^1(F)$, let
\[
\hat{\phi}_q(x)=\int_F \phi(y)\psi(2qxy)dy,
\]
where $dy$ is normalized as in \cite{Ge2}, ch 2, so that this extends to
an isometry of $L^2(F)$. This implies that $dy$ depends on $\psi, q$
in a way we will make explicit later.
(In the case where $F={\mathbb{R}}$, if $\psi(x)=\exp(\pi i x)$ then
$dy$ is $\sqrt{|q|}$ times Lebesgue measure.)

Let $\epsilon \in \{\pm\}$ and let
\[
L^2(F)_\epsilon=
\{\phi \ |\ \phi(-x)=\epsilon \phi(x),\ \ \forall x\in F\}.
\]
For $\phi\in L^2(F)_\epsilon$, we have
\[
\int_F \phi(x)\psi(2qxy)dx =
\int_{F/\{\pm \}}
\phi(x) (\psi(2qxy)+\epsilon \psi(-2qxy))\overline{dx},
\]
if $dx=2\overline{dx}$.

We define the Weil representation $r=r_{\psi,q}$ of
$\overline{G(F)}$ on $L^2(F)$ as in \cite{Fl}, \S 2.2, so

(i) $r(n(b),\zeta)\phi(x)=\zeta\cdot \psi(bq(x))\phi(x)$,

(ii) $r(w,\zeta)\phi(x)
=\zeta\cdot \gamma(q)\hat{\phi}_q(x)$,

(iii) $r(h(a),\zeta)\phi(x)=\zeta\cdot |a|^{1/2}\omega(a)\phi(ax)$.

We define the Weil representation $r_\epsilon=r_{\psi,q,\epsilon}$ of
$\overline{G(F)}$ on $L^2(F)_\epsilon$ by

(i) $r_\epsilon(n(b),\zeta)\phi(x)=\zeta\cdot \psi(bq(x))\phi(x)$,

(ii) $r_\epsilon(w,\zeta)\phi(x)
=\zeta\cdot {\frac{1}{2}}\gamma(q)(\hat{\phi}_q(x)
+\epsilon \hat{\phi}_{-q}(x))$,

(iii) $r_\epsilon(h(a),\zeta)\phi(x)=\zeta\cdot |a|^{1/2}\omega(a)\phi(ax)$.

The displayed formulas agree with the restriction of $r$ to
$L^2(F)_\epsilon$. Indeed, the parity relation
$\phi(-y)=\epsilon\phi(y)$ gives
$\widehat\phi_{-q}(x)=\epsilon\widehat\phi_q(x)$; hence the average
in (ii) equals $\widehat\phi_q(x)$.
The even and odd subspaces are therefore $r$-invariant.

We have
\[
r_\epsilon(g,1)=(-c,-1)
r_\epsilon(n(a/c),1)\cdot r(w,1)\cdot r_\epsilon(h(-c),1)
\cdot r_\epsilon(n(d/c),1),
\]
where $g$ is as in the lemma above. This implies, for $\phi\in L^2(F)_\epsilon$,
\[
r_\epsilon(g,1)\phi(x)=(-c,-1)|c|^{-1/2}\omega(-c)\gamma(q)\,
\psi\bigl({\textstyle\frac{a}{c}}q(x)\bigr)
\int_F\phi(y)\,\psi\bigl({\textstyle\frac{d}{c}}q(y)\bigr)\,
\psi_q\bigl(-{\textstyle\frac{2}{c}}xy\bigr)\,dy ,
\]
where $\psi_q(x)=\psi(qx)$ and $q(x)=qx^2$.

Since $r_+$ is even and $r_-$ is odd, we point out that
the characters $\Theta_\epsilon$ of each of these
representations may be obtained from the
character of the sum (i.e., the character of $r$)
by virtue of the trick
\[
\Theta_+(-g,1)+\Theta_-(-g,1)=
\Theta_+(g,1)-\Theta_-(g,1)
\]
(up to the scalar $\omega(-1)$ by which $(-I,1)$ acts on $L^2(F)_+$, and the
cocycle $\beta(-I,g)$). The character $\Theta_+(g)+\Theta_-(g)$ is determined in \cite{H}.
We shall take another approach.

\subsection{Integrating the kernel}

We study the \emph{formal} diagonal integral of the kernel. The
character identity below remains a conjectural regularization statement,
not a proved trace formula for the Weil representation.


For $\phi\in L^2(F)_\epsilon$, we have
\[
r_\epsilon(g,1)\phi(x)=
\int_{F}
k_\epsilon (x,y)\phi(y)dy,
\]
where the kernel function (which also depends on $g$)
\[
\begin{array}{c}
k_\epsilon (x,y)=
(-c,-1)|{\frac{1}{c}}|^{1/2}
\omega(-c)\gamma(q)\,\psi\bigl({\frac{a\,q(x)+d\,q(y)}{c}}\bigr)\,
{\frac{\psi_{-q}({\frac{2}{c}}xy)+\epsilon \psi_q({\frac{2}{c}}xy)}{2}},\\[1ex]
k_\epsilon (x,x)=(-c,-1)|{\frac{1}{c}}|^{1/2}
\omega(-c)\gamma(q)\,
{\frac{\psi_{-q}(-{\frac{a+d-2}{c}}x^2)+\epsilon \psi_q({\frac{a+d+2}{c}}x^2)}{2}} .
\end{array}
\]
Note that $k_\epsilon(x,x)$ may be regarded as a function on either
$F$ or $F/\{\pm 1\}$.
Also, note if $\phi\in L^2(F)_\epsilon$ then
the functions $x\longmapsto k_\epsilon(x,y)\phi(x)$,
$y\longmapsto k_\epsilon(x,y)\phi(y)$ are
functions on $F/\{\pm 1\}$.

The restriction $r_\epsilon$ of $r$ to $L^2(F)_\epsilon$
is an irreducible admissible representation (see, for example,
\cite{Fl}, \S 2.2, or \cite{Ge2}). For each genuine $f\in C_c^\infty(\overline{G(F)})$,
the trace of the operator $r_\epsilon(f)$ exists and there is an
genuine character function $\Theta_\epsilon(g)$ on $\overline{G(F)}$
such that
\[
\tr(r_\epsilon(f))=\int_{\overline{G(F)}} \Theta_\epsilon (g)f(g)dg.
\]

We shall assume the following

\begin{conjecture}
On the regular elements of the ``big cell''
$N(F)wN(F)\overline{A(F)}$, the character $\Theta_\epsilon$
is given by the principal
value integral of the kernel function along the diagonal:
\[
\Theta_\epsilon(g,1)=
\lim_{N\rightarrow \infty}
\int_{\{x\in F\ |\ |x|<N\}}
k_\epsilon(x,x)dx.
\]
Here $dx$ is the normalized Haar measure as above.
\end{conjecture}

\begin{remark} This conjecture implies
\[
\begin{array}{c}
\Theta_\epsilon(g,1)\\
=(-c,-1)|{\frac{1}{c}}|^{1/2}\omega(-c)\gamma(q)
\lim_{N\rightarrow \infty}
\int_{\{x\in F\ |\ |x|<N\}}
{\frac{\psi_{-q}(-{\frac{a+d-2}{c}}x^2)
+\epsilon \psi_q({\frac{a+d+2}{c}}x^2)}{2}}dx.
\end{array}
\]
\end{remark}

Let
\[
\int_{\mathbb{R}} \exp(aix^2)dx
=\lim_{N\rightarrow \infty}
\int_{\{x\in {\mathbb{R}}\ |\ |x|<N\}} \exp(aix^2)dx.
\]
The following result is elementary and well-known:

\begin{lemma}\label{lem:gauss} Let $d_Lx$ denote the usual Lebesgue measure.
If $F=\mathbb{R}$ and $a\not= 0$ is real then
\[
\int_{F} \exp(aix^2)d_Lx =\exp(\pi i \sgn(a)/4)\sqrt{\frac{\pi}{|a|}}.
\]
\end{lemma}

\begin{remark}
If $\langle\cdot,\cdot\rangle$ is the usual inner product on
${\mathbb R}^n$ then
\[
\int_{{\mathbb R}^n} \exp(ai\langle\vec{x},\vec{x}\rangle)d_L x =e^{\pi i n\,\sgn(a)/4}({\frac{\pi}{|a|}})^{n/2}.
\]
Using simple generalizations like this it is hoped that it will
be possible to generalize some of the results below to
$\overline{Sp(2n,{\mathbb{R}})}$.
\end{remark}

\begin{proof} We may assume without loss of generality that $a>0$.
Let $C_R=C_{R,a}\cup C_{R,b}\cup C_{R,c}$
denote the sector of the circle of radius $R$,
consisting of $z=re^{i\theta}$ with (a) $\theta=0,\ 0\leq r\leq R$,
(b) $\theta=\pi/4,\ 0\leq r\leq R$, and (c) $0\leq \theta \leq \pi/4,\ r=R$.
We have
\[
\begin{array}{c}
\int_{C_{R,c}} e^{iaz^2}dz=
Ri\int_0^{\pi/4} e^{R^2(-\sin 2\theta +i\cos 2\theta )}e^{i\theta} d\theta\\
\ll  R\int_0^{\pi/4} e^{-R^2\sin 2\theta }d\theta\\
\ll  R\int_0^{\epsilon} e^{-2R^2x}dx +Re^{-\epsilon R^2}
\ll \ln(R)/R,
\end{array}
\]
picking $\epsilon =2\ln(R)/R^2$.
This allows us to use a standard residue theorem argument
applied to $C_R$ to obtain
\[
\int_F e^{iax^2}dx=\exp(\pi i/4)\int_F e^{-ay^2}dy.
\]
The result now follows from the fact that, for $a>0$,
\[
\int_0^\infty e^{-ax^2}dx={\frac{1}{2}}\sqrt{\frac{\pi}{a}}.
\]
\end{proof}

\begin{example}
For $a=1$, Lemma~\ref{lem:gauss} is the Fresnel integral
$\int_{\mathbb R}e^{ix^2}\,dx=e^{\pi i/4}\sqrt{\pi}$.
\end{example}

Let $\tau (x)=\exp (\pi i q x)$, so $\tau=\psi_q$,
and let $d_\tau x$ denote
the Haar measure normalized as in \cite{Ge2}, ch 2. If
\[
\gamma(\tau)=
\int_{\mathbb{R}} \tau(x^2)d_\tau x
\]
then $\gamma(\tau) =\gamma(q)$ (in the notation of \cite{Fl}, \S 2.2) and
\[
\gamma(\tau)=\exp ({\frac{\pi i}{4}} \sgn(q)),
\]
by \cite{Ge2}, page 36. Let $\tau_a(x)=\tau(ax)$. Then
\[
\int_{\mathbb{R}} \tau_a(x^2)d_\tau x
=|a|^{-1/2}\exp\left({\frac{\pi i}{4}}\sgn(qa)\right).
\tag{*}
\]
For nonzero real $u$, set
$\gamma_{\mathbb R}(u)=\exp(\pi i\sgn(u)/4)$, and put
$\alpha=q(a+d-2)/c$, $\beta=q(a+d+2)/c$.
Then, when both parameters are nonzero,
\[
\begin{array}{c}
\gamma(q)
\int_{\mathbb{R}}
\left(\psi_{-q}\left(-\frac{a+d-2}{c}x^2\right)
+\epsilon\psi_q\left(\frac{a+d+2}{c}x^2\right)\right)d_\tau x\\
=\gamma(q)\left[
\frac{\gamma_{\mathbb R}(\alpha)}
     {\left|\frac{a+d-2}{c}\right|^{1/2}}
+\epsilon
\frac{\gamma_{\mathbb R}(\beta)}
     {\left|\frac{a+d+2}{c}\right|^{1/2}}
\right].
\end{array}
\]
The real phases depend on the signs of $\alpha$ and $\beta$; they must
not be replaced by constants $\gamma(-q)$ and $\gamma(q)$ independently
of the signs of $a+d\pm2$ and $c$.

In the p-adic case, let
\[
\kappa(\tau)=\lim_{N\rightarrow - \infty}
\int_{\varpi^N{\cal O}_F}
\tau(u^2) du.
\]
For a \emph{fixed} Haar measure $du$, we have
$\kappa(\tau_{a^2})=|a|^{-1}\kappa(\tau)$, for all
$a\in F^\times$. Now normalize the measure $du$ as in \cite{Mo3}
(so $du=d_\tau u$ depends on $\tau$) and
call the result $\kappa_{nor}(\tau)$. It is known
that (see for example \cite{Mo3})
\[
\kappa_{nor}(\tau)=
\left\{
    \begin{array}{c}
    1,\ \ \ 2|f(\tau),\\
    G(\tau),\ \ \ {\rm otherwise}.
    \end{array}
\right.
\]
where the conductor of $\tau$ is given by
$\varpi^{f(\tau)}{\cal O}_F$ and where $G(\tau)$ denotes the Gauss sum
normalized so as to be an eighth root of unity as in \cite{Mo3}.
Now, the Haar measure $d_\tau u$ has been chosen to depend on
$\tau$ in such a way that $\operatorname{vol}({\cal O}_F,d_\tau u)=q^{f(\tau)/2}$.
This is the measure for which $(\hat{\phi}_q)^\wedge_q(x)=\phi(-x)$,
by \cite{Shal}.
Note that $f(\tau_a)=f(\tau)-\nu(a)$, where
$\nu : F^\times \rightarrow {\mathbb{Z}}$ is the normalized
valuation. Putting these together gives
$\kappa_{nor}(\tau_{a^2})=\kappa_{nor}(\tau)$. It is also known that
$\kappa_{nor}(\tau_{a})=({\frac{a}{\varpi}})\kappa_{nor}(\tau)$, for
$a\in {\cal O}_F^\times$, where $\left(\frac{\cdot}{\varpi}\right)$ is the
Legendre symbol. Of course,
\[
\operatorname{vol}({\cal O}_F,d_\tau u)^{-1}
\kappa_{nor}(\tau)
=\operatorname{vol}({\cal O}_F,du)^{-1}\kappa(\tau),
\]
so it is easy to convert from one to the other.
We assume that $du$ is the measure in the conjecture above.

\begin{proposition}[Conditional diagonal-kernel identity]
\label{prop:kernel-conditional}
Assume the principal-value conjecture above, and assume that
$c\ne0$, $a+d\ne\pm2$. With the Haar measure $d_{\psi_q}x$
used for the Fourier transform in \S\ref{sec:weil-rep}, set
\[
\mathcal K_{\psi_q}(\alpha)
=\operatorname{PV}\int_F\psi(\alpha x^2)\,d_{\psi_q}x,\qquad
\alpha=\frac{q(a+d-2)}{c},\qquad
\beta=\frac{q(a+d+2)}{c}.
\]
Then the formal kernel calculation gives
\[
\Theta_\epsilon(g,1)=
\frac{(-c,-1)|c|^{-1/2}\omega(-c)\gamma(q)}{2}
\left(\mathcal K_{\psi_q}(\alpha)
+\epsilon\mathcal K_{\psi_q}(\beta)\right).
\]
Here $\gamma(q)$ is the prefactor in the Weil representation
of \S\ref{sec:weil-rep}; it has not been redefined through the
Gaussian integral.
\end{proposition}

\begin{proof}
The displayed formula is the diagonal integral of
$k_\epsilon(x,x)$ from \S\ref{sec:weil-rep}, with the two
phases written as $\psi(\alpha x^2)$ and $\psi(\beta x^2)$.
No change of Haar measure is made in passing to
$\mathcal K_{\psi_q}$. The conjecture is needed to identify
this regularized expression with the distribution character.
\end{proof}

\begin{remyy}[Real specialization, conditional]
\label{rem26:real-kernel}
Take $F=\mathbb R$ and $g=r(\theta)$ regular elliptic.
Here
\[
\alpha=\frac{2q(\cos\theta-1)}{\sin\theta},\qquad
\beta =\frac{2q(\cos\theta+1)}{\sin\theta}.
\]
The Gaussian calculation above gives the \emph{conditional} expression
\[
\Theta_\epsilon(g,\zeta_g)=
\frac{\zeta_g(-\sin\theta,-1)\omega(-\sin\theta)\gamma(q)}{2}
\left(
\frac{\gamma_{\mathbb R}(\alpha)}
     {|\det(g-I_2)|^{1/2}}
+\epsilon
\frac{\gamma_{\mathbb R}(\beta)}
     {|\det(g+I_2)|^{1/2}}
\right).
\]
The phases $\gamma_{\mathbb R}(\alpha)$ and
$\gamma_{\mathbb R}(\beta)$ vary across the arcs
$(-\pi,0)$ and $(0,\pi)$, and across the other intervals in
the covering parameter. The unsigned formulas in the 1995
version suppressed this dependence. In particular, an
agreement on one interval does not extend to the others
merely by conjugation with the rotation $w$, since $w$
commutes with every $r(\theta)$.
A full identification with the character formulas in
\cite{Ad1} requires checking the lift $\zeta_g$ and the
Weil-index conventions. This verification has not been supplied.
\end{remyy}

\begin{remyy}[Non-archimedean specialization, conditional]
\label{rem26:padic-kernel}
If $F$ is p-adic and
$g=\left(\begin{smallmatrix}a&\eta b\\b&a\end{smallmatrix}\right)$
is regular elliptic, then $c=b$, and in
Proposition~\ref{prop:kernel-conditional} one has
\[
\alpha=\frac{2q(a-1)}{b},\qquad
\beta=\frac{2q(a+1)}{b}.
\]
The proposition must be used with the \emph{same} Haar measure
$d_{\psi_q}x$ in the two quadratic integrals.
The earlier 1995 expression with a prefactor
$|q/b|^{1/2}$ and coefficients
$\gamma(q)\gamma(-q)$ and $\gamma(q)^2$ did not
follow from the diagonal kernel: it mixed the kernel prefactor
with subsequent Gaussian normalizations.
No verified global-in-$g$ character formula is claimed here;
compare \cite{H}.
\end{remyy}

\begin{remark}
The preceding two remarks retain the intended historical comparison
with \cite{Ad1} and \cite{H}, but they are not used in
Chapter~\ref{ch:global}. In particular, neither constitutes a
proof of the principal-value conjecture or an assertion about
the complete character distribution of the Weil representation.
\end{remark}

\subsection{Matchings for non-tempered representations}

We deduce matchings for the non-tempered genuine representations from the results on discrete series.


Let $r_{even}$ denote the even part of the Weil representation
on $\overline{GL(2,F)}$ in \cite{Fl}, \S 2.2 (where it is denoted by
$r_\mu$, $\mu$ an even character which we take here to be the
trivial character).

\subsubsection{The real case}

Over ${\mathbb R}$, the even halves of the two oscillator representations are the non-tempered genuine representations.


Let $F={\mathbb{R}}$. Let $\overline{\pi}_{+,nt},\overline{\pi}_{-,nt}$
denote the two irreducible, genuine, non-tempered, unitary representations
of $\overline{G(F)}$. In the notation of \cite{Ad1}, \cite{Ge2}, we have
\[
\overline{\pi}_{+,nt}
=\overline{\pi}^+_{1/2}=\omega(\psi_+)_{even},
\ \ \ \ \
\overline{\pi}_{-,nt}
=\overline{\pi}^-_{1/2}=\omega(\psi_-)_{even}.
\]
On the Grothendieck group of the unitary dual of $\overline{G(F)}$,
we have
\[
r_{even}|_{\overline{G(F)}}=
\overline{\pi}_{+,nt}+\overline{\pi}_{-,nt}.
\]
Let $\pi_{nt}$ denote the (non-tempered) trivial representation of $G(F)$.

Using the notation above and that
of \cite{Ad1}-\cite{Ad2}, we have the following additional
identities in the Grothendieck group
of the unitary dual of $\overline{G(F)}$ or $G(F)$
(whichever is appropriate to the context):
\begin{equation}\label{eq:groth}
\begin{array}{c}
\pi'_{PS}(1/2,1/2)=\overline{\pi}_{+,nt}+\pi'_{DS}(-,1/2),\\
\pi'_{PS}(-1/2,1/2)=\overline{\pi}_{-,nt}+\pi'_{DS}(+,1/2),\\
\pi_{PS}(0,1)=\pi_{nt}+\pi_{DS}(+,1)+\pi_{DS}(-,1),\\
\pi_{PS}(0,-1)=\pi_{nt}+\pi_{DS}(+,1)+\pi_{DS}(-,1).
\end{array}
\end{equation}
(The last two lines refer to the two principal series that contain the
trivial representation $\pi_{nt}$; they have the same composition
factors.)
Here $\pi'_{DS}(\pm,1/2)$ are the genuine discrete series with
$\overline K$-types $\pm(3/2+2j)$, $j\ge0$, i.e.\
$\overline\pi{}^{\pm}_{3/2}$ in the notation of \S\ref{sec:realds}, and
$\pi_{DS}(\pm,1)=\pi_{\pm1}$. (For instance, the even oscillator
representation $\overline\pi_{+,nt}$ has $\overline K$-types
$1/2+2j$, $j\geq0$, and the remaining $\overline K$-types $-3/2-2j$ of the
principal series containing it are those of $\overline\pi{}^-_{3/2}$.)

\begin{lemma}\label{lem:ds-3/2} If
$f,\overline{f}$ correspond evenly on the elliptic set in
the sense of Definition~\ref{def:realds-corr}, then
$\Theta_{\pi'_{DS}(\pm,1/2)}(\overline{f})
=\Theta_{\pi_{DS}(\pm,1)}(f)$.
\end{lemma}

\begin{proof}
This is Lemma~\ref{lem:realds}(a) with $k=3$.
\end{proof}

\begin{lemma} If $f,\overline{f}$ correspond evenly on the elliptic set
in the sense of Definition~\ref{def:realds-corr}, then
\[
\Theta_{\overline{\pi}_{+,nt}}(\overline{f})
+\Theta_{\overline{\pi}_{-,nt}}(\overline{f})
=\Theta_{\pi_{nt}}(f).
\]
\end{lemma}

\begin{proof}
The characters of principal series representations
vanish on the regular elliptic set, and the hyperbolic orbital
integrals of $f$ and $\overline f$ vanish. Hence, by (\ref{eq:groth}) and
the Weyl integration formula,
$\Theta_{\overline{\pi}_{\pm,nt}}(\overline f)=-\Theta_{\pi'_{DS}(\mp,1/2)}(\overline f)$
and
$\Theta_{\pi_{nt}}(f)=-\Theta_{\pi_{DS}(+,1)}(f)-\Theta_{\pi_{DS}(-,1)}(f)$.
The result follows from Lemma~\ref{lem:ds-3/2}. (The identities
(\ref{eq:groth}) are quoted from \cite{Ad1}, \cite{Ge2}, and have not
been rechecked here.)
\end{proof}

\begin{remark}
Indeed, it seems likely that \cite{Fl} implies that
if $f,\overline{f}$ correspond evenly on the hyperbolic set in
the sense of \S\ref{sec:charsA}, then
\[
\Theta_{\overline{\pi}_{+,nt}}(\overline{f})
+\Theta_{\overline{\pi}_{-,nt}}(\overline{f})
=\Theta_{\pi_{nt}}(f).
\]
\end{remark}

\subsubsection{The \texorpdfstring{$p$}{p}-adic case}

Over a p-adic field we record the 1995 statements about the Weil representation. They have not been verified.


Let $F$ be p-adic. Let $r_{odd}$ denote the odd part of some Weil
representation (associated to some odd character)
on $\overline{GL(2,F)}$ as in \cite{Fl}, \S 2.2. This is an
irreducible genuine supercuspidal representation of
$\overline{GL(2,F)}$. Its restriction to
$\overline{SL(2,F)}$ decomposes into $4=[F^\times:F^{\times 2}]$
irreducible genuine supercuspidal components $\overline{\pi}_{sc,i}$
\cite{M3}.

Let $r_{even}$ denote the even part of the Weil
representation associated to the trivial character on
$\overline{GL(2,F)}$ as in \cite{Fl}, \S 2.2. This is an
irreducible genuine representation of
$\overline{GL(2,F)}$, which is not supercuspidal (its Jacquet module is
nonzero). Its restriction to
$\overline{SL(2,F)}$ decomposes into $4=[F^\times:F^{\times 2}]$
irreducible genuine infinite dimensional
non-tempered components $\overline{\pi}_{nt,i}$.

Let $\sigma$ denote the special representation of $SL(2,F)$
and $\pi_{nt}$ the trivial representation.

The 1995 draft derived the following from \cite{Fl} and the fact that the
restriction of the special (Steinberg) representation of $GL(2,F)$
to $SL(2,F)$ is irreducible.

\begin{remyy}[1995 draft; unverified] (a) If $f, \overline{f}$ correspond evenly on the elliptic set
then
\[
\sum_i \Theta_{\overline{\pi}_{sc,i}}(\overline{f})
=\Theta_\sigma(f).
\]
(b) If $f,\overline{f}$ correspond evenly on the elliptic set
then
\[
\sum_i \Theta_{\overline{\pi}_{nt,i}}(\overline{f})
=\Theta_{\pi_{nt}}(f).
\]
(c) Let $\overline{\sigma}_i$ denote the complement of
$\overline{\pi}_{nt,i}$ (in the associated full unitarily
induced representation). If $f,\overline{f}$ correspond
evenly on the elliptic set then
\[
\sum_i \Theta_{\overline{\sigma}_{i}}(\overline{f})
=-\Theta_{\pi_{nt}}(f).
\]
\end{remyy}

\begin{remark} If $f,\overline{f}$ correspond oddly on the elliptic set
then we have the supercuspidal
correspondence described in Appendix~\ref{ch:sc} (see Remark~\ref{rem:sc-caveat}).
\end{remark}

%\begin{center}
%\textbf{REFERENCES}
%\end{center}

%[A] J. Adams, ``Character of the oscillator representations'',
%preprint, 1995 (to appear in Isr. J. Math.)

%[A'] ---, Notes on the Flicker correspondence

%[A''] ---, ``Lifting of characters on orthogonal and metaplectic groups'',
%preprint, 1995

%[Fl] Y. Flicker, ``Automorphic forms on covering groups of GL(2)'',
%Inv. Math., 1980

%[G] S. Gelbart, \underline{Weil's representation and the spectrum of the
%metaplectic group}, SLNM 530, 1976

%[H] R. Howe, unpublished

%[HC] Harish-Chandra, \underline{Collected Works, vol III}, Springer-Verlag

%[Hum] J. Humpheries, \underline{Linear algebraic groups},
%Springer-Verlag Grad Texts

%[Mum] D. Mumford, \underline{The red book}, Springer LNM, 1358, 1988

%[J] D. Joyner, book

%[J'] ---, Notes on the ds correspondence for $\overline{SL(2,F)}$,
%$F$ p-adic

%[J''] ---, Notes on the ds correspondence for
%$\overline{SL(2,{\mathbb{R}})}$

%[Man] D. Manderscheid, ``Supercuspidal representations and the
%theta correspondence'', preprint

%[M] C. Moen, ``The dual pair $(U(1),U(1))$ over a p-adic field'',
%Pac J Math, vol 158, 1993

%[R] R. Rao, ``On some explicit formulas in the theory of Weil
%representations'', Pac J Math \underbar{157}(1993)335-371

%[S] J. Shalika, ``Representations of the two by two unimodular group
%over local fields'', IAS

%[V] V. S. Varadarajan, An introduction to harmonic analysis on semi-simple
%Lie groups, Cambridge U. Press, 1989
%% ===== end aabook2e.tex =====



%% ===== begin aabook5.tex =====
\chapter{Global consequences}\label{ch:global}

Throughout this chapter, $k$ is a number field, ${\mathbb{A}}={\mathbb{A}}_k$,
$G=SL(2)$, and $\overline{G({\mathbb{A}})}$ is the 2-fold cover of
\S\ref{sec:global}. Let ${\cal H}(G({\mathbb{A}}))$ and
${\cal H}(\overline{G({\mathbb{A}})})$ denote the spaces of global
functions of \S\ref{sec:globalalg}. We write $R_{cusp}$ for the
representation of $G({\mathbb{A}})$ on
$L^2_{cusp}(G(k)\backslash G({\mathbb{A}}))$, where $G(k)$ is
embedded in $\overline{G({\mathbb{A}})}$ by the splitting of
\S\ref{sec:global}. For a cuspidal automorphic representation $\pi$ of
$G({\mathbb{A}})$, $m(\pi)$ denotes its multiplicity in $R_{cusp}$.
As usual, global test functions are assumed $K_\infty$-finite.

Let ${\cal A}_0$ denote the space of genuine cusp forms on
$\overline{G({\mathbb{A}})}$, let $\Theta_0\subset{\cal A}_0$ be the closed
span of the cuspidal elementary theta series, i.e.\ the theta series
attached to the Weil representations $r_{\psi,\chi}$ of one-dimensional
quadratic forms (\cite{GePS}, \cite{Ge2}), and let ${\cal A}_{00}$ be the
orthogonal complement of $\Theta_0$ in ${\cal A}_0$. Waldspurger's
multiplicity one theorem \cite{Wa1} (see also \cite{PS}) states that every
irreducible genuine representation occurs in ${\cal A}_{00}$ with
multiplicity at most one. It is a statement about ${\cal A}_{00}$, not about
${\cal A}_0$. We write $\overline{R}_{00}$ for the representation of
$\overline{G({\mathbb{A}})}$ on the completion of ${\cal A}_{00}$, and
$m(\overline{\pi})\leq1$ for the multiplicity of a genuine irreducible
representation $\overline\pi$ in it.

For each place $v$ we write $K_{0,v}=G({\cal O}_v)$ if $v$ is finite,
and, if $v$ is finite and odd, $K_{1,v}=w_\varpi^{-1}K_{0,v}w_\varpi$ with
$w_\varpi$ as in \S\ref{sec:filtrations}. These are representatives of the
two conjugacy classes of hyperspecial maximal compact subgroups of
$G(k_v)$. Since $v$ is odd, $\overline{K_{i,v}}$ splits uniquely
(Lemma~\ref{lem:facts}(b)), and ${\cal H}(\overline{G(k_v)},K_{i,v})$
denotes the Hecke algebra of genuine functions that are bi-invariant under
the image of the splitting.

The strategy examined in this chapter is modeled on that of Gelbart and
Jacquet \cite{GJ}, who compare automorphic representations of $GL(2)$ and
of the multiplicative group of a quaternion algebra. One compares the
cuspidal spectrum of $G({\mathbb A})$ with the genuine cuspidal spectrum of
the cover by means of the two trace formulas, and uses Waldspurger's
multiplicity one theorem in place of strong multiplicity one, which fails
for both groups. \S\ref{sec:hypotheses} states the two hypotheses that this
strategy needs when the unramified places are matched by the Satake
correspondence of \S\ref{sec:spherical}. \S\ref{sec:main} proves that the
hypotheses imply a correspondence $\pi\mapsto\overline\pi$ and
multiplicity one (Proposition~\ref{thm:main}). \S\ref{sec:nonexist} proves
that no such $\overline\pi$ exists (Theorem~\ref{thm:nonexist}), so that
the hypotheses are inconsistent (Corollary~\ref{cor:inconsistent}).

\section{The hypotheses}\label{sec:hypotheses}

This section defines the exceptional set of places and the classes of
corresponding functions. It then states the two hypotheses.

For each place $v$ of $k$, let $\mathfrak{F}_v$ denote the set of
equivalence classes of irreducible unitary representations of $G(k_v)$
that occur as the local component at $v$ of a cuspidal automorphic
representation of $G({\mathbb A})$, and let $\overline{\mathfrak{F}}_v$
denote the set of equivalence classes of genuine irreducible unitary
representations of $\overline{G(k_v)}$ that occur as the local component at
$v$ of an irreducible subrepresentation of ${\cal A}_{00}$.

\begin{definition}\label{def:exceptional}
For a cuspidal automorphic representation $\pi=\otimes_v\pi_v$ of
$G({\mathbb{A}})$, let $S(\pi)$ denote the finite set of places $v$ of $k$
such that $v$ is archimedean, or $v$ is dyadic, or $v$ is odd and $\pi_v$
has no nonzero $K_{0,v}$-fixed vector.
\end{definition}

\begin{definition}\label{def:corresponding}
Let $S$ be a finite set of places of $k$ containing the archimedean and
dyadic places. For each $v\in S$, let ${\cal C}_v$ be a vector space of
pairs $(f_v,\overline{f}_v)$ with $f_v\in C_c^\infty(G(k_v))$ and
$\overline{f}_v\in C_c^\infty(\overline{G(k_v)})_{gen}$ ($K_v$-finite if
$v$ is archimedean). A pair $(f,\overline{f})$ with
$f=\otimes_v f_v\in{\cal H}(G({\mathbb{A}}))$ and
$\overline{f}=\otimes_v\overline{f}_v\in{\cal H}(\overline{G({\mathbb{A}})})$
is called {\bf $(S,{\cal C})$-corresponding} if
\begin{enumerate}
\item $(f_v,\overline{f}_v)\in{\cal C}_v$ for each $v\in S$;
\item for each $v\notin S$, $f_v\in{\cal H}(G(k_v),K_{0,v})$ and
$\overline{f}_v\in{\cal H}(\overline{G(k_v)},K_{0,v})_{gen}$ correspond
under the Satake isomorphism of \S\ref{sec:spherical};
\item for almost all $v$, $f_v$ and $\overline{f}_v$ are the unit
elements of the respective spherical Hecke algebras.
\end{enumerate}
The Satake isomorphism at $v\notin S$ depends on the choice of a
uniformizer $\varpi_v$ and of the constant $\eta_{\varpi_v}$
(Remark~\ref{rem:choices}). These choices are fixed once and for all.
\end{definition}

\begin{hypothesis}[local spectral transfer]\label{hyp:LH}
Let $v$ be a place of $k$ and $\pi_v\in\mathfrak{F}_v$. We say that
{\bf Hypothesis LH holds for
$\pi_v$ with the class ${\cal C}_v$} if ${\cal C}_v$ is a vector space of
pairs as in Definition~\ref{def:corresponding} with the following
properties.
\begin{enumerate}
\item[(LH0)] ${\cal C}_v=\bigcup_L{\cal C}_v(L)$ is a directed union of
subspaces ${\cal C}_v(L)$, indexed by open compact subgroups $L$ of
$G(k_v)$ if $v$ is finite (with the pairs in ${\cal C}_v(L)$ bi-invariant
under $L$, respectively under its image in a splitting of $\overline{L}$
over a subgroup of finite index), and by finite sets $L$ of
$K_v$-types if $v$ is archimedean (with the pairs in ${\cal C}_v(L)$
having $K_v$-types in $L$, respectively genuine $\overline{K_v}$-types in
a finite set determined by $L$).
\item[(LH1)] (Separation) For each $L$, and for each finite family
$\sigma_1,\dots,\sigma_r$ of pairwise inequivalent irreducible
admissible representations of $G(k_v)$, each of which either lies in
$\mathfrak{F}_v$ or is a transfer in the sense of (LH2), and none of whose
characters vanishes
identically on ${\cal C}_v(L)$, the linear functionals
$(f_v,\overline{f}_v)\longmapsto\tr \sigma_j(f_v)$ on
${\cal C}_v(L)$ are linearly independent.
\item[(LH2)] (Transfer) For each
$\overline{\tau}\in\overline{\mathfrak{F}}_v$, either
$\tr \overline{\tau}(\overline{f}_v)=0$ for all
$(f_v,\overline{f}_v)\in{\cal C}_v$, or there is an irreducible
admissible representation $t_v(\overline{\tau})$ of $G(k_v)$, called a
\emph{transfer} of $\overline\tau$, such that
\[
\tr \overline{\tau}(\overline{f}_v)=\tr t_v(\overline{\tau})(f_v)
\qquad\text{for all }(f_v,\overline{f}_v)\in{\cal C}_v .
\]
Moreover, $t_v$ is injective on the set of $\overline{\tau}$ for which it
is defined.
\item[(LH3)] (Detection) $\tr \pi_v(f_v)\neq 0$ for some
$(f_v,\overline{f}_v)\in{\cal C}_v$.
\item[(LH4)] (Infinitesimal characters; $v$ archimedean only) For each
infinitesimal character $\nu_v$ of $G(k_v)$, the
$\overline\tau\in\overline{\mathfrak{F}}_v$
for which $t_v(\overline\tau)$ is defined and has infinitesimal character
$\nu_v$ have only finitely many distinct infinitesimal characters.
\end{enumerate}
\end{hypothesis}

The transfer in (LH2) is unique. If $\sigma$ and $\sigma'$ are both
transfers of $\overline\tau$, choose $L$ such that $\tr\overline\tau$ does
not vanish identically on ${\cal C}_v(L)$. Then $\tr\sigma$ and
$\tr\sigma'$ are equal and nonzero on ${\cal C}_v(L)$, so
$\sigma\cong\sigma'$ by (LH1).

For example, for the elliptic matchings of \S\ref{sec:realds}, a genuine
discrete series representation of infinitesimal character $\overline\nu$
is matched with a representation of infinitesimal character
$2\overline\nu$ or $2\overline\nu\pm1$, so (LH4) holds for them.

\begin{remark}\label{rem:LH-families}
Conditions (LH1) and (LH2) are imposed only on the families
$\mathfrak{F}_v$ and $\overline{\mathfrak{F}}_v$, and not on all
irreducible representations, for the following reason. Let $v$ be real
and let ${\cal C}_v$ be one of the classes of
Definition~\ref{def:realds-corr}, which consist of pairs with vanishing
hyperbolic orbital integrals. Principal series characters vanish on the
regular elliptic set, so for every $f$ in such a class
\[
\tr\pi_0^+(f)+\tr\pi_0^-(f)=0,\qquad
\tr\pi_{nt}(f)+\tr\pi_1(f)+\tr\pi_{-1}(f)=0,
\]
and more generally $\tr\pi_m(f)+\tr\pi_{-m}(f)+\tr F_m(f)=0$ for $m\geq1$,
where $F_m$ is the irreducible representation of dimension $m$. Thus the
characters of all irreducible admissible representations are never
linearly independent on such a class. The trivial representation
$\pi_{nt}=F_1$ and the $F_m$, $m\geq2$, are not in $\mathfrak{F}_v$ (see
the end of \S\ref{sec:LH-status}). Similarly, by Lemma~\ref{lem:ds-3/2}
and the lemma following it, the even oscillator representation satisfies
$\tr\overline\pi_{+,nt}(\overline f)=-\tr\pi_{-1}(f)$ on the even class.
On the regular elliptic set the characters of the irreducible
representations of $SL(2,{\mathbb R})$, multiplied by
$e^{i\theta}-e^{-i\theta}$, are $-\sgn(m)e^{im\theta}$, $\mp1$,
$e^{im\theta}-e^{-im\theta}$, and $0$, and none of these is
$-e^{-i\theta}$. So $\overline\pi_{+,nt}$ has no transfer on the even
class. It does not lie in $\overline{\mathfrak{F}}_v$: its image under
Waldspurger's local correspondence is one-dimensional \cite{Wa2}, and
the local components of a cuspidal automorphic representation of $PGL(2)$
are infinite-dimensional.

Two limitations remain even with these restrictions. First, the relation
$\tr\pi_0^+(f)=-\tr\pi_0^-(f)$ shows that (LH1) fails for the $\pm$odd
classes of Definition~\ref{def:realds-corr} whenever both limits of
discrete series lie in $\mathfrak{F}_v$ and are detected. Second, if $f$
has a single $K$-type $n>0$ and vanishing hyperbolic orbital integrals,
then $\tr\pi_m(f)=-\tr\pi_{-m}(f)-\tr F_m(f)=0$ for every $m>0$ with
$m<n$, since neither $\pi_{-m}$ nor $F_m$ contains the $K$-type $n$. So
(LH3) for a holomorphic discrete series $\pi_m$ requires an index $L$
containing a $K$-type of $\pi_{-m}$ or of $F_m$. The same relations hold
at p-adic places for functions with vanishing hyperbolic orbital
integrals: the two constituents of a reducible unitary principal series
have opposite traces, and so do the trivial and the special
representation.
\end{remark}

\begin{definition}\label{def:infchar}
Let $\mathfrak{z}_\infty$ denote the center of the universal enveloping
algebra of the complexified Lie algebra of
$G_\infty=\prod_{v|\infty}G(k_v)$. Since $\overline{G_\infty}$ has the same
Lie algebra, $\mathfrak{z}_\infty$ acts by a character (the infinitesimal
character) on every irreducible admissible representation of $G_\infty$
or of $\overline{G_\infty}$. For a cuspidal automorphic representation
$\pi'$ of $G({\mathbb{A}})$, let $\nu_\infty(\pi')$ be the infinitesimal
character of $\pi'_\infty=\otimes_{v|\infty}\pi'_v$.

Suppose that, for each archimedean $v$, ${\cal C}_v$ satisfies (LH1) and
(LH2), with transfer $t_v$. An irreducible subrepresentation
$\overline\pi'$ of ${\cal A}_{00}$ {\bf lies over} the
character $\nu$ of $\mathfrak{z}_\infty$ if $t_v(\overline\pi'_v)$ is
defined for each archimedean $v$ and $\otimes_{v|\infty}t_v(\overline\pi'_v)$
has infinitesimal character $\nu$. For global functions $f$ and genuine
global functions $\overline f$, put
\[
\tr R_{cusp}^\nu(f)=\sum_{\pi':\ \nu_\infty(\pi')=\nu}m(\pi')\tr\pi'(f),
\qquad
\tr \overline{R}{}_{00}^{\,\nu}(\overline f)=
\sum_{\overline\pi'\text{ lying over }\nu}m(\overline\pi')\tr\overline\pi'(\overline f).
\]
These are subseries of the absolutely convergent series for
$\tr R_{cusp}(f)$ and $\tr\overline R_{00}(\overline f)$.
\end{definition}

\begin{example}
Let $k={\mathbb Q}$, and let $\pi'$ be a cuspidal automorphic
representation of $SL(2,{\mathbb A})$ occurring in the restriction of the
representation of $GL(2,{\mathbb A})$ generated by the cusp form $\Delta$
of weight $12$. Then $\pi'_\infty$ is one of the discrete series
$\pi_{\pm11}$ (lowest $K$-types $\pm12$), and $\nu_\infty(\pi')$ is the
infinitesimal character $11$. Take ${\cal C}_\infty$ to be one of the
elliptic matchings of \S\ref{sec:realds}. By Lemma~\ref{lem:realds}, a
genuine $\overline\pi'$ lies over $\nu=11$ only if
$\overline\pi'_\infty=\overline\pi{}^{\pm}_{k/2}$ with $k-2=11$ (even
matching) or with $k-1=11$ or $k-3=11$ (odd matchings). Since $k$ must
be odd, only the even matching contributes, and then
$\overline\pi'_\infty=\overline\pi{}^{\pm}_{13/2}$, of infinitesimal
character $11/2$: the finite set allowed by (LH4) is $\{11/2\}$.
\end{example}

\begin{hypothesis}[global trace identity]\label{hyp:GH}
For the given $S$ and $({\cal C}_v)_{v\in S}$, and for every character
$\nu$ of $\mathfrak{z}_\infty$,
\[
\tr R_{cusp}^\nu(f)=\tr \overline{R}{}_{00}^{\,\nu}(\overline{f})
\]
for all $(S,{\cal C})$-corresponding pairs $(f,\overline{f})$.
\end{hypothesis}

If $(f,\overline f)$ is $(S,{\cal C})$-corresponding and
$\overline\pi'\subset{\cal A}_{00}$
lies over no $\nu$, then some archimedean $\tr\overline\pi'_v$ vanishes on
${\cal C}_v$ by (LH2), so $\tr\overline\pi'(\overline f)=0$. Summing
Hypothesis~\ref{hyp:GH} over $\nu$ therefore gives
the unrefined identity $\tr R_{cusp}(f)=\tr\overline R_{00}(\overline f)$.

\begin{remark}\label{rem:why-refine}
(1) The unrefined identity cannot be reduced to a finite identity by
fixing only the level and the $K_\infty$-types: it is false that only
finitely many cuspidal representations have a vector fixed by a given
open compact subgroup $L_f$ with $K_\infty$-types in a given finite set.
For instance, for $k={\mathbb Q}$, $L_f=SL(2,\widehat{\mathbb Z})$, and the
trivial $SO(2)$-type, there are infinitely many spherical Maass cusp forms
of level one, with Laplace eigenvalues tending to infinity (Weyl's law
for the cuspidal spectrum). What is true is the finiteness for a
\emph{fixed infinitesimal character} (Lemma~\ref{lem:finiteness}), and
Hypothesis~\ref{hyp:GH} is refined accordingly.

(2) One would hope to derive the refined identity from the unrefined one
by spectral multipliers at the archimedean places: for a $W$-invariant
Paley--Wiener function $\alpha$ there is an operation
$f_\infty\mapsto f_{\infty,\alpha}$ on $K_\infty$-finite functions in
$C_c^\infty(G_\infty)$ with $\tr\pi(f_{\infty,\alpha})=\alpha(\nu_\pi)\tr\pi(f_\infty)$
(Arthur's multipliers, \cite{APW}). If each archimedean class ${\cal C}_v$
were stable under such multipliers, applied compatibly to $f_v$ and
$\overline f_v$ (compatibility being governed by the relation between
infinitesimal characters in (LH4)), then the unrefined identity would hold
with $f$ replaced by $f_\alpha$ for all $\alpha$, and one could hope to
isolate a single $\nu$ using the absolute convergence of the spectral
expansions. Neither the stability of the classes nor the isolation
argument is carried out here.
\end{remark}

\begin{remark}\label{rem:GH}
(1) Hypothesis~\ref{hyp:GH} is the kind of statement that a full
comparison of the trace formulas of Appendix~\ref{ch:tf} would be expected
to give, refined as in Remark~\ref{rem:why-refine}(2), if the geometric
sides agreed term by term for corresponding pairs. The spectral sides
contain, besides the cuspidal terms, the residual spectrum (the trivial
representation of $G({\mathbb A})$; the even theta representations of the
cover), the cuspidal theta series $\Theta_0$, and the terms coming from
Eisenstein series. A comparison would have to show that these
contributions cancel. The decomposition-of-measures argument of Langlands
\cite{L4}, as adapted by Flicker in \cite[\S\S 4.1--4.2]{Fl}, is designed
to prove such a cancellation (see also \cite[\S\S 10.5, 14.5]{R}). Whether
it applies here has not been checked, and
Corollary~\ref{cor:inconsistent} shows that it cannot succeed for the
correspondence of Definition~\ref{def:corresponding}.

(2) A term-by-term comparison of the geometric sides requires matching
of orbital integrals and of weighted orbital integrals (or, in the
invariant trace formula, of the invariant distributions $I_M$) at every
place. Chapter~\ref{ch:local} provides this only in part. For spherical
functions outside $S$, the Satake isomorphism gives the matching of
hyperbolic orbital integrals but not of elliptic ones
(Remark~\ref{rem:satake-oi}). Weighted orbital integrals of spherical
functions are matched at most for the unit elements
(Remark~\ref{rem:weighted-units}), while Definition~\ref{def:corresponding}
lets $f_v$ range over the whole spherical Hecke algebra for $v\notin S$.
Weighted orbital integrals cannot be matched for arbitrary corresponding
pairs (Remark~\ref{rem:weighted-noninv}), and nothing is matched at the
dyadic places.
\end{remark}

\section{Formal consequences of the hypotheses}\label{sec:main}

We state the consequence of the two hypotheses, prove two lemmas, and then
prove the proposition.

\begin{proposition}\label{thm:main}
Let $\pi=\otimes_v\pi_v$ be a cuspidal automorphic representation of
$G({\mathbb{A}})$ and let $S=S(\pi)$. Suppose that for each $v\in S$
Hypothesis~\ref{hyp:LH} holds for $\pi_v$ with a class ${\cal C}_v$, and
that Hypothesis~\ref{hyp:GH} holds for $S$ and $({\cal C}_v)_{v\in S}$.
Then:
\begin{enumerate}
\item[(a)] $m(\pi)=1$;
\item[(b)] there is a unique irreducible subrepresentation
$\overline{\pi}=\otimes_v\overline{\pi}_v$ of ${\cal A}_{00}$ such that
$\overline{\pi}_v$ corresponds to
$\pi_v$ under the Satake correspondence of \S\ref{sec:spherical} for
$v\notin S$, and $t_v(\overline{\pi}_v)\cong\pi_v$ for $v\in S$.
\end{enumerate}
\end{proposition}

By Theorem~\ref{thm:nonexist} below, conclusion (b) never holds. The
proposition is therefore used only through its contrapositive
(Corollary~\ref{cor:inconsistent}).

\begin{lemma}[finiteness]\label{lem:finiteness}
Let $L_f$ be an open compact subgroup of $G({\mathbb{A}}_f)$, $E$ a
finite set of $K_\infty$-types, and $\nu$ a character of
$\mathfrak z_\infty$. Then only finitely many cuspidal
automorphic representations $\pi'$ of $G({\mathbb{A}})$, counted with
multiplicity, have $\nu_\infty(\pi')=\nu$ and a nonzero vector which is
$L_f$-fixed and whose $K_\infty$-types lie in $E$. The same holds for
genuine cuspidal automorphic representations of
$\overline{G({\mathbb{A}})}$, with $L_f$ replaced by an open compact
subgroup of $\overline{G({\mathbb{A}}_f)}$ meeting $\mu_2$ trivially and
$E$ by a finite set of genuine $\overline{K_\infty}$-types.
\end{lemma}

\begin{proof}
Distinct copies of such representations in the cuspidal spectrum
contribute linearly independent vectors to the space of cusp forms which
are $L_f$-invariant, have $K_\infty$-types in $E$, and are annihilated by
the kernel of $\nu$ in $\mathfrak z_\infty$. That space is
finite-dimensional by Harish-Chandra's finiteness theorem \cite{HC68}
(see also \cite{BJ}, \cite{MW}). For the cover, one reduces as usual to
finitely many quotients $\Gamma\backslash\overline{G_\infty}$ with
$\Gamma$ arithmetic, and Harish-Chandra's argument applies. We have not
written out the details (Remark~\ref{rem26:finiteness-cover}). Without the
condition on $\nu$ the statement is false (Remark~\ref{rem:why-refine}(1)).
\end{proof}

\begin{remyy}\label{rem26:finiteness-cover}
The case of the cover in Lemma~\ref{lem:finiteness} has not been fully
verified. The reduction of genuine cusp forms of fixed level to finitely
many quotients $\Gamma\backslash\overline{G_\infty}$, and the application of
Harish-Chandra's theorem to the nonlinear group $\overline{G_\infty}$, are
standard but have not been written out here.
\end{remyy}

\begin{lemma}[tensor products]\label{lem:tensor}
For $v$ in a finite set $S$, let $V_v$ be a complex vector space and
$\Lambda_v$ a finite linearly independent set of linear functionals on
$V_v$. Then the functionals $\otimes_{v\in S}\lambda_v$
($\lambda_v\in\Lambda_v$) on $\otimes_{v\in S}V_v$ are linearly
independent.
\end{lemma}

\begin{proof} Choose $x_{\mu}\in V_v$ ($\mu\in\Lambda_v$) with
$\lambda(x_\mu)=\delta_{\lambda\mu}$ for $\lambda,\mu\in\Lambda_v$, and
evaluate a linear relation at $\otimes_v x_{\mu_v}$.
\end{proof}

\begin{proof}[Proof of Proposition~\ref{thm:main}]
Let $\nu=\nu_\infty(\pi)$.
For each $v\in S$ use (LH3) and (LH0) to fix an index $L_v$ such that
$\tr \pi_v$ does not vanish identically on ${\cal C}_v(L_v)$.
Restrict attention to $(S,{\cal C})$-corresponding pairs with
$(f_v,\overline{f}_v)\in{\cal C}_v(L_v)$ for $v\in S$. Let ${\cal H}^S$ be
the restricted tensor product of the spherical Hecke algebras
${\cal H}(G(k_v),K_{0,v})$ for $v\notin S$. For $h\in{\cal H}^S$, let
$\overline{h}$ denote the corresponding genuine function, and write
$f=f_S\otimes h$, $\overline{f}=\overline{f}_S\otimes\overline{h}$.

{\em Step 1 (finiteness).} If $\tr \pi'(f)\neq 0$ for some such $f$,
then $\pi'_v$ has a nonzero $K_{0,v}$-fixed vector for $v\notin S$, and
$\pi'_S=\otimes_{v\in S}\pi'_v$ has a nonzero vector fixed by
$\prod_{v\in S, v<\infty}L_v$ with $K_\infty$-types in $\prod_{v|\infty}L_v$.
By Lemma~\ref{lem:finiteness}, only a finite set ${\cal A}$ of cuspidal
automorphic representations $\pi'$ with $\nu_\infty(\pi')=\nu$ can
contribute to $\tr R^\nu_{cusp}(f)$. Similarly, a genuine
$\overline\pi'\subset{\cal A}_{00}$
contributing to $\tr\overline R{}^{\,\nu}_{00}(\overline f)$ has a nonzero
vector fixed by the image of a splitting of an open compact subgroup of
$\overline{G({\mathbb A}_f)}$ determined by the $L_v$ ($v\in S$ finite) and
the $\kappa_v(K_{0,v})$ ($v\notin S$), with genuine $\overline{K_\infty}$-types
in a finite set determined by the $L_v$ ($v|\infty$), and since it lies
over $\nu$, (LH4) allows only finitely many infinitesimal characters for
$\overline\pi'_\infty$. By Lemma~\ref{lem:finiteness} again, only a finite
set $\overline{{\cal A}}$ of such $\overline\pi'$ can contribute.
Hypothesis~\ref{hyp:GH} for this $\nu$
therefore reads
\begin{equation}\label{eq:finite-identity}
\sum_{\pi'\in{\cal A}}m(\pi')\,\omega_{\pi'}^S(h)
\prod_{v\in S}\tr \pi'_v(f_v)
=\sum_{\overline{\pi}'\in\overline{{\cal A}}}m(\overline{\pi}')\,
\omega_{\overline{\pi}'}^S(\overline{h})
\prod_{v\in S}\tr \overline{\pi}'_v(\overline{f}_v),
\end{equation}
where $\omega_{\pi'}^S$ (resp. $\omega_{\overline{\pi}'}^S$) is the
character of ${\cal H}^S$ (resp. of its genuine analog) by which it acts
on the line of spherical vectors of $\otimes_{v\notin S}\pi'_v$ (resp. of
$\otimes_{v\notin S}\overline{\pi}'_v$).

{\em Step 2 (isolating the components outside $S$).} By
\S\ref{sec:spherical}, $\omega_{\overline{\pi}'}^S(\overline{h})$ is a
character of ${\cal H}^S$ as a function of $h$, and it equals
$\omega_{\pi''}^S(h)$ for the unramified representation
$\otimes_{v\notin S}\pi''_v$ corresponding to
$\otimes_{v\notin S}\overline{\pi}'_v$ under the Satake correspondence.
Fix $f_S$ and regard both sides of (\ref{eq:finite-identity}) as
functions of $h$. Distinct characters of an algebra are linearly
independent (Artin), so, after grouping the finitely many terms
according to their characters, we may equate the coefficients of the
character $\omega_\pi^S$. Since a $K_{0,v}$-spherical irreducible
representation is determined by its Hecke character, and the Satake
correspondence is a bijection on spherical representations, this gives
\begin{equation}\label{eq:S-identity}
\sum_{\substack{\pi'\in{\cal A}\\ \pi'_v\cong\pi_v\ (v\notin S)}}
m(\pi')\prod_{v\in S}\tr \pi'_v(f_v)
=\sum_{\substack{\overline{\pi}'\in\overline{{\cal A}}\\
\overline{\pi}'_v\leftrightarrow\pi_v\ (v\notin S)}}
m(\overline{\pi}')\prod_{v\in S}\tr \overline{\pi}'_v(\overline{f}_v),
\end{equation}
where $\overline{\pi}'_v\leftrightarrow\pi_v$ means that they correspond
under the Satake correspondence. This holds for all
$(f_v,\overline{f}_v)\in{\cal C}_v(L_v)$, $v\in S$, and hence, by
linearity, on $\otimes_{v\in S}{\cal C}_v(L_v)$.

{\em Step 3 (transfer).} For $\overline\pi'\in\overline{\cal A}$ and
$v\in S$ we have $\overline\pi'_v\in\overline{\mathfrak{F}}_v$. By (LH2),
each term on the right of
(\ref{eq:S-identity}) either vanishes identically or equals
$m(\overline{\pi}')\prod_{v\in S}\tr t_v(\overline{\pi}'_v)(f_v)$.
The representations $\overline{\pi}'$ occurring on the right are
isomorphic to one another outside $S$ (the genuine Satake correspondence
is injective), so two distinct ones differ at some $v\in S$, and by the
injectivity of $t_v$ the corresponding tuples
$(t_v(\overline{\pi}'_v))_{v\in S}$ are distinct. On the left, distinct
$\pi'$ occurring in (\ref{eq:S-identity}) have distinct tuples
$(\pi'_v)_{v\in S}$, since they agree outside $S$, and
$\pi'_v\in\mathfrak{F}_v$.

{\em Step 4 (comparison of coefficients).} For each $v\in S$, let
$\Lambda_v$ be the set of distinct representations of $G(k_v)$ which
occur as $\pi'_v$ on the left or as $t_v(\overline{\pi}'_v)$ on the right
of (\ref{eq:S-identity}), and whose characters do not vanish identically
on ${\cal C}_v(L_v)$. It is finite by Step 1. Terms involving some other
representation vanish on $\otimes_{v\in S}{\cal C}_v(L_v)$. By (LH1), the
characters of the elements of $\Lambda_v$ are linearly independent on
${\cal C}_v(L_v)$, so by Lemma~\ref{lem:tensor} the product functionals
are linearly independent on $\otimes_{v\in S}{\cal C}_v(L_v)$. By the
choice of the $L_v$, $\pi_v\in\Lambda_v$ for every $v\in S$. Comparing
the coefficients of $\otimes_{v\in S}\tr \pi_v$ on the two sides of
(\ref{eq:S-identity}) gives
\[
m(\pi)=\sum_{\overline{\pi}'}m(\overline{\pi}'),
\]
the sum over the $\overline{\pi}'$ in $\overline{\cal A}$ with
$\overline{\pi}'_v\leftrightarrow\pi_v$ for
$v\notin S$ and $t_v(\overline{\pi}'_v)\cong\pi_v$ for $v\in S$. (Every
irreducible subrepresentation $\overline\pi'$ of ${\cal A}_{00}$ with these
properties lies over $\nu$, so it does lie in $\overline{\cal A}$.)
By Step 3, there is at most one such $\overline{\pi}'$, and
$m(\overline{\pi}')\leq 1$ by Waldspurger's theorem. Hence $m(\pi)\leq 1$.
Since $\pi$ is automorphic, $m(\pi)\geq 1$, which proves (a). Then the
sum on the right equals $1$, so it has exactly one term, which proves (b).
\end{proof}

\begin{remark}
Hypothesis~\ref{hyp:GH} is only used for $\nu=\nu_\infty(\pi)$ and for
pairs with $f_S$ in $\otimes_{v\in S}{\cal C}_v(L_v)$ and arbitrary
$h\in{\cal H}^S$.
\end{remark}

\section{The Satake correspondence is not global}\label{sec:nonexist}

This section proves that the correspondence in
Proposition~\ref{thm:main}(b) does not exist. The proof compares the
unramified parameters on the two groups.

Let $v$ be an odd place, $F=k_v$, and let $\psi$ be a nontrivial additive
character of $F$ of conductor ${\cal O}$. Let $\chi_\psi=\omega$ be the
function on $F^\times$ in formula (iii) of \S\ref{sec:weil}, for the
quadratic form $x^2$. We use the following standard properties
(\cite{Ge2}, \cite{Rao}): $\chi_\psi(x)$ depends only on the class of $x$
modulo squares, $\chi_\psi(xy)=(x,y)\chi_\psi(x)\chi_\psi(y)$, and
$\chi_\psi=1$ on ${\cal O}^\times$. By Lemma~\ref{lem:Acocycle}(a),
$(a(x),\zeta)\mapsto\zeta\,\chi_\psi(x)\mu(x)$ is a genuine character of
$\overline{A(F)}$ for every character $\mu$ of $F^\times$, and every
genuine character is of this form for a unique $\mu$.

\begin{lemma}\label{lem:square}
Let $\chi$ be an unramified character of $A(F)$, let $\overline\chi$ be
the corresponding genuine character (Definition~\ref{def:corrchar}), and
define $\mu$ by $\overline\chi(a(x),1)=\chi_\psi(x)\mu(x)$. Then $\mu$ is
an unramified character of $F^\times$ and
\[
\mu(\varpi)^2=\chi(a(\varpi))^2 .
\]
\end{lemma}

\begin{proof}
Write $\overline\psi(x)=\overline\chi(a(x),1)$, so that
$\overline\psi(xy)=(x,y)\overline\psi(x)\overline\psi(y)$ and
$\overline\psi=\chi_1$ on the elements of even valuation. Since $\chi$ is
unramified, $\chi_1$ is trivial on ${\cal O}^\times$, and so is $\chi_\psi$.
Hence $\mu$ is unramified. Next,
$\overline\psi(\varpi)^2=(\varpi,\varpi)\overline\psi(\varpi^2)
=(\varpi,\varpi)\chi_1(\varpi^2)$ and
$\chi_\psi(\varpi)^2=(\varpi,\varpi)\chi_\psi(\varpi^2)=(\varpi,\varpi)$.
Dividing, $\mu(\varpi)^2=\chi_1(\varpi^2)=\chi(a(\varpi^2))$.
\end{proof}

Thus the correspondence of Definition~\ref{def:corrchar} is
$\mu(\varpi)=\pm\chi(a(\varpi))$, with a sign that depends on $\psi$,
$\varpi$ and $\eta_\varpi$.

\begin{lemma}\label{lem:A00}
Let $\pi$ be a cuspidal automorphic representation of $G({\mathbb{A}})$,
let $S\supset S(\pi)$ be finite, and let $\overline{\pi}'$ be a genuine
cuspidal automorphic representation of $\overline{G({\mathbb{A}})}$ such
that $\overline{\pi}'_v$ corresponds to $\pi_v$ under the Satake
correspondence for every $v\notin S$. Then no irreducible constituent of
$\Theta_0$ is isomorphic to $\overline\pi'$. Consequently every copy of
$\overline\pi'$ in the genuine cuspidal spectrum lies in ${\cal A}_{00}$.
\end{lemma}

\begin{proof}
We use three facts.

(i) Every irreducible constituent $\sigma$ of $\Theta_0$ is a quotient of
a global Weil representation $\otimes_v r_{\psi_v,\chi_v}$, and so
$\sigma_v$ is the even part $r^+_{\psi_v,\chi_v}$ of the local Weil
representation for almost all $v$ \cite{GePS}, \cite{Ge2}.

(ii) Let $v$ be odd, with $\psi_v$ of conductor ${\cal O}_v$ and $\chi_v$
unramified. By formula (iii) of \S\ref{sec:weil}, the functional
$\phi\mapsto\phi(0)$ identifies the Jacquet module of $r^+_{\psi_v,\chi_v}$
with the genuine character $(h(a),\zeta)\mapsto\zeta|a|^{1/2}\omega(a)$ of
$\overline{A(k_v)}$. After normalizing by $\delta_B^{-1/2}(h(a))=|a|^{-1}$,
Frobenius reciprocity embeds $r^+_{\psi_v,\chi_v}$ (which is irreducible
and has a vector fixed by $\kappa_v(K_{0,v})$, namely the characteristic
function of ${\cal O}_v$) in $I_A(\overline\chi)$ with
$|\overline\chi(a(x),1)|=|x|^{-1/2}$. Under the Satake correspondence of
Definition~\ref{def:satake-corr}, which does not change $\chi_1$ on
$A_2(k_v)$, it therefore corresponds to the spherical constituent of
$I_A(\chi)$ for a character $\chi$ with $|\chi(a(x))|=|x|^{\pm1/2}$.

(iii) By Labesse--Langlands \cite{LL}, $\pi$ is a constituent of the
restriction to $SL(2,{\mathbb{A}})$ of a cuspidal automorphic
representation $\Pi$ of $GL(2,{\mathbb{A}})$. At each place $v$ where $\Pi_v$
is unramified, $\Pi_v\cong I(\mu_1,\mu_2)$ with
$|\mu_i(x)|=|x|^{\pm s_v}$ and $|s_v|<1/4$ \cite{GJ78}. Hence $\pi_v$ is
the spherical constituent of $I_A(\chi)$ with $\chi(a(x))=\mu_1\mu_2^{-1}(x)$,
and $|\chi(a(x))|=|x|^{\pm 2s_v}$ with $|2s_v|<1/2$.

Now let $\sigma$ be an irreducible constituent of $\Theta_0$. Choose an odd
place $v\notin S$ at which $\Pi_v$, $\psi_v$, $\chi_v$ are unramified and
$\sigma_v\cong r^+_{\psi_v,\chi_v}$. All but finitely many places qualify.
If $\sigma\cong\overline\pi'$, then $r^+_{\psi_v,\chi_v}$ would correspond
to $\pi_v$ under the Satake correspondence, which by (ii) and (iii)
is impossible, because a spherical representation determines the
character $\chi$ of the principal series containing it up to
$\chi\mapsto\chi^{-1}$. Hence $\overline\pi'$ is isomorphic to no
constituent of $\Theta_0$. Since $\Theta_0$ is a closed invariant
subspace of the (discretely decomposable) genuine cuspidal spectrum, the
orthogonal projection to $\Theta_0$ of any copy of $\overline\pi'$ is
zero, so every copy lies in ${\cal A}_{00}$.
\end{proof}

\begin{remyy}\label{rem26:A00-norm}
The proof of Lemma~\ref{lem:A00} uses the bound $|s_v|<1/4$ only to
exclude $|s_v|=1/4$ at a single place $v\notin S$ where all the data are
unramified. The trivial bound $|s_v|<1/2$ does not suffice. The
normalizations in (ii) follow directly from
$r_+(h(a),\zeta)\phi(0)=\zeta|a|^{1/2}\omega(a)\phi(0)$
and $\delta_B^{1/2}(h(a))=|a|$.
The source-dependent identification of the resulting spherical
constituent remains to be checked against \cite{Ge2} and
\cite{GePS}.
\end{remyy}

\begin{theorem}\label{thm:nonexist}
Let $\pi$ be a cuspidal automorphic representation of $G({\mathbb A})$ and
let $S\supseteq S(\pi)$ be a finite set of places. There is no genuine
cuspidal automorphic representation $\overline\pi'$ of
$\overline{G({\mathbb A})}$ such that $\overline\pi'_v$ corresponds to
$\pi_v$ under the Satake correspondence of
Definition~\ref{def:satake-corr} for every $v\notin S$. This holds for
every choice of the uniformizers $\varpi_v$ and of the constants
$\eta_{\varpi_v}$.
\end{theorem}

\begin{proof}
We use the following results.

(W) (Waldspurger \cite{Wa1}, \cite{Wa2}; see also \cite{PS}.) Fix a
nontrivial character $\psi$ of $k\backslash{\mathbb A}$. To each
irreducible subrepresentation $\overline\pi'$ of ${\cal A}_{00}$ there is
attached a cuspidal automorphic representation
$\sigma=\mathrm{Wald}_\psi(\overline\pi')$ of $PGL(2,{\mathbb A})$. For
almost all $v$, if $\overline\pi'_v$ is the spherical constituent of
$I_A(\overline\chi_v)$ with
$\overline\chi_v(a(x),1)=\chi_{\psi_v}(x)\mu_v(x)$, then $\sigma_v$ is the
unramified principal series $\pi(\mu_v,\mu_v^{-1})$ of $PGL(2,k_v)$.

(GJ) (Gelbart--Jacquet \cite{GJ78}.) For a cuspidal automorphic
representation $\Pi$ of $GL(2,{\mathbb A})$ there is an automorphic
representation $\mathrm{Ad}(\Pi)$ of $GL(3,{\mathbb A})$ whose Satake
parameter at each unramified place is
$\{\alpha_v\beta_v^{-1},1,\alpha_v^{-1}\beta_v\}$ if $\{\alpha_v,\beta_v\}$
is that of $\Pi_v$. It is self-dual, and it is either cuspidal or an
isobaric sum of unitary cuspidal representations of $GL(1)$ and $GL(2)$.
For a sufficiently large finite set $S'$ of places,
$L^{S'}(s,\Pi\times\Pi^\vee)=\zeta_k^{S'}(s)L^{S'}(s,\mathrm{Ad}\,\Pi)$.

(JS) (Jacquet--Shalika \cite{JS}; see also \cite{Sha}.) For unitary
cuspidal representations $\tau,\tau'$ of general linear groups, the
partial Rankin--Selberg $L$-function $L^{S'}(s,\tau\times\tau'^\vee)$ is
holomorphic and nonzero at $s=1$ if $\tau\not\cong\tau'$, and has a simple
pole at $s=1$ if $\tau\cong\tau'$.

Suppose that $\overline\pi'$ exists. By Lemma~\ref{lem:A00} it is a
subrepresentation of ${\cal A}_{00}$. Let $\sigma$ be as in (W). Let $\Pi$
be a cuspidal automorphic representation of $GL(2,{\mathbb A})$ whose
restriction contains $\pi$ \cite{LL}, and for $v$ outside a finite set
write $\Pi_v=I(\mu_{1,v},\mu_{2,v})$ and
$t_v=\mu_{1,v}(\varpi_v)\mu_{2,v}(\varpi_v)^{-1}$, so that $\pi_v$ is the
spherical constituent of $I_A(\chi_v)$ with $\chi_v(a(\varpi_v))=t_v$.
By Lemma~\ref{lem:square}, $\mu_v(\varpi_v)^2=t_v^2$ for almost all $v$,
whatever choices were made in Definition~\ref{def:corrchar}. The Satake
parameters are therefore
\[
\mathrm{Ad}(\Pi)_v:\ \{t_v,1,t_v^{-1}\},\qquad
\mathrm{Ad}(\sigma)_v:\ \{t_v^2,1,t_v^{-2}\}.
\]
The nine products of two parameters of $\mathrm{Ad}(\Pi)_v$ are
\[
\{t_v^{2},1,t_v^{-2}\}\ \cup\ \{t_v,1,t_v^{-1}\}\ \cup\ \{t_v,1,t_v^{-1}\},
\]
so for a sufficiently large finite set $S'$
\begin{equation}\label{eq:L-identity}
L^{S'}\bigl(s,\mathrm{Ad}(\Pi)\times\mathrm{Ad}(\Pi)\bigr)
=L^{S'}\bigl(s,\mathrm{Ad}(\sigma)\bigr)\,
L^{S'}\bigl(s,\mathrm{Ad}(\Pi)\bigr)^2 .
\end{equation}
Since $\sigma$ and $\Pi$ are cuspidal,
$L^{S'}(s,\sigma\times\sigma^\vee)$ and $L^{S'}(s,\Pi\times\Pi^\vee)$ have
simple poles at $s=1$ by (JS). The identities
$L^{S'}(s,\sigma\times\sigma^\vee)
=\zeta_k^{S'}(s)L^{S'}(s,\mathrm{Ad}\,\sigma)$
and
$L^{S'}(s,\Pi\times\Pi^\vee)
=\zeta_k^{S'}(s)L^{S'}(s,\mathrm{Ad}\,\Pi)$,
together with the simple poles in (JS) and the simple pole of
$\zeta_k^{S'}(s)$, show that
$L^{S'}(s,\mathrm{Ad}\,\sigma)$ and
$L^{S'}(s,\mathrm{Ad}\,\Pi)$ are holomorphic (indeed nonzero)
at $s=1$. This also covers dihedral $\Pi$, for which
$\mathrm{Ad}(\Pi)$ need not be cuspidal.
Consequently, the right side of
(\ref{eq:L-identity}) is holomorphic at $s=1$. On the other hand, write
$\mathrm{Ad}(\Pi)$ as an isobaric sum of unitary cuspidal representations
$\tau_1,\dots,\tau_r$. Since $\mathrm{Ad}(\Pi)$ is self-dual, the left
side of (\ref{eq:L-identity}) is $\prod_{i,j}L^{S'}(s,\tau_i\times\tau_j^\vee)$.
By (JS) each factor with $\tau_i\cong\tau_j$ has a pole at $s=1$, and the
other factors are holomorphic and nonzero there. So the left side has a
pole at $s=1$. This is a contradiction.
\end{proof}

\begin{remyy}\label{rem26:nonexist}
The proof of Theorem~\ref{thm:nonexist} was written during the 2026
revision. The preceding proof applies equally to dihedral and non-dihedral
$\Pi$; it uses the isobaric decomposition of $\mathrm{Ad}(\Pi)$,
not its cuspidality. Carbon's 2026 referee report independently
checked the pole-counting deduction. Nevertheless, the precise
normalizations of (W), (GJ) and (JS) have not been fully
rechecked from the primary sources in this revision. In particular the
compatibility of Waldspurger's correspondence with the parametrization
$\overline\chi=\chi_\psi\mu$ at unramified places, in the normalization of
this text, should be checked against \cite{Wa2}.
\end{remyy}

\begin{corollary}\label{cor:inconsistent}
Let $\pi$ be a cuspidal automorphic representation of $G({\mathbb A})$
and $S=S(\pi)$. There are no classes $({\cal C}_v)_{v\in S}$ such that
Hypothesis~\ref{hyp:LH} holds for each $\pi_v$ with ${\cal C}_v$ and
Hypothesis~\ref{hyp:GH} holds for $S$ and $({\cal C}_v)_{v\in S}$. In
particular, if Hypothesis~\ref{hyp:LH} holds at the places of $S$, then
the trace identity of Hypothesis~\ref{hyp:GH} is false for
$\nu=\nu_\infty(\pi)$.
\end{corollary}

\begin{proof}
Combine Proposition~\ref{thm:main}(b) with Theorem~\ref{thm:nonexist}.
\end{proof}

\begin{remyy}\label{rem:even}
\emph{The even correspondence.} Lemma~\ref{lem:square} and (W) explain
Theorem~\ref{thm:nonexist}. Write $t'=\mu(\varpi)$ for the unramified
parameter of a genuine representation and $t=\chi(a(\varpi))$ for that of
a representation of $SL(2,F)$. The Satake correspondence of
\S\ref{sec:spherical} is $t=\pm t'$. Waldspurger's correspondence followed
by restriction from $PGL(2,F)$ to $SL(2,F)$ sends
$I_A(\chi_\psi\mu)$ to the spherical constituent of $I_A(\chi)$ with
$\chi(a(x))=\mu(x)^2$, that is, $t=t'^{\,2}$. On dual groups this is the
projection $SL(2,{\mathbb C})\rightarrow PGL(2,{\mathbb C})$. It does not
depend on $\psi$, because replacing $\psi$ by $x\mapsto\psi(ax)$ multiplies
$\mu$ by the quadratic character $(a,\cdot)_2$.

The two rules behave differently. Under $t=t'^{\,2}$ the genuine
complementary series ($|t'|=q^{-s}$, $0<s<1/2$) corresponds to the
complementary series of $SL(2,F)$ ($0<2s<1$), and the even Weil
representation ($s=1/2$) to the trivial representation. Under
$t=\pm t'$ the even Weil representation corresponds to an irreducible
complementary series representation, as in the proof of
Lemma~\ref{lem:A00}. At the real place, the elliptic matchings of
\S\ref{sec:realds} compare $r(\theta)$ with the element over
$r(2\theta)$ and double the infinitesimal character. They are of the
even type, and Lemma~\ref{lem:realds}(a) is the archimedean form of the
Shimura correspondence (weight $k/2$ to weight $k-1$). The hyperbolic
matchings of \S\ref{sec:charsA}--\S\ref{sec:spherical} and
\S\ref{sec:realA} are of the odd type. Hypothesis~\ref{hyp:GH} combines
the two types, and Corollary~\ref{cor:inconsistent} shows that this
combination is impossible.

One may ask whether Proposition~\ref{thm:main} has an analog in which
every place is matched by the even correspondence. We have not found one
that yields multiplicity one for $SL(2)$, for three reasons.
\begin{enumerate}
\item The map $t'\mapsto t'^{\,2}$ is two-to-one, so the transfer is not
injective: $\mu$ and its product with the unramified quadratic character
have the same image. In Step 2 of the proof the right side becomes a sum
over genuine representations whose images under (W) are quadratic twists
of one another, and the uniqueness argument of Step 3 is lost.
\item Only those $\pi$ are reached for which $\Pi$ can be twisted to have
trivial central character. For the other $\pi$ there is no
$\overline\pi$, so an identity between the full cuspidal traces cannot
hold.
\item Characters of the cover are related to characters of $PGL(2)$
(\cite{Wa2}, \cite{Ad2}, \cite{Li19}), and the restriction of a character
of $PGL(2,F)$ to $SL(2,F)$ is the sum of the characters in an $L$-packet.
A comparison of this kind therefore controls sums of multiplicities over
$L$-packets, weighted by the signs in Waldspurger's description of
${\cal A}_{00}$, and not the multiplicity of a single $\pi$. By
\cite{LL}, $m(\pi)>1$ would require two cuspidal representations of
$GL(2,{\mathbb A})$ that are twists of each other locally at every place
but not globally. On the cover these give different near-equivalence
classes, about whose relation Waldspurger's multiplicity one theorem says
nothing.
\end{enumerate}
This discussion is heuristic and has not been verified in detail.
\end{remyy}

\begin{example}
For the running example $\pi$ of Chapter~\ref{ch:intro} (from $\Delta$),
let $\tau(p)=p^{11/2}(\alpha_p+\alpha_p^{-1})$, so that $t_p=\alpha_p^2$.
The genuine representation attached to $\Delta$ by the Shimura
correspondence is generated by a cusp form of weight $13/2$, in agreement
with the example following Definition~\ref{def:infchar}. Its image under
(W) is the representation of $PGL(2,{\mathbb A})$ generated by $\Delta$,
so its unramified parameters are $t'_p=\pm\alpha_p$, and
$t_p=t_p'^{\,2}$. Theorem~\ref{thm:nonexist} says that no genuine cusp
form has parameters $t'_p=\pm\alpha_p^2$ for almost all $p$.
\end{example}

\section{Where the global comparison fails}\label{sec:where-fails}

This section traces the proposed comparison through the two trace
formulas and locates the terms that match and the terms that do not. It
explains Corollary~\ref{cor:inconsistent} rather than proving anything
new. The statements about intertwining operators and about the residual
spectrum of the cover are quoted from memory
(Remark~\ref{rem26:where-fails}). Proofs of them are in the companion
text \cite{JC}, as \S\ref{sec:cover-supplies} explains.

\subsection*{The proposed argument}

The argument has five steps.
\begin{enumerate}
\item Write the trace formula for $G({\mathbb A})$ and for the cover
(Appendix~\ref{ch:tf}). Each equates a geometric side, a sum over
conjugacy classes in $G(k)$ of orbital integrals and weighted orbital
integrals, with a spectral side, the trace of the discrete spectrum plus
terms coming from Eisenstein series. Both geometric sides are indexed by
classes in $G(k)$, because the cover splits over $G(k)$.
\item Choose $f=\otimes_vf_v$ and $\overline f=\otimes_v\overline f_v$ whose
local orbital integrals correspond \emph{directly}, that is, at an element
of $G(k_v)$ and at a lift of the same element, so that the two geometric
sides agree term by term.
\item Conclude that the spectral sides agree. Show that the Eisenstein
terms cancel each other, so that the traces of the discrete spectra agree.
\item Separate the identity according to the eigencharacters of the Hecke
algebras outside $S$, and obtain
$m(\pi)=\sum m(\overline\pi')$ over the genuine $\overline\pi'$ that
correspond to $\pi$.
\item Apply Waldspurger's multiplicity one theorem to the right side.
\end{enumerate}
Steps 4 and 5 are sound: they are Proposition~\ref{thm:main}. The failure
is in Steps 2 and 3. We look at the terms one family at a time.

\subsection*{The torus: the direct comparison is globally consistent}

Fix a nontrivial character $\psi=\otimes_v\psi_v$ of
$k\backslash{\mathbb A}$, and let $\chi_{\psi_v}$ be as in
\S\ref{sec:nonexist}, now at every place $v$. By Weil's product formula
for the Weil index \cite{We3},
\[
\prod_v\chi_{\psi_v}(x)=1\qquad\text{for }x\in k^\times .
\]
Hence $(a(x),\zeta)\mapsto\zeta\prod_v\chi_{\psi_v}(x_v)\mu_v(x_v)$ is a
genuine character of $\overline{A({\mathbb A})}$ trivial on $A(k)$ exactly
when $\mu=\otimes_v\mu_v$ is a Hecke character of $k$. So the genuine
automorphic characters of the cover of the torus correspond bijectively to
the automorphic characters $a(x)\mapsto\mu(x)$ of $A({\mathbb A})$.

Locally, put
$\overline\phi_v(a(x),1)=\chi_{\psi_v}(x)^{-1}\phi_v(a(x))$. Then
$\overline\phi_v^{\,\wedge}(\chi_{\psi_v}\mu_v)=\phi_v^{\wedge}(\mu_v)$ for
every $\mu_v$. For $p$ odd this is the matching of
Definition~\ref{def:corrA} with $\eta_\varpi$ chosen according to $\psi_v$,
for $k_v={\mathbb R}$ it is that of Lemma~\ref{lem:realA-match}, and the
same formula makes sense at the dyadic places. Suppose that $f_v$ and
$\overline f_v$ correspond in this way on the hyperbolic set at every
place. An element $\gamma=a(x)$, $x\in k^\times$, of $G(k)$ has the local
components $(a(x),1)$ in the cover, since $s_v(a(x))=1$ in
(\ref{eq:sv}). So, by the product formula,
\[
\prod_vF_{\overline f_v}(\gamma,1)=\prod_v\chi_{\psi_v}(x)^{-1}F_{f_v}(\gamma)
=\prod_vF_{f_v}(\gamma)\qquad(x\neq\pm1).
\]
The regular hyperbolic terms of the two geometric sides agree. On the
spectral side, $\tr I_A(\chi_\psi\mu)(\overline f)=\tr I_A(\mu)(f)$ for
every Hecke character $\mu$, as in Proposition~\ref{prop:hyp-ps}.

Thus the direct comparison is the correct comparison for the torus: for
its automorphic characters, for the orbital integrals over it, and for the
traces of the representations induced from it. This is why it looks
right.

\subsection*{First failure: intertwining operators and the residual spectrum}

The Eisenstein terms of the trace formula are integrals over $\mu$ and
over $s\in i{\mathbb R}$ of the traces above, weighted by logarithmic
derivatives of the intertwining operators $M(w,\mu|\cdot|^s)$
(Appendix~\ref{ch:tf}). On unramified vectors these operators act by
scalars. For $G({\mathbb A})$ the scalar is, up to finitely many places,
\[
c(s,\mu)=\frac{L(s,\mu)}{L(1+s,\mu)} ,
\]
and for the cover it is
\[
\overline c(s,\mu)=\frac{L(2s,\mu^2)}{L(1+2s,\mu^2)} ,
\]
the same as for the principal series $\pi(\mu|\cdot|^s,\mu^{-1}|\cdot|^{-s})$
of $PGL(2)$. The direct comparison matches the traces at $(\mu,s)$ on the
two groups, but the weights $c'/c$ and $\overline c\,'/\overline c$ are
different functions of $(\mu,s)$: one involves the zeros of $L(s,\mu)$,
the other those of $L(2s,\mu^2)$. So the Eisenstein terms of the two
formulas do not cancel in Step 3. Under the squaring correspondence,
which sends $(\mu,s)$ on the cover to $(\mu^2,2s)$ on $SL(2)$, the two
scalars agree exactly, since $c(2s,\mu^2)=\overline c(s,\mu)$.

The residual spectrum is read off from the poles. The function $c(s,\mu)$
has a pole for $\mathrm{Re}\,s>0$ only at $s=1$, $\mu=1$. Its residue gives
the trivial representation, with unramified parameter $|t|=q^{-1}$. The
function $\overline c(s,\mu)$ has poles at $s=1/2$ for every quadratic
character $\mu$. Their residues give the even theta representations, one
for each square class in $k^\times$, with $|t'|=q^{-1/2}$. The direct
comparison would pair a theta representation with a representation of
$SL(2,{\mathbb A})$ having $|t|=q^{-1/2}$ at almost all places. There is no
such representation in the discrete spectrum, by the proof of
Lemma~\ref{lem:A00}. The squaring correspondence sends every theta
representation to the trivial representation, since $(q^{-1/2})^2=q^{-1}$.

On the geometric side, the counterpart of the weights $c'/c$ is the
family of weighted orbital integrals and unipotent terms. A direct
matching of orbital integrals does not control these
(Remark~\ref{rem:weighted-noninv}), and the computation above shows that
they cannot all agree, since the spectral sides do not.

\subsection*{Second failure: the discrete spectra are disjoint}

Under the direct comparison, no irreducible constituent of the discrete
spectrum of $G({\mathbb A})$ corresponds, at almost all places, to an
irreducible constituent of the genuine discrete spectrum of the cover.
\begin{itemize}
\item A cuspidal $\pi$ has no cuspidal partner, by
Theorem~\ref{thm:nonexist}.
\item The trivial representation has $|t|=q^{-1}$. A partner would have
$|t'|=q^{-1}$ at almost all places, and the corresponding genuine
principal series constituent is not unitary, because the genuine
complementary series ends at exponent $1/2$.
\item The even theta representations, and the constituents of
$\Theta_0$, have $|t'|=q^{-1/2}$ at almost all places. A partner would
have $|t|=q^{-1/2}$, which excludes the trivial representation, and
excludes cuspidal representations by the bound of Gelbart and Jacquet, as
in Lemma~\ref{lem:A00}.
\end{itemize}
So suppose that, for some class of test functions, the traces of the
discrete spectra did agree. Separating by Hecke eigencharacters as in
Step~2 of the proof of Proposition~\ref{thm:main}, every eigencharacter
occurs on one side only, and its coefficient must vanish. The class would
have to annihilate every representation it was meant to compare. This is
Corollary~\ref{cor:inconsistent} seen from the spectral side.

\subsection*{Diagnosis: the same torus, a different coroot}

The principal series of $SL(2,F)$ and of its cover are both parametrized
by characters $\mu$ of $F^\times$, through $a(x)\mapsto\mu(x)$ and through
$\chi_\psi\mu$. They differ in the rank-one data attached to the root. For
$SL(2,F)$, reducibility and the intertwining operator are governed by
$\mu$, the restriction of the character to the coroot $x\mapsto a(x)$. For
the cover they are governed by $\mu^2$, that is, by the restriction to
$x\mapsto a(x^2)$: the coroot is doubled. On the dual side both tori are
${\mathbb C}^\times$. The dual group of $SL(2)$ is $PGL(2,{\mathbb C})$,
where the root takes the value $t$, and that of the cover is
$SL(2,{\mathbb C})$, where it takes the value $t'^{\,2}$. The direct
comparison is an identification $t=\pm t'$ of the dual tori, compatible
with the Weyl groups. It does not extend to a homomorphism of the dual
groups, because the only nontrivial homomorphism
$SL(2,{\mathbb C})\rightarrow PGL(2,{\mathbb C})$ is the projection, which
is $t=t'^{\,2}$ on the tori.

This is the sense in which the direct comparison is off by a square. It is
correct on the torus and wrong by a factor of $2$ in the direction of the
root. The parts of the trace formula that see only the torus match: the
hyperbolic orbital integrals, the traces of principal series, and the
continuous spectrum as a set. The parts that see the root do not: the
intertwining operators, the residual spectrum, and the cuspidal spectrum.
The elliptic terms of the geometric side are examined in
\S\ref{sec:cover-elliptic}. They fail as well, at the elliptic elements
close to $-1$. The unipotent terms and the weighted orbital integrals
have not been examined. At the real place the discrete series
are attached to the compact torus, where the root cannot be avoided, and
the matchings that transfer them (Lemma~\ref{lem:realds}) are forced to
double the parameter.

\begin{remyy}\label{rem26:where-fails}
The formula for $\overline c(s,\mu)$, the description of the residual
spectrum of the cover by even theta representations, and the form of
Weil's product formula used above are quoted from memory (compare
\cite{GS}, \cite{GePS}, \cite{Ge2}, \cite{We3}) and have not been checked
against the sources. The claim that the hyperbolic matching with
$\chi_{\psi_v}^{-1}$ agrees with Definition~\ref{def:corrA} for a suitable
$\eta_\varpi$ follows from Lemma~\ref{lem:square} and its proof, but the
normalizations of $\chi_{\psi_v}$ have not been compared with those of
\S\ref{sec:weil}. The first three statements are proved in \cite{JC}
(\S\ref{sec:cover-supplies}), which has not been refereed by a
specialist. For $F={\mathbb Q}_3$ and the real field the comparison with
Definition~\ref{def:corrA} and Lemma~\ref{lem:realA-match} is carried out
in Example~\ref{ex:eta-psi}.
\end{remyy}

\begin{example}
Let $k={\mathbb Q}$ and $\mu=1$. For $SL(2)$ the scalar is
$\xi(s)/\xi(1+s)$, where $\xi$ is the completed Riemann zeta function, and
its pole at $s=1$ gives the constant functions. For the cover it is
$\xi(2s)/\xi(1+2s)$, and its pole at $s=1/2$ gives the theta function
$\theta(z)=\sum_ne^{2\pi in^2z}$ of weight $1/2$. The direct comparison
would pair $\theta$ with an automorphic form on $SL(2)$ whose Hecke
eigenvalue at each odd prime $p$ is that of the spherical representation
with $t=p^{-1/2}$. No such form exists in the discrete spectrum. The
squaring correspondence pairs $\theta$ with the constant function.
\end{example}

\section{The local hypothesis at hyperspecial places}

We verify Hypothesis~\ref{hyp:LH} at the odd places where $\pi_v$ is
spherical for a hyperspecial maximal compact subgroup.

\begin{proposition}\label{prop:LH-spherical}
Let $v$ be an odd finite place, $i\in\{0,1\}$, and let $\pi_v$ be an
irreducible unitary representation of $G(k_v)$ with a nonzero
$K_{i,v}$-fixed vector. Let ${\cal C}_v$ be the space of pairs
$(f_v,\overline{f}_v)$ with $f_v\in{\cal H}(G(k_v),K_{i,v})$ and
$\overline{f}_v\in{\cal H}(\overline{G(k_v)},K_{i,v})_{gen}$
corresponding under the Satake isomorphism (Definition~\ref{def:satake-corr};
for $i=1$, transported from $i=0$ by conjugation by $w_\varpi$ in the
2-fold cover of $GL(2,k_v)$ of \S\ref{sec:local}). Then properties
(LH0)--(LH3) of Hypothesis~\ref{hyp:LH} hold for $\pi_v$ with the class
${\cal C}_v$, with (LH1) and (LH2) valid for all irreducible admissible
representations.
\end{proposition}

\begin{proof} (LH0) holds with the single index $L=K_{i,v}$. If $\sigma$ is an
irreducible admissible representation of $G(k_v)$ whose character does
not vanish on ${\cal H}(G(k_v),K_{i,v})$, then $\sigma$ has a
$K_{i,v}$-fixed line, on which ${\cal H}(G(k_v),K_{i,v})$ acts by an
algebra character $\omega_\sigma$, and $\tr \sigma(f)=\omega_\sigma(f)$
for $f\in{\cal H}(G(k_v),K_{i,v})$. Since $\sigma$ is determined by
$\omega_\sigma$, inequivalent such $\sigma$ give distinct characters,
which are linearly independent (Artin). This proves (LH1). Similarly,
if $\overline{\tau}$ is a genuine irreducible representation whose
character does not vanish on the genuine spherical Hecke algebra, then
$\tr \overline{\tau}(\overline{f})=\omega_{\overline{\tau}}(\overline{f})$
for an algebra character $\omega_{\overline{\tau}}$. Composing with the
Satake isomorphism of \S\ref{sec:spherical} gives an algebra character
of ${\cal H}(G(k_v),K_{i,v})$, which is $\omega_\sigma$ for a unique
spherical irreducible representation $\sigma=t_v(\overline{\tau})$.
Since $\overline{\tau}$ is determined by $\omega_{\overline{\tau}}$,
$t_v$ is injective. This proves (LH2). Finally, (LH3) holds with
$f_v$ the unit element, since $\tr \pi_v(f_v)=1$. Condition (LH4) concerns
archimedean places only. For $i=1$, conjugation by a lift of $w_\varpi$
carries the splitting of $\overline{K_{0,v}}$ to the splitting of
$\overline{K_{1,v}}$, because the splittings are unique.
\end{proof}

\begin{example}
Let $v$ be odd and let $\chi$ be the unramified quadratic character of
$k_v^\times$. Then $I_A(\chi)$ is unitary and reducible, with two
constituents. One has a $K_{0,v}$-fixed vector and the other a
$K_{1,v}$-fixed vector, so Proposition~\ref{prop:LH-spherical} applies to
both (with $i=0$ and $i=1$). For the running example, this is the case at
a prime with $\tau(p)=0$.
\end{example}

\section{What the local chapters provide toward Hypothesis LH}
\label{sec:LH-status}

We record what Chapters~\ref{ch:reps} and~\ref{ch:local} and
Appendix~\ref{ch:sc} contribute to Hypothesis~\ref{hyp:LH} at the places of
$S(\pi)$, and what is missing. In view of
Corollary~\ref{cor:inconsistent}, the interest of this list is local.

\begin{enumerate}
\item {\em Odd places, $\pi_v$ hyperspecial-spherical.} Hypothesis LH
holds (Proposition~\ref{prop:LH-spherical}). This covers the two
components of a reducible unramified unitary principal series.

\item {\em Odd places, $\pi_v$ a ramified principal series.}
Proposition~\ref{prop:hyp-ps}, with the correspondence of
characters of \S\ref{sec:charsA}, gives the character identity
$\tr I_A(\overline{\chi})(\overline{f})=\tr I_A(\chi)(f)$ for
functions corresponding oddly on the hyperbolic set. For (LH2) one would
also need to control the traces of
the genuine discrete series representations on that class, and for
(LH1) a separation statement. Neither is established here.
When $\chi=\chi^w$ is ramified, $I_A(\chi)$ is reducible and its two
components are not separated by such functions.

\item {\em Odd places, $\pi_v$ supercuspidal.} Appendix~\ref{ch:sc} and
\S\ref{sec:elliptic} propose a correspondence through functions
corresponding oddly on the elliptic set (Lemma~\ref{lem:ell-odd}, which
requires the character relation on every anisotropic Cartan subgroup).
The character relation proposed in Appendix~\ref{ch:sc} is false when
$q\equiv3\pmod4$ (Remark~\ref{rem:sc-caveat}).

\item {\em Odd places, $\pi_v$ special.} See
Remark~\ref{rem:savin-caveat}.

\item {\em Archimedean places.} \S\ref{sec:realA} and
\S\ref{sec:realds} give character identities for principal series
(through the hyperbolic matching) and for discrete series and limits of
discrete series (through the three elliptic matchings of
Lemma~\ref{lem:realds}, for functions with vanishing hyperbolic orbital
integrals). Condition (LH4) holds for these matchings (see the remark
following Hypothesis~\ref{hyp:LH}), and for the hyperbolic matching,
since corresponding principal series have the same infinitesimal
character. For the elliptic classes, Remark~\ref{rem:LH-families} shows
that (LH1) fails for the $\pm$odd classes when both limits of discrete
series are present, and that (LH3) constrains the $K$-types. A
verification of (LH1)--(LH3) for the even class is not
carried out here. Complex places are not treated.

\item {\em Dyadic places.} Nothing is established. See
\S\ref{sec:dyadic}.
\end{enumerate}

The trivial representation never occurs as a local component of a
cuspidal automorphic representation of $G({\mathbb{A}})$ (a consequence of
the strong approximation theorem), and the finite-dimensional
representations of dimension greater than one are not unitary. So these
do not lie in $\mathfrak{F}_v$.

\begin{example}
For the running example, $\pi_\infty$ is one of the discrete series
$\pi_{\pm11}$. By Lemma~\ref{lem:realds}(a) with $k=13$, the even elliptic
matching transfers $\overline\pi{}^{\pm}_{13/2}$ to $\pi_{\pm11}$, and (LH4)
holds. By Remark~\ref{rem:LH-families}, detecting $\pi_{11}$ requires
test functions with a $K$-type of $\pi_{-11}$ or of $F_{11}$, in addition
to $K$-types of $\pi_{11}$. Nothing is known at the place $2$.
\end{example}

\section{Summary of results and open questions}\label{sec:summary}

This section collects what the text proves, what it assumes, and what
remains open.

\emph{Proved here.}
\begin{enumerate}
\item Theorem~\ref{thm:nonexist}: no genuine cuspidal automorphic
representation of the cover corresponds to a given cuspidal $\pi$ under
the Satake correspondence of \S\ref{sec:spherical} at almost all places.
It uses \cite{LL}, \cite{GJ78}, \cite{JS} and \cite{Wa1}, \cite{Wa2} (see
Remark~\ref{rem26:nonexist}).
\item Proposition~\ref{thm:main} and Corollary~\ref{cor:inconsistent}:
Hypotheses~\ref{hyp:LH} and~\ref{hyp:GH} would imply $m(\pi)=1$ and the
existence of such a correspondence, so they cannot both hold.
\item Lemma~\ref{lem:A00}: a genuine cuspidal representation corresponding
to $\pi$ at almost all places would lie in ${\cal A}_{00}$.
\item Proposition~\ref{prop:LH-spherical}: the local hypothesis holds at
odd places where $\pi_v$ is spherical for a hyperspecial maximal compact
subgroup.
\item The local results of Chapter~\ref{ch:local}: the structure of the
genuine characters of the split torus and the matching of functions on it
over p-adic fields with $p>2$ and over ${\mathbb R}$
(Lemmas~\ref{lem:Acocycle}, \ref{lem:genuine-chars}, \ref{lem:Wequiv},
\ref{lem:FT-match}, \ref{lem:realA-match}). The resulting identity for
principal series (Proposition~\ref{prop:hyp-ps}). The elliptic identity of
Lemma~\ref{lem:ell-odd}, given a character relation on every anisotropic
Cartan subgroup, and the three real elliptic matchings
(Lemma~\ref{lem:realds}), with their consequence for non-tempered
representations.
\item Proposition~\ref{prop:wrong-orbit}: on the split torus the direct
comparison is the lifting of characters between the cover and $PGL(2)$,
with the tori of $PGL(2)$ and $SL(2)$ identified by
$\mathrm{diag}(s,1)\mapsto a(s)$.
\end{enumerate}

\emph{Quoted from the companion text \cite{JC}}
(Chapter~\ref{ch:cover}). The factor of the odd matching is the function
$\Delta_\psi$ that defines stable conjugacy on the cover
(Lemma~\ref{lem:Delta}), and it satisfies a product formula on rational
elements. At the elliptic elements $\gamma$ of $K_0$ with $2+\tr\gamma$
of odd valuation, the unit elements of the two Hecke algebras have their
stable and $\kappa$-orbital integrals interchanged
(Proposition~\ref{prop:unit-swap}). So the direct comparison has no
fundamental lemma on the elliptic set (Corollary~\ref{cor:no-fl}). The
orbits of $GL(2,k)$ on the continuous and residual spectrum of the cover
are the fibers of the squaring correspondence
(\S\ref{sec:cover-spectral}).

\emph{Assumed or unverified.} The character formula of
\S\ref{sec:weil} rests on a Conjecture. The statements labeled
Remark~2026.$n$ are unverified. The inputs of
Theorem~\ref{thm:nonexist} are quoted from memory
(Remark~\ref{rem26:nonexist}).

\emph{Open questions.}
\begin{enumerate}
\item Is there a comparison of the trace formulas of $SL(2)$ and of its
cover, built on the even correspondence at every place, and what does it
say about multiplicities (Remark~\ref{rem:even})?
\item Do the odd matchings of \S\ref{sec:charsA}--\S\ref{sec:spherical}
have a representation-theoretic meaning beyond the identity of
Proposition~\ref{prop:hyp-ps}? They do not preserve unitarity or
reducibility (Remark~\ref{rem:choices}). A partial answer is in
\S\ref{sec:cover-factor}: on characters they are multiplication by the
stable character of the Weil representation. Which representations,
other than principal series, satisfy the relation
\eqref{eq:direct-char} is not known.
\item Is there a character relation for supercuspidal representations of
the kind proposed in Appendix~\ref{ch:sc} when $q\equiv1\pmod4$
(Remark~\ref{rem:sc-caveat})? By \S\ref{sec:cover-factor} the question
should be asked with a factor that is invariant under conjugation, such
as $\Delta_\psi$, and for every $q$. By
Proposition~\ref{prop:unit-swap} the unit elements do not match with the
factor $\Delta_\psi$ at the elements close to $-1$, so the factor on the
elliptic set, if there is one, is not $\Delta_\psi$ alone.
\item Establish matchings at the dyadic places, perhaps by way of the
Iwahori-level results of \cite{GS12}, \cite{TW}.
\item Check the labels of Remark~\ref{rem26:ds-labels} and the
normalizations of Remarks~\ref{rem26:A00-norm} and~\ref{rem26:nonexist}.
\item Is there a matching on the elliptic set that exchanges stable and
$\kappa$-orbital integrals, with factor $\Delta_\psi$ times a factor
that changes sign between the two classes in a stable class, as
Proposition~\ref{prop:unit-swap} and Example~\ref{ex:epsilon} suggest?
It would pair the unstable part of the trace formula of $SL(2)$, which
\cite{LL} attach to the tori, with the stable part of the trace formula
of the cover.
\item Compare the weighted orbital integrals and the unipotent terms of
the two trace formulas (\S\ref{sec:cover-table}).
\end{enumerate}

\begin{example}
For the running example $\pi$ (from $\Delta$), the genuine partner is the
weight $13/2$ form of the Shimura correspondence. It is related to $\pi$
by the even correspondence at every place where both are unramified, and
by the even elliptic matching at the real place.
\end{example}

%% ===== end aabook6.tex =====



%% ===== begin Chapter 6 (new in Draft 7) =====
\chapter{The comparison seen from the trace formula of the cover}\label{ch:cover}
\chaptermark{The trace formula of the cover}

Chapter~\ref{ch:global} shows that the direct comparison is not a
correspondence of automorphic representations, and \S\ref{sec:where-fails}
locates the failure in the terms of the trace formula that involve the
root. That section works with a sketch of the trace formula of the cover
(Appendix~\ref{ch:tf}), quotes several facts from memory, and leaves the
elliptic terms unexamined. A companion text \cite{JC} now writes out the
trace formula of the cover in full, in its invariant form, and begins its
stabilization. This chapter reads the direct comparison through that
text. It has three aims.
\begin{itemize}
\item To say which statements of this text are proved in \cite{JC}
(\S\ref{sec:cover-supplies}).
\item To explain the direct comparison in the language of stable
conjugacy. The factor in the odd matching turns out to be the function
that defines stable conjugacy on the cover (\S\ref{sec:cover-factor}),
and the direct comparison turns out to be the known correspondence with
$PGL(2)$ in which the tori have been identified in the wrong way
(\S\ref{sec:cover-orbit}). This is a second explanation of the square,
on the geometric side, to set beside the one on the dual side in
\S\ref{sec:where-fails}.
\item To settle what happens on the elliptic set. The direct comparison
fails there as well. It fails at the elliptic elements close to $-1$,
where the unit elements of the two Hecke algebras have their stable and
unstable orbital integrals interchanged (\S\ref{sec:cover-elliptic}).
\end{itemize}
\S\ref{sec:cover-spectral} returns to the spectral side, and
\S\ref{sec:cover-table} collects the result in a table. The proofs are in
\cite[Chapter~8]{JC}. We repeat the short ones and explain the others.
Nothing in this chapter is used in Chapters~\ref{ch:local}
and~\ref{ch:global}. The status of the statements is in
Remark~\ref{rem26:cover-status}.

\paragraph{Notation.} The two texts use the same cocycle $\beta$, so
elements $(g,\zeta)$ of the local cover mean the same in both. The text
\cite{JC} writes $\phi$ and $f$ where we write $f$ and $\overline f$, and
it writes $\mu=\chi_\psi\cdot(\chi\circ p)$ where we write
$\overline\chi=\chi_\psi\mu$. We keep the notation of this text. For a
regular semisimple $\gamma\in G(F)$ and $x\in\overline{G(F)}$ over
$\gamma$ put
\[
J_f(\gamma)=|D(\gamma)|^{1/2}O_f(\gamma),\qquad
J_{\overline f}(x)=|D(\gamma)|^{1/2}O_{\overline f}(x),\qquad
D(\gamma)=(\tr\gamma)^2-4 .
\]
On the split torus these are $F_f$ and $F_{\overline f}$. We write
$\delta_a=\left(\begin{smallmatrix}a&0\\0&1\end{smallmatrix}\right)$ and
\[
D^+(\gamma)=\det(1+\gamma)=2+\tr\gamma .
\]

In this chapter $\chi_\psi(x)=\gamma_\psi(x)/\gamma_\psi(1)$, where
$\gamma_\psi(a)$ is the Weil index of the character $x\mapsto\psi(ax^2)$
of $F$, as in \cite[Lemma~3.3]{JC}. It has the properties used in
\S\ref{sec:nonexist}: it depends only on the class of $x$ modulo squares,
$\chi_\psi(xy)=(x,y)_2\chi_\psi(x)\chi_\psi(y)$, and $\chi_\psi=1$ on
${\cal O}^\times$ when $p$ is odd and $\psi$ has conductor ${\cal O}$. By
\cite[Lemma~5.7]{JC} the function $\omega$ in formula (iii) of
\S\ref{sec:weil-rep} is $\chi_{\psi^{-1}}$ in this normalization, for the
Weil representation in which $n(b)$ acts by $\psi(bx^2)$. Replacing $\psi$
by $\psi^{-1}$ multiplies $\chi_\psi$ by the quadratic character
$(-1,\cdot)_2$. This does not affect any absolute value, and it is the
normalization that Remarks~\ref{rem26:A00-norm} and~\ref{rem26:nonexist}
ask to have checked.

\section{What the companion text supplies}\label{sec:cover-supplies}

This section lists the statements of the present text that are proved in
\cite{JC}.

The text \cite{JC} is a corrected and expanded version of notes from
1995 that came from Chapter~9 of \cite{J2}. It treats the 2-fold cover of
$SL(2)$ over a number field. It proves the trace formula in its
non-invariant form (coarse and fine geometric expansions, spectral side),
then in its invariant form, and then determines how $GL(2)$ acts on the
cover and on each term. It is the worked-out form of what
Appendix~\ref{ch:tf} sketches, and the terms (a)--(e) of that appendix
are Theorems~5.16 and~6.11 of \cite{JC}.

\begin{enumerate}
\item \emph{Weil's product formula, and the automorphic character
$\chi_\psi$ of the cover of the torus} (used in \S\ref{sec:where-fails}):
\cite[Lemma~3.3]{JC}. Weil's theorem itself is quoted there from
\cite{We3}.
\item \emph{The torus is globally consistent} (\S\ref{sec:where-fails}):
\cite[Theorem~3.4 and Lemma~3.5]{JC}. The trace formula of the cover of a
torus is Poisson summation, and multiplying by $\chi_\psi$ identifies it
with Poisson summation on the torus. This is the same argument as in
\S\ref{sec:where-fails}.
\item \emph{The scalar of the intertwining operator},
$\overline c(s,\mu)=L(2s,\mu^2)/L(1+2s,\mu^2)$:
\cite[Lemma~5.6]{JC}. The unramified computation is done in the cover,
with Kubota's cocycle.
\item \emph{The residual spectrum of the cover is the sum of the even
theta representations, one for each square class}:
\cite[Theorem~5.8]{JC}. That the residues exist, are square integrable
and are indexed by the square classes is proved there. Their
irreducibility is quoted from \cite{Ge2}, \cite{GePS}, \cite{Wa1}.
\item \emph{The Jacquet module of the even Weil representation}
(fact~(ii) in the proof of Lemma~\ref{lem:A00}): \cite[Lemma~5.7]{JC}
proves, from the cocycle, that the torus acts in the Weil representation
by $\zeta\,\chi_{\psi^{-1}}(a)|a|^{1/2}\phi(ax)$, so that $r^+$ maps into
the principal series with exponent $-\frac12$.
\item \emph{The kernel and the character of the Weil representation}
(\S\ref{sec:weil}, where the character is obtained by a formal
computation and stated as a conjecture). The kernel, with its constant,
is \cite[Lemma~7.4]{JC}. The constant is checked there by composing two
kernels and recovering $\beta$. The character of $r_+-r_-$ on the
regular set is \cite[Proposition~7.5]{JC}:
\[
\Theta_+(g,\zeta)-\Theta_-(g,\zeta)=\zeta\,(c,D^+(g))_2\,\chi_\psi(D^+(g))\,
|D^+(g)|^{-1/2}\qquad(c\neq0).
\]
The character of $r_++r_-$ is the same formula with $D^-(g)=\tr g-2$ in
place of $D^+(g)$, \cite[(7.7)]{JC}. The trick of \S\ref{sec:weil-rep},
which passes from one to the other by translating by $-I$, is the
operation $\overline f\mapsto\overline f_-$ of
\S\ref{sec:cover-elliptic}. The proof in \cite{JC} that the formal
diagonal integral is the character uses the regularity theorem for the
characters of the Weil representation, which is quoted there.
\end{enumerate}

In each case the statement in the present text keeps the status it has
here. The list says where a proof can be found. It does not replace the
checks asked for in Remarks~\ref{rem26:A00-norm},
\ref{rem26:nonexist} and~\ref{rem26:where-fails}, because \cite{JC} has
not been refereed by a specialist either, and because the normalizations
of the two texts were compared only as far as stated in the paragraph on
notation above.

\begin{example}
Let $k={\mathbb Q}$. The example at the end of \S\ref{sec:where-fails}
uses $\xi(2s)/\xi(1+2s)$ and its pole at $s=1/2$. In \cite[Example~5.B]{JC}
the same function appears with the Euler factor at $2$ removed,
$(1-2^{-2s})\zeta(2s)/((1-2^{-1-2s})\zeta(2s+1))$, and its residue at
$s=1/2$ is computed to be $2/\pi^2$. The residual representation is
identified there with the even Weil representation that contains
$\theta(z)=\sum_ne^{2\pi in^2z}$.
\end{example}

\section{The factor in the odd matching}\label{sec:cover-factor}

This section identifies the factor in the odd matching on the split
torus with a function that is defined on all regular elements of the
cover, and explains what that function is for.

\subsection*{Stable conjugacy on the cover}

Two semisimple elements of $SL(2,F)$ are \emph{stably conjugate} if they
are conjugate under $GL(2,F)$. A stable conjugacy class of regular
elliptic elements is the union of two conjugacy classes, those of
$\gamma$ and of $\delta_a\gamma\delta_a^{-1}$, where $a$ is not a norm
from $E=F[\gamma]$. Labesse and Langlands \cite{LL} organize the trace
formula of $SL(2)$ by stable classes. The sum of the orbital integrals
over the two classes is the \emph{stable} orbital integral, and their
difference is the \emph{$\kappa$-orbital integral}.

The group $GL(2,F)$ also acts on the cover, by automorphisms
$\mathrm{Ad}(h)$ that lift conjugation by $h$ (through the cover of
$GL_2(F)$ of \S\ref{sec:local}). One might define stable conjugacy on the
cover in the same way. This does not work. For example
\[
\mathrm{Ad}(\delta_a)(a(t),\zeta)=(a(t),(a,t)_2\,\zeta),
\]
so when $t$ is not a square, $\mathrm{Ad}(\delta_a)$ carries a lift of
$a(t)$ to the \emph{other} lift, for suitable $a$. A genuine function
takes opposite values on the two lifts, so no genuine function other
than $0$ could be constant on such classes. In general, if $h\in GL(2,F)$
commutes with $\gamma$, then $\mathrm{Ad}(h)$ multiplies the lifts of
$\gamma$ by the sign $(\det h,D^+(\gamma))_2$ \cite[Lemma~7.2]{JC}. The
remedy, due to Adams \cite{Ad2} and Li, is to twist: put
\[
h\star x=(\det h,D^+(p(x)))_2\cdot\mathrm{Ad}(h)x .
\]
This is an action of $GL(2,F)$ on the elements of the cover with
$\tr p(x)\neq-2$. Its orbits are the \emph{stable conjugacy classes} of
the cover. They map bijectively to the stable classes of $SL(2,F)$, and
they never contain both lifts of an element.

\begin{lemma}[{\cite[Lemma~7.3]{JC}}]\label{lem:Delta}
There is exactly one function $\Delta_\psi$ on the regular set of
$\overline{G(F)}$, with values in the fourth roots of unity, that is
invariant under conjugation by $\overline{G(F)}$ and is given by
\[
\Delta_\psi(g,\zeta)=\zeta\,(c,D^+(g))_2\;\chi_\psi(D^+(g))
\]
when the lower left entry $c$ of $g$ is not $0$. It is genuine, it is
invariant under the action $\star$ of $GL(2,F)$, and on the split torus
\[
\Delta_\psi(a(x),\zeta)=\zeta\,\chi_\psi(x)^{-1}.
\]
Two regular elements of the cover are stably conjugate if and only if
their images in $SL(2,F)$ are stably conjugate and $\Delta_\psi$ takes
the same value on them. Up to the positive factor $|D^+|^{1/2}$,
$\Delta_\psi$ is the character of $r_+-r_-$.
\end{lemma}

So $\Delta_\psi$ is the function that tells the two lifts of an element
apart in a way that is compatible with stable conjugacy. In the theory of
Adams and Li it is the transfer factor between the cover and $SO(3)$.

\subsection*{The odd matching}

In \S\ref{sec:where-fails} the globally consistent form of the odd
matching on the torus is
$\overline\phi(a(x),1)=\chi_\psi(x)^{-1}\phi(a(x))$. By
Lemma~\ref{lem:Delta} this is
\begin{equation}\label{eq:odd-Delta}
\overline\phi(\tilde a)=\Delta_\psi(\tilde a)\,\phi(a)\qquad
\text{for }\tilde a\in\overline{A(F)}\text{ over }a .
\end{equation}
The factor of the odd matching is the transfer factor of Adams and Li,
restricted to the split torus. This has two consequences.

First, it explains the choices. Definition~\ref{def:corrA} contains a
constant $\eta_\varpi$ with $\eta_\varpi^2=(\varpi,\varpi)_2$, and
Lemma~\ref{lem:realA-match} contains the factor $-i$. By
Remark~\ref{rem:choices} there are exactly two bijections
$\chi\mapsto\overline\chi$, and neither is preferred. In
\eqref{eq:odd-Delta} the choice is the choice of the additive character
$\psi$. Replacing $\psi$ by $\psi_a(x)=\psi(ax)$ multiplies
$\Delta_\psi(x)$ by $(a,D^+(p(x)))_2$ \cite[Lemma~7.3(d)]{JC}. For $a$ a
unit that is not a square this exchanges the two bijections. The choices
at different places are not independent: globally one needs one
character $\psi$ of $k\backslash{\mathbb A}$.

Second, it shows how to extend the odd matching beyond the torus. The
section $\gamma\mapsto(\gamma,1)$ is not compatible with conjugation, so
a condition such as $\Phi_f(t)=\Phi_{\overline f}(t,1)$ in
Definition~\ref{def:ell-corr}(a), or a character relation such as
\eqref{eq:ell-char-rel}, depends on the coordinates. This is the
reason for the failure recorded in Remark~\ref{rem:sc-caveat}: for
$q\equiv3\pmod4$ the relation is not invariant under conjugation. The
invariant form of the same idea is the following.

\begin{definition}\label{def:direct}
Functions $f\in C_c^\infty(G(F))$ and
$\overline f\in C_c^\infty(\overline{G(F)})_{gen}$ \textbf{match directly}
on a set $U$ of regular elements of $G(F)$, invariant under conjugation,
if
\[
J_{\overline f}(x)=\Delta_\psi(x)\,J_f(p(x))\qquad\text{for all }x\text{ with }
p(x)\in U .
\]
\end{definition}

On the hyperbolic set this is the odd matching, by \eqref{eq:odd-Delta}.
The corresponding relation between characters is
\begin{equation}\label{eq:direct-char}
\Theta_{\overline\pi}(x)=\overline{\Delta_\psi(x)}\,\Theta_\pi(p(x)),
\end{equation}
because then $\Theta_{\overline\pi}(x)J_{\overline f}(x)=\Theta_\pi(p(x))
J_f(p(x))$ and the Weyl integration formulas of \S\ref{sec:elliptic}
give $\tr\overline\pi(\overline f)=\tr\pi(f)$. For principal series
\eqref{eq:direct-char} holds, with $\overline\chi=\chi_\psi\mu$ and
$\chi=\mu$, and it gives Proposition~\ref{prop:hyp-ps}. By the last
statement of Lemma~\ref{lem:Delta}, the relation \eqref{eq:direct-char}
says: \emph{the character of $\overline\pi$ is the character of $\pi$
times the character of $r_+-r_-$ for $\psi^{-1}$, times
$|\det(1+\gamma)|^{1/2}$.} This is a partial answer to the second open
question of \S\ref{sec:summary}. The odd matchings do have a meaning for
representations: they multiply characters by the stable character of the
Weil representation.

The factor $\Delta_\psi$ is $1$ at $(\gamma,1)$ in the cases where the
relation \eqref{eq:ell-char-rel} looks right. For $F={\mathbb R}$ this
holds for the rotations $r(\theta)$, $0<\theta<\pi$, in the
normalization of \cite[Example~7.B]{JC}. For $p$ odd it holds for
$\gamma\in K_0$ with $c$ and $D^+(\gamma)$ units. It is not $1$ at the
element $h$ of Remark~\ref{rem:sc-caveat}. That element has
$\tr h=-2+\varpi e$, so $D^+(h)=\varpi e$ has valuation $1$, and
$\Delta_\psi(h,\zeta)^2=(\varpi,-1)_2=(-1)^{(q-1)/2}$. For $q\equiv3
\pmod4$ the value is $\pm i$. We return to these elements in
\S\ref{sec:cover-elliptic}.

\begin{example}\label{ex:eta-psi}
Let $k={\mathbb Q}$ and let $\psi$ be the character of
${\mathbb Q}\backslash{\mathbb A}$ with $\psi_\infty(x)=e^{-2\pi ix}$ and
$\psi_p(x)=e^{2\pi i\{x\}_p}$. By \cite[Example~3.B]{JC},
$\chi_{\psi_\infty}(x)$ is $1$ for $x>0$ and $i$ for $x<0$, and for $p$
odd $\chi_{\psi_p}(u)=1$ and
$\chi_{\psi_p}(pu)=\left(\frac up\right)\epsilon_p$ for a unit $u$, with
$\epsilon_p=1$ if $p\equiv1\pmod4$ and $\epsilon_p=i$ if
$p\equiv3\pmod4$.

At the real place \eqref{eq:odd-Delta} gives the factor $1$ for $x>0$ and
$-i$ for $x<0$. This is Lemma~\ref{lem:realA-match}. So the real matching
of \S\ref{sec:realA} belongs to $\psi_\infty(x)=e^{-2\pi ix}$.

At $p=3$, with $\varpi=3$, \eqref{eq:odd-Delta} gives the factor $1$ on
$A_2(F)$ and $\chi_{\psi_3}(3u)^{-1}=-i\left(\frac u3\right)$ at
$a(3u)$. Definition~\ref{def:corrA} gives
$\eta_3^{-1}(3u,3)_2=-\eta_3^{-1}\left(\frac u3\right)$, because
$(3,3)_2=-1$. They agree for $\eta_3=-i$. The choice $\eta_3=i$ in the example after
Remark~\ref{rem:choices} belongs to the additive character
$\psi_3^{-1}$. At $p=5$ the same computation gives $\eta_5=1$ for
$\psi_5$, and $\eta_5=-1$ for $x\mapsto\psi_5(ux)$ with $u$ a unit that
is not a square.

So for this $\psi$ the constants are $-i$ at the real place, $-i$ at
$3$ and $1$ at $5$. They are linked. Another character of
${\mathbb Q}\backslash{\mathbb A}$ is $x\mapsto\psi(bx)$ with
$b\in{\mathbb Q}^\times$, and it changes the local factors by the Hilbert
symbols $(b,\cdot)_{2,v}$ at all places at once. Their product is $1$ on
${\mathbb Q}^\times$. A choice of the constants place by place, without
this constraint, would not be consistent globally.
\end{example}

\section{The right factor and the wrong orbit}\label{sec:cover-orbit}

This section compares the direct comparison with the character identity
between the cover and $PGL(2)$ that is known to hold, and finds that the
two differ only in how the tori are identified.

For the pair formed by the cover and $SO(3)\cong PGL(2)$ there is a
lifting of characters, due to Adams \cite{Ad2} over ${\mathbb R}$ and to
Schultz over p-adic fields (see \cite[\S7.7]{JC} for references and for
the status of the statement). For a principal series it reads
\begin{equation}\label{eq:adams}
\Theta_{\overline\pi}(x)=\bigl(\Theta_+(x)-\Theta_-(x)\bigr)\cdot
\Theta_\sigma\bigl(t(x)\bigr),
\end{equation}
where $\sigma$ is a representation of $PGL(2,F)$, $\Theta_\pm$ are the
characters of the two halves of a Weil representation, and $t(x)$ is the
conjugacy class in $PGL(2,F)$ with eigenvalue ratio $\epsilon$ when
$p(x)$ has eigenvalues $\epsilon^{\pm1}$. For $x$ over $a(s)$, $t(x)$ is
the class of $\mathrm{diag}(s,1)$.

Compare this with \eqref{eq:direct-char}. The first factor is the same:
up to $|D^+|^{1/2}$ it is $\Delta$. The second factor is a character of
$PGL(2,F)$ at $\mathrm{diag}(s,1)$ in one case and a character of
$SL(2,F)$ at $a(s)$ in the other.

\begin{proposition}[{\cite[Proposition~8.3]{JC}}]\label{prop:wrong-orbit}
For a character $\mu$ of $F^\times$ let $\sigma(\mu)$ be the principal
series of $PGL(2,F)$ induced from
$\mathrm{diag}(t_1,t_2)\mapsto\mu(t_1/t_2)$. Let $s\in F^\times$,
$s\neq\pm1$.

(a) $D(a(s))=D^+(a(s))\cdot D'(s)$, where $D'(s)=(1-s)^2/s$ is the Weyl
discriminant of $\mathrm{diag}(s,1)$ in $PGL(2,F)$.

(b) $|D^+(a(s))|^{-1/2}\,\Theta_{\sigma(\mu)}(\mathrm{diag}(s,1))
=\Theta_{I_A(\mu)}(a(s))$, where $I_A(\mu)$ is the principal series of
$SL(2,F)$ induced from $a(x)\mapsto\mu(x)$.

(c) The restriction of $\sigma(\mu)$ to $SL(2,F)$ is $I_A(\mu^2)$.
\end{proposition}

\begin{proof}
(a) $D(a(s))=(s-s^{-1})^2=(s-1)^2(s+1)^2/s^2$ and
$D^+(a(s))=2+s+s^{-1}=(1+s)^2/s$. (b) Both characters are
$\mu(s)+\mu(s)^{-1}$ divided by the square root of the absolute value of
the Weyl discriminant, which is $D'(s)$ on the left and $D(a(s))$ on the
right. Apply (a). (c) The torus of $SL(2)$ maps to that of $PGL(2)$ by
$a(x)\mapsto\mathrm{diag}(x,x^{-1})\equiv\mathrm{diag}(x^2,1)$, where the
inducing character is $\mu(x^2)$, and $PGL(2,F)$ is the product of its
Borel subgroup and the image of $SL(2,F)$.
\end{proof}

By (b), on the split torus the right side of \eqref{eq:adams} for
$\sigma=\sigma(\mu)$ is the right side of \eqref{eq:direct-char} for
$\pi=I_A(\mu)$. So for principal series the direct comparison \emph{is}
the lifting \eqref{eq:adams}, read as follows: the element
$\mathrm{diag}(s,1)$ of $PGL(2)$ is replaced by the element $a(s)$ of
$SL(2)$, and the representation $\sigma(\mu)$ of $PGL(2)$ by the
representation $I_A(\mu)$ of $SL(2)$.

There are two ways to go from $PGL(2)$ to $SL(2)$, and they give
different answers.
\begin{itemize}
\item \emph{Restriction.} This is what happens to automorphic forms. By
(c), $\sigma(\mu)$ restricts to $I_A(\mu^2)$. An element $a(s)$ of
$SL(2)$ becomes $\mathrm{diag}(s^2,1)$ in $PGL(2)$, which
\eqref{eq:adams} compares with a lift of $a(s^2)$. This is the squaring
correspondence $t=t'^{\,2}$, and the even matching that compares
$r(\theta)$ with the element over $r(2\theta)$.
\item \emph{Identifying the tori.} The map
$\mathrm{diag}(s,1)\mapsto a(s)$ is an isomorphism of the diagonal tori
that commutes with the Weyl groups. It sends $\sigma(\mu)$ to $I_A(\mu)$
and compares $a(s)$ with a lift of $a(s)$. This is the direct
comparison, $t=\pm t'$. The map is not induced by a homomorphism between
$PGL(2)$ and $SL(2)$ in either direction.
\end{itemize}
The direct comparison therefore uses the right transfer factor together
with the wrong correspondence of orbits. This is the geometric form of
the diagnosis of \S\ref{sec:where-fails}. There the direct comparison is
an identification of the dual tori that does not extend to the dual
groups. Here it is an identification of the tori that does not extend to
the groups.

The squares have a special position on the cover, which explains the
shape of Definition~\ref{def:ell-corr}(b). If $\gamma=\delta^2$ with
$\delta$ in the same torus, then
$1+\gamma=\delta(\delta+\delta^{-1})=\delta\cdot\tr\delta$, so
\[
D^+(\delta^2)=(\tr\delta)^2 .
\]
Hence the sign $(\det h,D^+)_2$ is trivial at the elements over squares:
there $GL(2,F)$ does not exchange the two lifts, and the element
$\tau(\delta)=(\delta,1)^2$ of Definition~\ref{def:ell-corr}(b) is a
canonical lift of $\delta^2$. Away from the classes with $D^+$ a square
or $\theta$ times a square, the average of an orbital integral of a
genuine function over the conjugates under $GL(2,F)$ is zero
\cite[Lemma~7.6(b)]{JC}. A character of $PGL(2,F)$, restricted to
$SL(2,F)$, is invariant under $GL(2,F)$. So a comparison of the cover
with $SL(2)$ that passes through $PGL(2)$ can only see the cover near
the elements over squares. This is why an even matching requires
$\Phi_{\overline f}$ to vanish outside $\tau(T(F))$.

\begin{example}
For the running example, $\Delta$ has parameters
$\tau(p)=p^{11/2}(\alpha_p+\alpha_p^{-1})$. The representation of
$PGL(2,{\mathbb Q}_p)$ is $\sigma(\mu)$ with $\mu(p)=\alpha_p$, and
\eqref{eq:adams} pairs it with the genuine principal series
$I_A(\chi_\psi\mu)$, with $t'=\alpha_p$. Restriction gives $I_A(\mu^2)$,
with $t=\alpha_p^2$. Identifying the tori gives $I_A(\mu)$, with
$t=\alpha_p$, which is not a local component of any cuspidal
representation attached to $\Delta$.
\end{example}

\section{The elliptic terms}\label{sec:cover-elliptic}

This section answers the question left open at the end of
\S\ref{sec:where-fails}: what the direct comparison does on the elliptic
set.

\subsection*{What is consistent}

The function $\Delta_\psi$ has a product formula. If $\gamma\in G(k)$ is
regular and $x_v$ are the local components of its image in
$\overline{G({\mathbb A})}$, then
\begin{equation}\label{eq:Delta-product}
\prod_v\Delta_{\psi_v}(x_v)=1
\end{equation}
\cite[Lemma~8.4]{JC}. In the restricted-product coordinates of
\S\ref{sec:global} one may take $x_v=(\gamma,1)$, and then the product is
$\prod_v(c,D^+)_{2,v}\cdot\prod_v\chi_{\psi_v}(D^+)$, which is $1$ by
Hilbert reciprocity and Weil's product formula. For $\gamma$ in the
torus this is the product formula $\prod_v\chi_{\psi_v}(x)=1$ of
\S\ref{sec:where-fails}. So the consistency found there for the torus is
a special case: if $f=\otimes f_v$ and $\overline f=\otimes\overline f_v$
match directly at \emph{every} place, on the whole regular set, then
\[
\prod_vJ_{\overline f_v}(x_v)=\prod_vJ_{f_v}(\gamma)
\]
for every regular $\gamma\in G(k)$, elliptic or hyperbolic, and the
regular elliptic terms of the two trace formulas agree
\cite[Corollary~8.5]{JC}.

\subsection*{What fails}

In a comparison of trace formulas the functions at almost all places are
not at our disposal. They are the unit elements $1_K$ and $e_K$ of the
spherical Hecke algebras, which correspond under the Satake
correspondence of \S\ref{sec:spherical}. On the hyperbolic set the unit
elements match directly. On the elliptic set they do not.

\begin{proposition}[{\cite[Proposition~8.6]{JC}}]\label{prop:unit-swap}
Let $p$ be odd, $\psi$ of conductor ${\cal O}$, $K=K_0$, and let
$\gamma\in K$ be regular, with lift $\gamma^*=\kappa(\gamma)$ in the
image of the splitting. Let $u$ be a unit that is not a square.

(a) If $D^+(\gamma)$ is a unit, then $\Delta_\psi(\gamma^*)=1$.

(b) Suppose that $D^+(\gamma)=2+\tr\gamma$ has odd valuation. Then
$\gamma$ is elliptic, $F[\gamma]$ is ramified, and the stable class of
$\gamma$ consists of the classes of $\gamma$ and of
$\gamma'=\delta_u\gamma\delta_u^{-1}$. Both lie in $K$, and
\[
J_{1_K}(\gamma')=J_{1_K}(\gamma),\qquad
J_{e_K}(\delta_u\star\gamma^*)=-J_{e_K}(\gamma^*).
\]
So for $SL(2,F)$ the stable orbital integral of the unit element at
$\gamma$ is $2J_{1_K}(\gamma)\neq0$ and its $\kappa$-orbital integral is
$0$. For the cover the stable orbital integral of the unit element at
$\gamma^*$ is $0$ and its $\kappa$-orbital integral is
$2J_{e_K}(\gamma^*)$.

(c) If $D(\gamma)$ has valuation $1$, then
$J_{e_K}(\gamma^*)=J_{1_K}(\gamma)>0$.
\end{proposition}

The proof of (b) is short, and it shows what is going on. The
automorphism $\mathrm{Ad}(\delta_u)$ preserves the image of the splitting
of $K$, because the splitting is unique. So
$\mathrm{Ad}(\delta_u)\gamma^*=\gamma'^*$, and $e_K$ takes the same value
on the conjugates of $\gamma^*$ and of $\gamma'^*$. But the element that
is \emph{stably} conjugate to $\gamma^*$ is
$\delta_u\star\gamma^*=(u,D^+(\gamma))_2\,\gamma'^*=-\gamma'^*$, because $u$
is a non-square unit and $D^+(\gamma)$ has odd valuation. The twist that
makes stable conjugacy well defined moves the element out of the image
of the splitting, and $e_K$ changes sign.

So at these elements the unit elements have their stable and
$\kappa$-orbital integrals \emph{interchanged}. Any matching that
respects stable conjugacy, as the matching with factor $\Delta_\psi$
does, would have to send the stable orbital integral $2J_{1_K}(\gamma)$
of $1_K$ to the stable orbital integral $0$ of $e_K$. Hence:

\begin{corollary}[{\cite[Corollary~8.7]{JC}}]\label{cor:no-fl}
Let $p$ be odd and let $\gamma\in K_0$ be regular with $D^+(\gamma)$ of
odd valuation. The unit elements $1_K$ and $e_K$ do not match directly
on any invariant set that contains $\gamma$. The same holds with
$\Delta_\psi$ replaced by any nowhere vanishing genuine function that is
constant on stable conjugacy classes.
\end{corollary}

The elements in question satisfy $\gamma\equiv-1$ modulo the prime of
$F[\gamma]$. The element $h$ of Remark~\ref{rem:sc-caveat} is
one of them, with $D^+(h)=\varpi e$. So the counterexample to the
supercuspidal character relation in Appendix~\ref{ch:sc} and the failure
of the direct comparison on the elliptic set are the same phenomenon,
and it is not special to $q\equiv3\pmod4$. Near $+1$ nothing of the
kind happens: the cover is a product near the identity, and
$\Delta_\psi=1$ there on $(\gamma,1)$.

Why $-1$? On the cover the $\kappa$-orbital integral of
$\overline f$ at $x$ is, up to the sign $(-1,-1)_2$, the stable orbital
integral at the element $(-I,1)x$ over $-\gamma$ of the translate
$\overline f_-(y)=\overline f((-I,1)y)$ \cite[Lemma~7.6(c)]{JC}. For
$SL(2)$ there is no such identity: the $\kappa$-orbital integrals are new
functions on the elliptic tori, and \cite{LL} attach to them the tori as
endoscopic groups. For the cover the two endoscopic groups are both
$SO(3)$, and they are exchanged by the central element $-1$. The trace
formula of the cover has the form $S(\overline f)+S(\overline f_-)$ on
its elliptic regular terms, with one stable distribution $S$
\cite[Theorem~7.7]{JC}. The elements close to $-1$ are where the two
halves trade places.

\begin{remyy}\label{rem26:li-fl}
The vanishing in Proposition~\ref{prop:unit-swap}(b) is what the
fundamental lemma for the cover and $SO(3)$ predicts. Write
$\gamma=e/\overline e$ with $e\in F[\gamma]^\times$. The class in
$PGL(2,F)$ that corresponds to $\gamma^*$ is the image of $e$, and
$D^+(\gamma)=\tr(e)^2/N(e)$ has odd valuation exactly when $\det e$ does.
Such an element is not conjugate into the center times
$GL(2,{\cal O})$, so the unit element of the Hecke algebra of
$PGL(2,F)$ has orbital integral $0$ there. The normalizations of the
transfer factors of Li and of \cite{JC} have not been compared, so this
is a consistency check only.
\end{remyy}

\subsection*{The size of the failure}

The failure can be measured on rational elements. Let $S$ contain the
archimedean and dyadic places and the places where $\psi_v$ is ramified.
Let $f=f_S\otimes\bigotimes_{v\notin S}1_{K_v}$ and
$\overline f=\overline f_S\otimes\bigotimes_{v\notin S}e_{K_v}$, with
$f_v$ and $\overline f_v$ matching directly at each $v\in S$. Let
$\gamma\in G(k)$ be regular, lying in $K_v$ and with $D(\gamma)$ of
valuation at most $1$ for all $v\notin S$. Then
\begin{equation}\label{eq:epsilon}
\prod_vJ_{\overline f_v}(x_v)=\varepsilon^S(\gamma)\prod_vJ_{f_v}(\gamma),
\qquad\varepsilon^S(\gamma)=\prod_{v\notin S}\Delta_{\psi_v}(\gamma_v^*)^{-1},
\end{equation}
where $\gamma_v^*=\kappa_v(\gamma)$ \cite[Proposition~8.8]{JC}. The
factor $\varepsilon^S(\gamma)$ is a fourth root of unity. It is $1$ when
$D^+(\gamma)$ is a unit outside $S$, by
Proposition~\ref{prop:unit-swap}(a). For a hyperbolic $\gamma$ the two
sides agree. For elliptic $\gamma$ the factor is not $1$ in general.

\begin{example}\label{ex:epsilon}
Let $k={\mathbb Q}$, $S=\{\infty,2\}$, $\psi$ as in
Example~\ref{ex:eta-psi}, and
\[
\gamma=\left(\begin{array}{cc}2&-1\\3&-1\end{array}\right)\in
SL(2,{\mathbb Z}),\qquad\tr\gamma=1,\quad D^+(\gamma)=3,\quad
D(\gamma)=-3 .
\]
It has order $6$, and ${\mathbb Q}[\gamma]={\mathbb Q}(\sqrt{-3})$. In
the coordinates $(\gamma,1)$,
\[
\Delta_{\psi_\infty}(\gamma,1)=1,\qquad\Delta_{\psi_2}(\gamma,1)=i,\qquad
\Delta_{\psi_3}(\gamma,1)=-i,
\]
and $\Delta_{\psi_p}(\gamma,1)=1$ for $p\geq5$. The product is $1$, as
\eqref{eq:Delta-product} says. At $p=3$ the splitting is
$\kappa_3(\gamma)=(\gamma,s_3(\gamma))$ with
$s_3(\gamma)=(3,-1)_2=-1$ by (\ref{eq:sv}), so
$\Delta_{\psi_3}(\gamma_3^*)=i$ and
\[
\varepsilon^S(\gamma)=-i .
\]
For test functions that match directly at $\infty$ and at $2$ and are
the unit elements elsewhere, the term of $\gamma$ in the trace formula
of the cover is $-i$ times its term in the trace formula of $SL(2)$. For
$\gamma'=\delta_{-1}\gamma\delta_{-1}^{-1}$, which is conjugate to
$\gamma$ under $GL(2,{\mathbb Q})$ and not under $SL(2,{\mathbb Q})$, the
factor is $+i$. So where the trace formula of $SL(2)$ has the sum of the
two orbital integrals of $f$, the trace formula of the cover has $-i$
times their difference.

The sign $s_3(\gamma)=-1$ is the value at $\gamma$ of the Jacobi symbol
$\left(\frac d{c}\right)$ that occurs in the multiplier of the theta
function. So the multiplier system of forms of half-integral weight is
what prevents the direct comparison on the elliptic set. The values
above are computed by the script \texttt{comparison\_n2\_check.py} of
\cite{JC}.
\end{example}

\section{The spectral side again}\label{sec:cover-spectral}

This section adds three points to the account of the spectral side in
\S\ref{sec:where-fails}. They use the invariant trace formula of the
cover, \cite[Corollary~6.14]{JC}:
\[
\sum_{\overline\pi}m(\overline\pi)\tr\overline\pi(\overline f)
-\frac14\sum_{b\in k^\times/k^{\times2}}\tr I_A(\chi_{\psi_b})(\overline f)
-\frac12\sum_{\mu}\frac1{2\pi}\int_{\mathbb R}\frac{m'}{m}(\mu,it)\,
\tr I_A(\chi_\psi\mu|\cdot|^{it})(\overline f)\,dt .
\]
Here $\overline\pi$ runs over the discrete spectrum, cuspidal and
residual, $\mu$ runs over the Hecke characters of $k$ that are trivial on
the positive reals embedded diagonally at the archimedean places,
$\psi_b(x)=\psi(bx)$, and $m(\mu,s)$ is the scalar
$\overline c(s,\mu)$ of \S\ref{sec:where-fails} completed at the bad
places.

\subsection*{The term at $s=0$}

The characters of the torus fixed by the Weyl group contribute a term
$\frac14\tr(M(w,0)I_A(\cdot))$ to the spectral side of either trace
formula. For the cover the fixed characters are the
$\chi_{\psi_b}=\chi_\psi\cdot(b,\cdot)_2$, and they correspond under the
direct comparison to the quadratic characters $(b,\cdot)_2$ of the torus
of $SL(2)$, which are the fixed characters there. The operators do not
correspond. For the cover $M(w,0)=-1$ for every $b$
\cite[Proposition~6.13]{JC}. This is the second term of the display. For
$SL(2)$ and $b\neq1$, the representation $I_A((b,\cdot)_2)$ is reducible
at every place, into $\pi_v^+\oplus\pi_v^-$, and $M(w,0)$ acts on
$\otimes_v\pi_v^{\epsilon_v}$ by $\prod_v\epsilon_v$ \cite{LL}. For
directly matching functions the two terms therefore differ by
\[
\frac14\Bigl(\tr\bigl(M(w,0)I_A((b,\cdot)_2)(f)\bigr)
+\tr I_A((b,\cdot)_2)(f)\Bigr)
=\frac12\sum_{\prod_v\epsilon_v=1}\ \prod_v\tr\pi_v^{\epsilon_v}(f_v) .
\]
In \cite{LL} this is the contribution of the trivial character of the
torus of $k(\sqrt b)$, one of the endoscopic groups of $SL(2)$. On the
cover there is nothing to match it. The sum on the right has the same
shape as the trace of $\overline f$ on the theta series of one square
class, $\sum_{\prod\epsilon_v=1}\prod_v\tr r_{\epsilon_v}(\overline f_v)$
\cite[Proposition~7.12]{JC}. On the cover the representations of this
shape are in the discrete spectrum. On $SL(2)$ they are in a term of the
continuous spectrum.

\subsection*{Where the multiplicities are}

Remark~\ref{rem:even}(3) argues heuristically that a comparison through
$PGL(2)$ controls sums over $L$-packets and not the multiplicity of one
representation. The structure of the two trace formulas gives a more
definite form of this. For $SL(2)$ the multiplicity of a representation
inside its $L$-packet is decided by the unstable part of the trace
formula, which has two sources: the $\kappa$-orbital integrals on the
elliptic tori, and the operators $M(w,0)$ at the quadratic characters.
Both lead to tori as endoscopic groups \cite{LL}. The cover has neither.
Its $\kappa$-orbital integrals are stable orbital integrals of
$\overline f_-$, its operators $M(w,0)$ are the scalar $-1$, and its
endoscopic groups are two copies of $SO(3)$. What decides multiplicities
on the cover is the pair $(\overline f,\overline f_-)$, that is, the
central character and the sign in Waldspurger's theorem. The terms of
the trace formula of $SL(2)$ that would have to be controlled to prove
multiplicity one have no image in the trace formula of the cover.

\subsection*{The fibers of the squaring map}

The map $t'\mapsto t'^{\,2}$ is two-to-one (Remark~\ref{rem:even}(1)),
and the bijection $\chi\mapsto\overline\chi$ depends on a choice between
two (Remark~\ref{rem:choices}). Both are the action of $GL(2)$ on the
cover. For $a\in F^\times$,
\[
I_A(\chi_\psi\mu)\circ\mathrm{Ad}(\delta_a)\cong
I_A(\chi_\psi\mu\cdot(a,\cdot)_2)=I_A(\chi_{\psi_a}\mu)
\]
\cite[Proposition~8.9]{JC}. For $p$ odd and $a=u$ a non-square unit,
$\mathrm{Ad}(\delta_u)$ preserves the splitting of $K_0$ and exchanges the
unramified genuine representations with parameters $t'$ and $-t'$.
Globally, $GL(2,k)$ acts on the genuine automorphic spectrum, and its
orbits on the continuous and residual spectrum are the sets with
$\mu^2$ fixed. In particular it permutes the theta representations of the
different square classes transitively. So \emph{the fibers of the
squaring correspondence are the orbits of $GL(2)$.} The direct comparison
needs a section of the squaring map, and there is none that $GL(2)$
respects. Each choice of $\psi$ gives one, and $\psi\mapsto\psi_u$
exchanges the two.

\begin{example}
Let $k={\mathbb Q}$ and $S=\{\infty,2\}$. The theta representations
unramified outside $S$ belong to the square classes of $\pm1,\pm2$
\cite[Example~5.A]{JC}. The matrices $\mathrm{diag}(b,1)$,
$b\in\{\pm1,\pm2\}$, permute them transitively. Under the squaring
correspondence all four go to the trivial representation of
$SL(2,{\mathbb A})$. For the running example, the genuine cusp form of
weight $13/2$ attached to $\Delta$ and its conjugates under
$\mathrm{diag}(b,1)$ correspond to $\Delta$ twisted by the quadratic
characters $(b,\cdot)_2$, and these twists have the same restriction to
$SL(2)$.
\end{example}

\section{The comparison term by term}\label{sec:cover-table}

This section collects the results of \S\ref{sec:where-fails} and of this
chapter in one table.

\begin{center}
\renewcommand{\arraystretch}{1.2}
\begin{tabular}{p{0.27\textwidth}p{0.33\textwidth}p{0.32\textwidth}}
\hline
Term & Direct comparison ($t=\pm t'$) & Squaring ($t=t'^{\,2}$)\\
\hline
automorphic characters of the torus & correspond, through $\chi_\psi$ &
two-to-one\\
regular hyperbolic orbital integrals & agree
(\S\ref{sec:where-fails}) & compared at $a(s)$ and over $a(s^2)$\\
traces of principal series & agree (Proposition~\ref{prop:hyp-ps}) &
agree at $(\mu,s)$ and $(\mu^2,2s)$\\
regular elliptic orbital integrals & differ by $\varepsilon^S(\gamma)$;
no fundamental lemma (Corollary~\ref{cor:no-fl}) & expected to agree,
through $PGL(2)$ (Li)\\
weights $m'/m$ of the continuous spectrum & differ: $L(s,\mu)$ against
$L(2s,\mu^2)$ & agree\\
the term at $s=0$ & differ by a term of \cite{LL}
(\S\ref{sec:cover-spectral}) & not examined\\
residual spectrum & no partner & every theta representation goes to the
trivial representation\\
cuspidal spectrum & no partner (Theorem~\ref{thm:nonexist}) &
Shimura--Waldspurger, then restriction\\
weighted orbital integrals, unipotent terms & not examined & not
examined\\
\hline
\end{tabular}
\end{center}

\medskip
The first three rows are the torus. The direct comparison is right for
them and for nothing else.

\begin{remyy}\label{rem26:cover-status}
\emph{Status of Chapter~\ref{ch:cover}.} The chapter was written by
Claude in November 2026, together with Chapter~8 of \cite{JC}, and
neither has been refereed. Lemma~\ref{lem:Delta},
Proposition~\ref{prop:unit-swap}, Corollary~\ref{cor:no-fl} and the
identities \eqref{eq:Delta-product} and \eqref{eq:epsilon} are quoted
from \cite{JC}, where they are proved. Proposition~\ref{prop:wrong-orbit}
is proved here. The numerical values in Examples~\ref{ex:eta-psi}
and~\ref{ex:epsilon} are computed by the scripts of \cite{JC}. The
identification of $\eta_3$ in Example~\ref{ex:eta-psi} uses
Definition~\ref{def:corrA} as printed and the normalization of
$\chi_\psi$ fixed at the beginning of this chapter. The statements about
the trace formula of $SL(2)$ in \S\ref{sec:cover-spectral} (the form of
the term at $s=0$, the action of $M(w,0)$, and the role of the tori) are
quoted from memory of \cite{LL} and have not been checked. The lifting
\eqref{eq:adams} is quoted. The column ``Squaring'' of the table in
\S\ref{sec:cover-table} records what is expected from the work of
Waldspurger and Li, and it is not proved in either text. References to
\cite{JC} by number are to its Draft~3.
\end{remyy}

\begin{example}
For the running example, the table reads as follows at an odd prime $p$.
The unit elements $1_{K_p}$ and $e_{K_p}$ correspond under the Satake
correspondence, their traces agree on $I_A(\mu)$ and
$I_A(\chi_{\psi_p}\mu)$, and their hyperbolic orbital integrals agree. At
the elliptic element $\gamma$ of Example~\ref{ex:epsilon} and $p=3$
their orbital integrals are equal, both being $J_{1_K}(\gamma)$, while a
direct matching would require the ratio $\Delta_{\psi_3}(\gamma^*)=i$.
\end{example}
%% ===== end Chapter 6 =====

\appendix

\chapter{On a supercuspidal correspondence}\label{ch:sc}

\section{Introduction}
In this appendix, we discuss a proposed character relation between
``oddly corresponding'' supercuspidal
representations of $G(F)=SL(2,F)$ and its 2-fold metaplectic cover
$\overline{SL(2,F)}$, where $F$ is a p-adic field with $p>2$.
Following a suggestion of P. Kutzko, we use the Frobenius formula
to obtain character relations between ``corresponding'' supercuspidal
representations of $G(F)$ consistent with the global
correspondence in Chapter~\ref{ch:global}.
The parameterization of the supercuspidals for the correspondence
which we use is that of D. Manderscheid \cite{M1}-\cite{M2}, recalled
below (see also \cite{M3}).

\begin{example}
Let $\sigma$ be a cuspidal representation of $SL(2,{\mathbb F}_q)$, inflated to
$K_0$. Then $\pi={\rm c\text{-}Ind}_{K_0}^{G(F)}\sigma$ is a depth-zero
supercuspidal representation of $G(F)$, and the 1995 draft relates it to
${\rm c\text{-}Ind}_{\overline{K_0}}^{\overline{G(F)}}(\sigma\otimes\varepsilon)$
(Remark~\ref{prop:164}).
\end{example}

\section{Filtrations and classification}\label{sec:filtrations}

This section recalls Manderscheid's description of the supercuspidal representations of $G(F)$ and of its cover by compact induction from maximal compact subgroups.


Let
$F$ denote a p-adic field, let
$\varpi =\varpi_F$ be a fixed uniformizer of $F$, let $P:=P_F:=\varpi_F{\cal{O}}_F$,
and ${\mathbb{F}} :={\cal{O}}_F/P_F$.
Let ${\cal{B}}$ denote the tree for $G(F)$, $x_0$ the vertex fixed by
$K_0:=G({\cal{O}}_F)$ and $x_1$ the vertex fixed by $
K_1:=w_\varpi^{-1}G({\cal{O}}_F)w_\varpi$, where
$w_\varpi:=
\left(\begin{array}{cc}
0&\varpi \\ 1& 0\end{array}\right)$.
We may find a metric $d$ on the ambient affine space of ${\cal{B}}$
for which the subgroup of $K_0$ fixing the ball
\begin{equation}
{\cal{B}}_m(x_0):=\{\text{ vertices }x\ |\ d(x,x_0)\leq m\}
\end{equation}
pointwise, $\operatorname{Fix}\ {\cal{B}}_m(x_0)$, is the standard filtration subgroup:
\begin{equation}
\operatorname{Fix}\ {\cal{B}}_m(x_0)=K_{0,m}:=
\{\left(\begin{array}{cc} a&b \\ c& d\end{array}\right)\in K_0\ |\
a-1\equiv d-1\equiv b\equiv c\equiv 0\  ({\rm mod\ } P^m)\}.
\end{equation}
Similarly,
\begin{equation}
\operatorname{Fix}\ {\cal{B}}_m(x_1)=K_{1,m}=w_\varpi^{-1}K_{0,m}w_\varpi.
\end{equation}
For $i=0,1$, we call an irreducible representation $\sigma$ of
$K_i/K_{i,m}\cong G({\cal{O}}_F/\varpi^m{\cal{O}}_F)$
{\bf strongly cuspidal of level $m>1$} if
$\sigma$ does not contain the trivial character
of the subgroup $(N\cap K_{i,m-1-i})/K_{i,m}$ (as in the 1995 draft; see \cite{M1} for
the precise subgroup). In case $m=1$ we call a
{\bf strongly cuspidal representation} of
$K_i/K_{i,1}\cong G({\mathbb{F}})$
any irreducible representation which is cuspidal as a representation of
$G({\mathbb{F}})$. Since $p> 2$, the metaplectic group
$\overline{SL(2,F)}$ splits over $K_0$ and $K_1$ (though the splittings
are different), so a similar definition may be given for
{\bf strongly cuspidal representation} of
$\overline{K_i}/\overline{K_{i,m}}$. The papers \cite{M1},
\cite{M2} show that each irreducible strongly cuspidal representation
of $K_i/K_{i,m}$ ($i=0,1;m\geq 1$) compactly induces to an irreducible
supercuspidal representation of $G(F)$, and similarly for its metaplectic
covers. Moreover, all irreducible supercuspidal representations
of $G(F)$ (or of $\overline{G(F)}$) arise in this way.

Let $pr$ denote the natural projection $pr:\overline{G(F)}\rightarrow G(F)$.

\begin{remyy}[1995 draft; false when $q\equiv3\pmod4$, see Remark~\ref{rem:sc-caveat}] Fix a strongly cuspidal representation $\sigma$
of level $m\geq 1$ of $K_i$, $i=0,1$. Let
$\overline{\sigma}:=\sigma \otimes \varepsilon$ on
$\overline{K_i}\cong K_i\times \mu_2$, where $\varepsilon$ is the
identity character of $\mu_2$, and let
$$
\begin{array}{c}
\overline{\pi}
:=\operatorname{c\text{-}Ind}_{\overline{K_i}}^{\overline{G(F)}}\overline{\sigma},\\
\pi
:=\operatorname{c\text{-}Ind}_{K_i}^{G(F)}\sigma.
\end{array}
$$
We have
\begin{equation}
\chi_{\overline{\pi}}(x)=\zeta\cdot\chi_{\pi}(pr(x)),\label{eq:sc-relation}
\end{equation}
if $x=(pr(x),\zeta )\in \overline{G(F)}$ with $pr(x)$
elliptic regular.
\label{prop:164}
\end{remyy}

\begin{remyy}[1995 draft; depends on Remark~\ref{prop:164}] Suppose that $p>2$. The character relation above
determines a bijective correspondence
$\overline{\pi}\longmapsto \pi$
between the equivalence classes of irreducible genuine supercuspidal
representations of $\overline{SL(2,F)}$, and the equivalence
classes of irreducible supercuspidal representations of $SL(2,F)$.
\end{remyy}

\begin{remark}\label{rem:sc-caveat}
Three points in Remark~\ref{prop:164} and its proof need attention.
First, $\overline{\sigma}$ must be genuine, which is why it is defined as
$\sigma\otimes\varepsilon$, with $\varepsilon$ the identity character of
$\mu_2$ (the function $\sigma\otimes1$ is not genuine). Second, the isomorphism
$\overline{K_i}\cong K_i\times\mu_2$ is given by a splitting
$\kappa_i(k)=(k,s_i(k))$ which, by Lemma~\ref{lem:facts}(b), need not be
$k\mapsto(k,1)$. With respect to it,
$\overline{\sigma}((k,\zeta))=\zeta\, s_i(k)\,\sigma(k)$, whereas the proof
below evaluates $\chi_{\overline{\sigma}}^\circ$ as if $s_i$ were
trivial. Third, the function $(x,\zeta)\mapsto\zeta\,\chi_\pi(x)$ in
(\ref{eq:sc-relation}) is not invariant under conjugation in
$\overline{G(F)}$ when $q\equiv3\pmod4$, so (\ref{eq:sc-relation}) is
false in that case. Indeed, let
$w=\left(\begin{smallmatrix}0&-1\\1&0\end{smallmatrix}\right)$. For
$g=\left(\begin{smallmatrix}a&b\\c&d\end{smallmatrix}\right)\in G(F)$ with
$abc\neq0$, a computation with (\ref{eq:beta}) and (\ref{eq:conj}) gives
\[
(w,1)\,(g,1)\,(w,1)^{-1}=\bigl(wgw^{-1},\,(a,-bc)_2\bigr).
\]
Let $e\in{\cal O}^\times$ and
$h=\left(\begin{smallmatrix}-1+\varpi e&-e\\ \varpi&-1\end{smallmatrix}\right)\in K_0$.
Then $\det h=1$, and $(\tr h)^2-4=-4\varpi e+\varpi^2e^2$ has valuation $1$,
so $h$ is regular elliptic. Here
$(a,-bc)_2=(-1+\varpi e,\varpi e)_2=(-1,\varpi)_2=(-1)^{(q-1)/2}$. If
$q\equiv3\pmod4$, relation (\ref{eq:sc-relation}) and the invariance of
$\chi_{\overline\pi}$ and $\chi_\pi$ under conjugation would give
\[
\chi_\pi(h)=\chi_{\overline\pi}(h,1)=\chi_{\overline\pi}(whw^{-1},-1)
=-\chi_\pi(whw^{-1})=-\chi_\pi(h),
\]
so $\chi_\pi(h)=0$. But the reduction of $h$ modulo $\varpi$ is $-1$ times a
regular unipotent element, so $h$ fixes exactly one neighbor of $x_0$ in
the tree, and at that neighbor the same holds. Hence the fixed point set
of $h$ is a single edge. For $\pi$ induced from a cuspidal representation
$\sigma$ of $K_0/K_{0,1}$, the Frobenius formula then has a single term,
$\chi_\pi(h)=\chi_\sigma(\bar h)$, and cuspidal characters of
$SL(2,{\mathbb F}_q)$ do not vanish at $-1$ times a regular unipotent
element. (The sign is checked by machine for $q=3,7,11$.) When
$q\equiv1\pmod4$, the function $s_v$ of \S\ref{sec:global} is trivial on
the regular elliptic elements of $K_0$, and whether a relation of the form
(\ref{eq:sc-relation}) holds is open. In general the Frobenius formula for
the cover carries the sign $s_i(gxg^{-1})$ in each term, and this sign
varies with $g$. The two statements are not used
in Chapter~\ref{ch:global}. Note also that over a p-adic field the
regular elliptic elements lie in several conjugacy classes of anisotropic
Cartan subgroups (\S\ref{sec:elliptic}). A corrected relation must be
checked on each of them.
\end{remark}

The 1995 argument for Remark~\ref{prop:164} is given in the next subsection,
with the gaps noted in Remark~\ref{rem:sc-caveat}.

\begin{example}
For $m=1$, a strongly cuspidal representation of $K_0/K_{0,1}\cong
SL(2,{\mathbb F}_q)$ is an irreducible cuspidal representation of
$SL(2,{\mathbb F}_q)$, for instance an irreducible constituent of the
restriction of a cuspidal representation of $GL(2,{\mathbb F}_q)$ of
dimension $q-1$. Its compact induction to $G(F)$ is an irreducible
supercuspidal representation of depth zero.
\end{example}

\section{The Frobenius formula}

To compare characters of supercuspidal representations, we express them by the Frobenius formula for compactly induced representations.


First we prove the following metaplectic version of the Frobenius formula.

\begin{proposition} If $\overline{\pi} \in \Pi_{gen}(\overline{G(F)})$
is induced from a finite dimensional representation $\overline{\sigma}$
of $\overline{K_i}$ then
$$
\chi_{\overline{\pi}} (x)=\sum_{\substack{ z \text{ in }\\
\overline{K_i}\setminus\overline{G(F)}/\overline{K_i}}}
\left[\sum_{\substack{ g \text{ in }\\
\overline{K_i}\setminus\overline{K_i}z\overline{K_i}}}
\chi_{\overline{\sigma}}^\circ (gxg^{-1})\right]
$$
converges for each $x\in \overline{G(F)}^{reg}$. (The convergence is not
uniform.)
\label{prop:166}
\end{proposition}

We first prove a preliminary lemma.
\begin{lemma} Let $x\in \overline{G(F)}^{reg}$ and let
$$
\chi_{\overline{\sigma}}^z(x):=\sum_{\substack{ g \text{ in }\\
\overline{K_i}\setminus\overline{K_i}z\overline{K_i}}}
\chi_{\overline{\sigma}}^\circ (gxg^{-1}).
$$
There exists a finite subset $Z$ of representatives $z$ of
$\overline{K_i}\setminus\overline{G(F)}/\overline{K_i}$
and an open neighborhood $U$ of $x$ such that
$\chi_{\overline{\sigma}}^z(y)=0$ for all $y\in U$ and all $z\notin Z$.
\end{lemma}
\begin{proof} Note $\chi_{\overline{\sigma}}^\circ$ is supported
on $\overline{K_i}$ and $K_i$ is the stabilizer in $G(F)$ of the vertex
$x_i$ of the tree ${\cal{B}}$.
It follows that $\chi_{\overline{\sigma}}^\circ (gxg^{-1})\not= 0$
if and only if $gxg^{-1}\in \overline{K_i}$ if and only if
$p(x)$ fixes the point $p(g)^{-1}x_i$ in the tree for $G(F)$.
It follows that
$\chi_{\overline{\sigma}}^z (x)= 0$ unless $p(x)$ is elliptic
in $G(F)$. Let $\Delta (x)=q^{-m}$ for $m\geq 0$. It is known that
(\cite{L4}, \S 5)
\begin{equation}
|\operatorname{Fix}(x)|=
\left\{ \begin{array}{c}
\frac{q^{m+1}+q^m-2}{q-1},\ \ \ \ \text{if $F[x]/F$ is unramified},\\
\frac{2(q^m-1)}{q-1},\ \ \ \text{if $F[x]/F$ is ramified}.
\end{array} \right.
\end{equation}
From this formula, it follows that there is a neighborhood $U'$
of $p(x)$ in $G(F)^{reg}$ and a number $N=N(U')<\infty$
such that, for all $y\in U'$, we have $|\operatorname{Fix}(y)|\leq N$.
In other words, the fixed point formula above implies that the
number of fixed points is locally uniformly bounded.

From this we obtain a finite set $Z$ as desired.
\end{proof}

\begin{proof}[Proof of Proposition~\ref{prop:166}]
The proposition follows immediately from
the lemma exactly as in \cite[p.~213]{Kut}.
\end{proof}

\begin{proof}[1995 argument for Remark~\ref{prop:164}]
The idea is as follows:
We compare the character of the supercuspidal representations
$\overline{\pi}\in \Pi_{gen}(\overline{G(F)})$ and
$\pi\in\Pi(G(F))$ corresponding to a given finite dimensional representation
$\sigma$ of $K_i$, $i=0,1$, using the Frobenius formula.
As the number of terms in the sum for $\chi_\pi$ and in
the sum for $\chi_{\overline{\pi}}$ are the same, by
Proposition \ref{prop:166} it suffices to compare the
characters $\chi_\sigma^\circ$
and $\chi_{\overline{\sigma}}^\circ$ term-by-term.
For $x$ in $\overline{K_0}$,
we want to show that two analogous terms are equal by
conjugating the elliptic element $x$ around by an element in the cover of
$K_0$ or $A(F)\cup \{t,t^{-1}\}$ or $N(F)$ (in that order), where
$t=\operatorname{diag}(1,\varpi)$ is a diagonal matrix in $GL(2,F)$.
Note we may pass between $K_0$ and $K_1$ by conjugating by either
$t$ or its inverse.

If $T:=$ Cent$(x,\overline{G(F)})$ then $T$ is elliptic and we may embed $T$ into a maximal
compact subgroup $\overline{K}$ of $\overline{G(F)}$.
To effect the term-by-term comparison, we recall that $\overline{K}_0$
splits, so we must first worry about conjugation by
an element $(r,\zeta)$ with $r$ in $A(F)\cup \{t,t^{-1}\}$.
We first calculate
$$
(r,\zeta)\cdot (k,\zeta')\cdot ( r,\zeta)^{-1}
=(r,\zeta)\cdot (k,\zeta')\cdot (r^{-1},(r_1,r_2)_2\zeta^{-1} ),
$$
for $k\in K_0$ elliptic
and $r=\left(\begin{array}{cc} r_1&0\\0&r_2\end{array}\right)\in
GL(2,F)$. The details of this straightforward calculation are
included for the
reader's convenience. Writing
$k=\left(\begin{array}{cc} a&b\\c&d\end{array}\right)$ and
recalling the definition of $p:GL(2)\rightarrow SL(2)$ from \S\ref{sec:local}, we have
\begin{equation}
\begin{array}{c}
rkr^{-1}=\left(\begin{array}{cc} r_1a&r_1b\\r_2c&r_2d\end{array}\right)\cdot
\left(\begin{array}{cc} r_1^{-1}&0\\0&r_2^{-1}\end{array}\right)
=\left(\begin{array}{cc} a&r_1r_2^{-1}b\\r_1^{-1}r_2c&d\end{array}\right)\\
p(r)k=
\left(\begin{array}{cc} r_1a&r_1b\\r_1^{-1}c&r_1^{-1}d\end{array}\right)\\
p(rk)=
\left(\begin{array}{cc} r_1a&r_1b\\r_1^{-1}c&r_1^{-1}d\end{array}\right)\\
p(rk)^{d(1,r_1^{-1}r_2^{-1})}=
\left(\begin{array}{cc} r_1a&r_2^{-1}b\\r_2c&r_1^{-1}d\end{array}\right)
\end{array}
\end{equation}
If $k$ is elliptic then $c\not= 0$. We have, by (\ref{eq:mult}), 
$$
(r,\zeta)\cdot (k,\zeta')
=(rk,\beta_1 (r,k)\zeta\zeta' )=(rk,(r_1,c)_2\zeta\zeta' ),
$$
so
$$
(r,\zeta)\cdot (k,\zeta')\cdot ( r,\zeta)^{-1}
=(rkr^{-1}, (r_1,c)_2(r_1,r_2)_2(r_1,r_1)_2^{-2}\beta_1 (rk, r^{-1})\zeta' ).
$$
Since
$\beta_1 (rk, r^{-1})=(r_1^{-1},r_1^{-2}r_2c)_2(r_1^{-1}r_2^{-1},r_1)_2$,
we find that
$$
(r,\zeta)\cdot (k,\zeta')\cdot ( r,\zeta)^{-1}
=(rkr^{-1}, (r_1,r_1r_2)_2\zeta' ).
$$
Since $r\in A(F)\cup \{t,t^{-1}\}$, we obtain the desired comparison
$$
\chi_{\overline{\sigma}}^\circ ((r,\zeta)
\cdot (k,\zeta')\cdot ( r,\zeta)^{-1})
=\zeta'\cdot \chi_{\sigma}^\circ (rkr^{-1}),
$$
in this case.
\par
Next we must conjugate our elliptic regular $x=(k,\zeta')$ by
$(n,\zeta)$, where $n\in N(F)$. We obtain, as the reader may verify,
$(n,\zeta)\cdot (k,\zeta')\cdot (n,\zeta)^{-1}=(nkn^{-1},\zeta')$.
\end{proof}
%% ===== end aabook4.tex =====



%% ===== begin aabook2a.tex =====
\begin{example}
The fixed-point count in the proof of the lemma above is easy to check in small cases.
For $F[x]/F$ unramified and $m=0$ it gives
$(q+1-2)/(q-1)=1$: a regular elliptic element of $K_0$ with regular
elliptic reduction fixes only the vertex $x_0$. For $m=1$ it gives
$(q^2+q-2)/(q-1)=q+2$, the number of vertices at distance at most $1$ from
$x_0$.
\end{example}


\chapter{The trace formulas}\label{ch:tf}

We recall the trace formula in the case of $SL(2)$, as formulated by
Arthur. Then we state the metaplectic analog of the formula which
can be derived analogously. We shall be very brief since there are
already several expository works explaining the ideas needed here
(e.g., \cite{A15}, \cite{MW}).

This appendix reproduces a summary from the 1995 draft. It is not used in
the proofs of Chapter~\ref{ch:global}, and it is incomplete in the
following respects. The spectral side omits the discrete terms that are
not cuspidal: the residual spectrum (the trivial representation of
$SL(2,{\mathbb A})$; the even theta representations of the cover) and the
terms $\frac14\tr\bigl(M(w,\chi)I_B(\chi,f)\bigr)$ for characters with
$\chi^w=\chi$. The dependence on the truncation parameter $T$ is not made
explicit. For precise statements see \cite{A2}, \cite{A3}, \cite{A6}, and,
for the cover, \cite{Li21}.

\section{The trace formula for \texorpdfstring{$SL(2)$}{SL(2)}}

We shall merely write down the terms, with some definitions, from the
work of Arthur \cite{A2}, \cite{A3}, \cite{A7}.

\subsection{The definition of \texorpdfstring{$J_M(u,f)$}{J\_M(u,f)}}

The weighted orbital integrals $J_M$ are the geometric ingredients of the trace formula.


Let $f\in C_c^\infty (G({\mathbb {A}}_S))$ and let
$\gamma \in M({\mathbb {A}}_S)$ have Jordan decomposition
$\gamma =\sigma u$. Suppose that
$M({\mathbb {A}}_S)_\gamma =G({\mathbb {A}}_S)_\gamma $.
Assume that $S$ is a ``sufficiently large'' finite set which contains
the archimedean places. In this case, define
\begin{equation}
J_M(\gamma ,f)=|D(\gamma )|^{1/2}
\int_{M({\mathbb {A}}_S)\backslash G({\mathbb {A}}_S)}
\int_{{\cal O}_\gamma (M({\mathbb {A}}_S))}f(x^{-1}\mu x)v_M(x)d\mu dx,
\end{equation}
where ${\cal O}_\gamma (M({\mathbb {A}}_S))
=\{m^{-1}\gamma m\ |\ m\in M({\mathbb {A}}_S)\}$.
Here $v_M$ is the ``volume'' of the convex hull of the points
$\{H_B(x),H_{B^{-}}(x)\}$ in ${\mathfrak a}$. If $x=nk$ then
$H_B(x)=0$, so $v_M(nk)=|H_{B^{-}}(x)|=|\log\delta (wx)|$.
In any case, this integral can be written as
\[
J_M(\gamma ,f)=|D(\gamma )|^{1/2}
\int_{G({\mathbb {A}}_S)_\gamma \backslash G({\mathbb{A}}_S)}
f(x^{-1}\gamma x)v_M(x)dx.
\]
This is a function on $M({\mathbb {A}}_S)$.

Suppose now $\gamma \in M({\mathbb {A}}_S)$ is arbitrary and $a\in
A_M({\mathbb {A}}_S)$ is chosen so that
$M({\mathbb {A}}_S)_{a\gamma }=G({\mathbb {A}}_S)_{a\gamma }$.
Following \cite{A7}, (8.1)
we define $J_M(\gamma ,f)$ to be the limit at $a=1$
of $|D(a\gamma )|^{1/2}$ times the integral
\begin{equation}
\int_{G({\mathbb {A}}_S)_\sigma \backslash G({\mathbb {A}}_S)}
\int_{M({\mathbb {A}}_S)_\gamma \backslash G({\mathbb {A}}_S)_\sigma }
f(y^{-1}\sigma x^{-1}a u xy) \left( \sum_{L\in {\cal L}(M)}r_M^L(\gamma
,a)v_L(xy)\right) dxdy,
\end{equation}
where $\gamma\in M({\mathbb A}_S)$ is as above and $r_M^L$ is as in \cite{A7}.

\begin{example}
Let $G=SL(2),\;M=A,\;a=\left(
\begin{array}{cc}
t & 0 \\
0 & t^{-1}
\end{array}
\right) $, where $t\neq \pm 1$. In this case,
\[
r_A(\gamma ,a)=1,
\]
and
\[
v_A(x)=\sum_{v\in S}v_{A,v}(x_v),
\]
where $x=\prod_{v\in S}a_v\cdot \left[
\begin{array}{cc}
1 & n_v \\
0 & 1
\end{array}
\right] \cdot k_v,$ for
$a_v\in {A(k_v)}$, $k_v\in {K_v}$
where $v_{A,v}$ is the local weight of \S\ref{sec:weighted-defs}
(\cite{A6}, Intro., \cite{Sho}, chap. 6). For $z\in Z(G(k))$, $J_A(z,f)$ is obtained from the limit as
$a\rightarrow 1$ of
\begin{equation}
\ |D(a)|^{1/2}
\int_{A({\mathbb {A}}_S)\backslash G({\mathbb {A}}_S)}
f(y^{-1} az y)v_A(y)dy
\end{equation}
\cite{Sho}, section 6.1, p. 64, after removing the divergent multiple of
$J_G(az,f)$ that comes from the term $L=G$ of the sum above
(\S\ref{sec:weighted-central}).
\end{example}

\subsection{Definitions}

We recall the objects on the spectral side of the trace formula.


Let ${\cal X}$ denote the equivalence classes of pairs $(M,\rho)$ as in
\cite{A2}, \cite{A3}.
For each standard parabolic $P=MN$ of $G$ defined over $k$
and each $\chi \in {\cal X}$, we define
$L^2(M(k)\backslash M({\mathbb{A}})^1)_\chi $
as follows. If there is no $P_1\in {\cal P}_\chi $
with $P_1\subset P$ then we define
$L^2(M(k)\backslash M({\mathbb{A}})^1)_\chi $ to be $0$.
Otherwise, let
$\Psi (P_1,P,\chi,W )$ denote the space of
all complex-valued functions $\psi :N({\mathbb{A}})M_1(k)\backslash G(%
{\mathbb{A}})\rightarrow {\mathbb C}$, $P_1=N_1M_1\in {\cal P}_\chi $
with $P_1\subset P$, such that

\begin{enumerate}
\item  $\psi (amk)$, as a function of
$a\in A_1(k)\backslash A_1({\mathbb{A}})$ (fixed
$m\in M_1(k)\backslash M_1({\mathbb{A}})^1,\;k\in K$),
is compactly supported,

\item  $\psi (amk)$, as a function of $k\in K_\infty $, belongs to the
vector space defined by the span of the matrix coefficients of some
irreducible representation $W$ of $K_\infty $,

\item  $\psi (amk)$, as a function of $m\in M_1(k)\backslash M_1(\mathbb A)^1$, belongs to
$L^2(M_1(k)\backslash M_1(\mathbb A)^1)_{cusp,\chi }$.
\end{enumerate}

Given $\psi \in \Psi (P_1,P,\chi,W )$, $P_1\subset P$, let
\[
\psi ^{\wedge ,M}(m)=
\sum_{\delta \in P_1(k)\cap M(k)\backslash
M(k)}\psi (\delta m),\;\;\;\;\;m\in M(k)\backslash
M({\mathbb{A}})^1,
\]
and by \cite[Lemma 1]{L1} the function $\psi ^{\wedge ,M}$ lies in
$L^2(M(k)\backslash M({\mathbb{A}})^1)$. We call this a {\bf theta function}.
Let $L^2(M(k)\backslash M({\mathbb{A}})^1)_\chi $ denote
the $L^2$-closure of the span of all the $\psi ^{\wedge ,M}$ with $\psi \in
\Psi (P_1,P,\chi,W )$, some $W$ and $P_1\subset P$ as above.
Lemma 2 of \cite{L1} implies
\begin{equation}
L^2(M(k)\backslash M({\mathbb{A}})^1)=\bigoplus_{\chi \in
{\cal X}}L^2(M(k)\backslash M({\mathbb{A}})^1)_\chi .
\end{equation}

Let $I_P(\pi)$ be the induced representation as in \cite{A2},
\cite{A3}. If $f$ is a smooth function and $\chi\in {\cal X}$,
let $I_P(\pi,f)_\chi$ denote the
$L^2(M(k)\backslash M({\mathbb{A}})^1)_\chi$-isotypic component
of $I_P(\pi,f)$.

\begin{definition}
Let $G=SL(2)$. Define $M^T_P$ as follows:

if $\chi =(G,\rho )$, let $M_G^T(\pi )_\chi =1$;

if $P=B,\chi =(A,\rho )$ unramified, let
\[
\begin{array}{c}
M_B^T(\pi )_\chi \\
=\operatorname{vol}({\mathbb{R}}/{\mathbb{Z}})
\lim_{\varsigma \rightarrow 0}
\left[
\frac{e^{\varsigma (T)}
M(1,\pi )^{-1}M(1,\pi ,\varsigma )}{\varsigma (\alpha^{\vee })}
+\frac{e^{-\varsigma (T)}
M(s,\pi )^{-1}M(s,\pi ,\varsigma )}{-\varsigma (\alpha^{\vee })}
\right]\\
=M(s,\pi)^{-1}
\lim_{\varsigma \rightarrow 0}
{\frac{M(s,\pi)-e^{-2\varsigma (T)}M(s,\pi,\varsigma)}{\varsigma(\alpha^\vee)}} ,
\end{array}
\]
where $\alpha $ denotes the unique positive root.
\end{definition}

\begin{lemma}
For $G=SL(2)$, each $(G,\pi)\in\chi\in{\cal{X}}$ is
unramified. Moreover, $(A,\mu)\in\chi\in{\cal{X}}$ is unramified
if and only if $\mu^w\not=\mu$.
Here ``unramified'' is being used in Arthur's sense of
a trivial Weyl-group stabilizer, not in the local
$K$-spherical sense used in Chapter~\ref{ch:local}.
In the case of the above definition, at a function $\phi$ on $K$,
\[
\begin{array}{c}
(M^T_B(\pi)_\chi \phi)(k)=
M(s,\pi)^{-1}
\lim_{\varsigma \rightarrow 0}
{\frac{M(s,\pi)\phi(k)-e^{-2\varsigma (T)}e^{\varsigma (H(k))}
M(s,\pi_\varsigma)\phi_\varsigma(k)}{\varsigma(\alpha^\vee)}} \\
=M(s,\pi)^{-1}
\lim_{\varsigma \rightarrow 0}
{\frac{M(s,\pi)\phi(k)-e^{-2\varsigma (T)}M(s,\pi_\varsigma)\phi_\varsigma (k)}{\varsigma(\alpha^\vee)}} \\
=M(s,\pi)^{-1}M'(s,\pi)\phi(k)-2 \phi(k) T,
\end{array}
\]
where the prime denotes differentiation with respect to
$\varsigma$. Thus
\[
M_B^T(\pi )_\chi =
M(s,\pi)^{-1}M'(s,\pi)-2\cdot I\cdot T,
\]
where $I$ denotes the identity.
\end{lemma}

\subsection{The trace formula}

The geometric side of the trace formula is the sum of the following terms:

(a) the sum over classes ${\cal{O}}\subset G(k)$
such that ${\cal{O}}\cap Z(G(k))\not= \emptyset$ of
\begin{equation}
J_{\cal{O}}(f)=
\sum_{\gamma\in {\cal O}}
a^G(S,\gamma)J_G(\gamma,f)
+{\frac{1}{2}}\sum_{\gamma\in {\cal O}\cap A(k)}
a^A(S,\gamma)J_A(\gamma,f),
\end{equation}
where $J_M$ is defined above and
the constants $a^M(S,\gamma )$ satisfy, if $\gamma$ is semisimple,
the property
\[
a^M(S,\gamma )=\left\{
\begin{array}{c}
\operatorname{vol}(M(k)_\gamma \backslash M({\mathbb {A}})_\gamma ^1),
\;\;\;\;{\rm if\;}\gamma \;{\rm is}\;
k-{\rm elliptic\;in}\;M(k), \\
0,\;\;\;\;\;\;\;\;\;\;\;\;\;\;\;\;\;\;{\rm otherwise},\;\;\;\;\;
\end{array}
\right.
\]

(b) the sum over elliptic regular classes ${\cal{O}}\subset G(k)$
of
\begin{equation}
J_{\cal{O}}(f)=
\sum_{\gamma\in {\cal O}}
a^G(S,\gamma)J_G(\gamma,f),
\end{equation}
where $J_G$ (an orbital integral) and $a^G$ (a volume)
are as above,

(c) the sum over ``hyperbolic'' regular classes
${\cal{O}}\subset G(k)$
of
\begin{equation}
J_{\cal{O}}(f)=
{\frac{1}{2}}\sum_{\gamma\in {\cal O}\cap A(k)}
a^A(S,\gamma)J_A(\gamma,f),
\end{equation}
where $J_A$ and $a^A$ are as above,

The spectral side of the trace formula is the sum of the following terms:

(d) the sum over classes $\chi\in {\cal{X}}(G)$ of
\begin{equation}
J_\chi(f)=
\tr(\pi(f)),\ \ \ (\pi,G)\in\chi,
\end{equation}

(e) the sum over classes $\chi\in {\cal{X}}-{\cal{X}}(G)$ of
\begin{equation}
J_\chi(f)=
{\frac{1}{2}}\int_{\Pi_u(A(k)\backslash A({\mathbb{A}}))}
\tr(M_B^T(\pi)_\chi
I_B(\pi,f)_\chi)d\pi,
\end{equation}
where $M_B^T(\pi)_\chi$ is (essentially a logarithmic derivative
of an intertwining operator) given by the definition and lemma above.

\section{The trace formula for \texorpdfstring{$\overline{SL(2)}$}{SL(2)-bar}}

The analogs of the definitions above for the 2-fold cover of $SL(2)$
are so similar that we do not repeat them (precise definitions
are all given in \cite{J2}). The companion text \cite{JC} writes this
trace formula out in full, with proofs, and puts it in invariant form.
See \S\ref{sec:cover-supplies}.

The geometric side of the trace formula is the sum of the following terms:

(a) the sum over classes ${\cal{O}}\subset G(k)$
such that ${\cal{O}}\cap Z(G(k))\not= \emptyset$ of
\begin{equation}
J_{\cal{O}}(f)=
\sum_{\gamma\in {\cal O}}
a^G(S,\gamma)J_G(\gamma,f)+{\frac{1}{2}}
\sum_{\gamma\in {\cal O}\cap A(k)}
a^A(S,\gamma)J_A(\gamma,f),
\end{equation}
where $J_M$ and $a^M$
are given by the ``obvious'' analogs of the formulas
in the previous section,

(b) the sum over elliptic regular classes ${\cal{O}}\subset G(k)$
of
\begin{equation}
J_{\cal{O}}(f)=
\sum_{\gamma\in {\cal O}}
a^G(S,\gamma)J_G(\gamma,f),
\end{equation}
where $J_G$ (an orbital integral) and $a^G$ (a volume)
are given by the ``obvious'' analogs of the formulas
in the previous section,

(c) the sum over ``hyperbolic'' regular classes
${\cal{O}}\subset G(k)$
of
\begin{equation}
J_{\cal{O}}(f)=
{\frac{1}{2}}\sum_{\gamma\in {\cal O}\cap A(k)}
a^A(S,\gamma)J_A(\gamma,f),
\end{equation}
where $J_A$ and $a^A$
are given by the ``obvious'' analogs of the formulas
in the previous section,

The spectral side of the trace formula is the sum of the following terms:

(d) the sum over the cuspidal classes $\chi\in {\cal{X}}(G)$ of
\begin{equation}
J_\chi(f)=
\tr(\pi(f)),\ \ \ (\pi,G)\in\chi,
\end{equation}

(e) the sum over classes $\chi\in {\cal{X}}-{\cal{X}}(G)$ of
\begin{equation}
J_\chi(f)=
{\frac{1}{2}}\int_{\Pi_u(A(k)\backslash\overline{A({\mathbb{A}})})}
\tr(M_B(\pi)_\chi I_B(\pi,f)_\chi)d\pi,
\end{equation}
where $M_B(\pi)_\chi =M_B^T(\pi)_\chi|_{T=0}$.
%% ===== end aabook5.tex =====


%% ===== begin aabook6.tex (rewritten in Draft 1) =====

\begin{example}
Only the genuine part of the cuspidal spectrum contributes to the
spectral side for a genuine $\overline f$. Indeed, if $\phi$ is
non-genuine, so that $\phi(x(1,\zeta))=\phi(x)$, then replacing $g$ by
$g(1,\zeta)$ in $\int\overline f(g)\phi(xg)\,dg$ multiplies the integrand by
$\zeta$, and averaging over $\zeta\in\mu_2$ gives
$\overline R(\overline f)\phi=0$.
\end{example}


\chapter*{Glossary}
\addcontentsline{toc}{chapter}{Glossary}

\begin{description}
\item[Anisotropic (elliptic) Cartan subgroup.] A maximal torus of $SL(2)$
over $F$ whose $F$-points form a compact group. It is the norm-one group of
a quadratic extension of $F$, and its regular elements are the regular
elliptic elements.
\item[Character.] The locally integrable function on the regular set that
represents the trace distribution $f\mapsto\tr\pi(f)$ of an admissible
representation.
\item[Complementary series.] Irreducible unitary representations obtained
by inducing a non-unitary character of the torus, for exponents in a
certain interval.
\item[Corresponding pair (matching).] A pair of test functions, one on
$SL(2)$ and one genuine on the cover, whose orbital integrals are related by
a prescribed rule, so that their traces agree on related representations.
\item[Cuspidal automorphic representation.] An irreducible representation
of the adelic group occurring in the space of cusp forms, the square
integrable functions on $G(k)\backslash G({\mathbb A})$ with vanishing constant
terms.
\item[Discrete series.] Irreducible representations whose matrix
coefficients are square integrable. A limit of discrete series is a
tempered representation at the boundary of the discrete series parameters.
\item[Direct comparison.] The comparison of $SL(2)$ with its cover in
which an element of $SL(2)$ is compared with a lift of the same element.
The odd matchings are of this kind. Two functions match directly if
their orbital integrals differ by the factor $\Delta_\psi$
(Definition~\ref{def:direct}).
\item[Dyadic place.] A place of $k$ lying over $2$.
\item[Elementary theta series.] Theta series on the metaplectic cover
built from the Weil representation of a one-dimensional quadratic form.
\item[Even and odd matching.] Two ways of relating a function on $SL(2)$
to a genuine function on the cover. An odd matching compares orbital
integrals at an element and at a lift of the same element. An even
matching compares the orbital integral at an element with the one at the
square of a lift of it.
\item[Genuine.] A function on the cover, or a representation of it, on
which the central element $(1,-1)$ acts by $-1$.
\item[Hilbert symbol.] The pairing $(x,y)_2=\pm1$ on $F^\times$ that is $1$
exactly when $x$ is a norm from $F(\sqrt y)$. It defines the Kubota
cocycle.
\item[Hyperspecial maximal compact subgroup.] For $SL(2)$ over a p-adic
field, a conjugate of $SL(2,{\cal O})$ under $GL(2,F)$. There are two
$SL(2,F)$-conjugacy classes.
\item[Infinitesimal character.] The character by which the center of the
universal enveloping algebra acts on an irreducible representation of a
real or complex group.
\item[Iwahori subgroup.] The subgroup of $SL(2,{\cal O})$ of matrices that are
upper triangular modulo the maximal ideal.
\item[Kubota cocycle.] The explicit 2-cocycle with values in $\{\pm1\}$,
given by Hilbert symbols of matrix entries, that defines the metaplectic
cover.
\item[Metaplectic cover.] The nontrivial central extension of $SL(2,F)$, or
of $SL(2,{\mathbb A})$, by $\{\pm1\}$.
\item[Multiplicity one.] The property that each irreducible representation
occurs at most once in the cuspidal spectrum.
\item[Near equivalence.] Two automorphic representations are nearly
equivalent if their local components agree at almost all places.
\item[Orbital integral.] The integral of a test function over a conjugacy
class. It is normalized by a power of the Weyl discriminant. A weighted
orbital integral inserts a weight function, as in Arthur's trace formula.
\item[Principal series.] Representations induced from a character of the
Borel subgroup, or, on the cover, from a genuine character of the cover of
the torus.
\item[Satake isomorphism.] The isomorphism between a spherical Hecke
algebra and the Weyl-invariant functions on the torus, given by integrating
along the unipotent radical.
\item[Spherical Hecke algebra.] The convolution algebra of compactly
supported functions that are bi-invariant under a hyperspecial maximal
compact subgroup (on the cover, under its image in a splitting).
\item[Splitting.] A homomorphism from a subgroup into the cover that lifts
the inclusion. Its existence means the cover is trivial over that subgroup.
\item[Strong multiplicity one.] The property that a cuspidal representation
is determined by its local components at almost all places. It fails for
$SL(2)$ and for its cover.
\item[Supercuspidal representation.] An irreducible representation whose
matrix coefficients are compactly supported modulo the center. For $SL(2)$
these are compactly induced from maximal compact subgroups.
\item[Trace formula.] An identity between a sum of orbital integrals and a
sum of traces of automorphic representations, valid for all test
functions.
\item[Stable conjugacy.] Two elements of $SL(2,F)$ are stably
conjugate if an element $h$ of $GL(2,F)$ conjugates one to the other. On
the cover the action of $h$ is twisted by a sign, the Hilbert symbol of
$\det h$ and $\det(1+\gamma)$. With the twist, the two lifts
of an element are never stably conjugate (\S\ref{sec:cover-factor}). The
stable orbital integral is the sum of the orbital integrals over the
conjugacy classes in a stable class. For an elliptic class there are
two of them, and their difference is the $\kappa$-orbital integral.
\item[Transfer factor.] A function by which orbital integrals on two
groups are multiplied before they are compared. For the cover it is the
genuine function $\Delta_\psi$ of Lemma~\ref{lem:Delta}.
\item[Transfer.] A map from representations of the cover to
representations of $SL(2)$ under which the traces of corresponding pairs
agree.
\item[Weil representation.] The oscillator representation of the
metaplectic cover on functions on $F$. It splits into an even and an odd
part.
\item[Weyl integration formula.] The expression of the integral of a
function over the group as a sum, over Cartan subgroups, of integrals of
its orbital integrals.
\end{description}

\chapter*{Acknowledgments}
\addcontentsline{toc}{chapter}{Acknowledgments}

The first author thanks C. Moen and J. Adams for many helpful conversations,
and some anonymous referees of the 1995 draft for their comments.
One of them prodded him to treat the case $p=2$ separately.

This 2026 revision was prepared with the help of AI systems.
Carbon (Google Gemini) refereed the 1995 draft. Claude (Anthropic)
assessed those referee reports, merged the original source files into a
single file, made the corrections, and rewrote the Introduction,
\S\ref{sec:dyadic}, and Chapter~\ref{ch:global} so that the main result
is stated conditionally (Draft~1). The author made the decision to state
the result conditionally. A referee report on Draft~1, which the author
believes was written by ChatGPT (OpenAI), found further errors, among them
a false finiteness lemma in Chapter~\ref{ch:global}. Claude assessed that
report and made the Draft~2 revisions, including the
infinitesimal-character refinement of Hypothesis~\ref{hyp:GH},
Lemma~\ref{lem:A00}, and the corrections to the local character formulas
in Chapter~\ref{ch:local}. Claude also made the Draft~3 revision, which
brings the text into line with the author's style guide without changing
the mathematics. A referee report on Draft~3 by Claude found that the
hypotheses of the main theorem of Draft~3 are inconsistent, together with
further errors. Claude made the Draft~4 revision, which states the main
result as Theorem~\ref{thm:nonexist}, restricts Hypothesis~\ref{hyp:LH},
corrects the errors found, and moves the supercuspidal and trace formula
chapters to appendices. The statement and proof of
Theorem~\ref{thm:nonexist} are due to Claude and have not been checked by a
specialist. In Draft~5 the author supplied the motivating question of
\S\ref{sec:motivation} and the new title, and Claude added
\S\ref{sec:motivation} and \S\ref{sec:where-fails}. At the author's request,
Claude is listed as second author from Draft~5 on, in view of its
contributions to Chapter~\ref{ch:global}. Draft~6 incorporated
corrections and editorial suggestions from a subsequent referee
report by Carbon (Google Gemini), including the real-place
$GL(2)$ cocycle, the real Gaussian phase, and the Weil kernel.
ChatGPT (OpenAI) prepared the Draft~6 source, checked the
internal algebra of these changes, and verified several
bibliographic details against public records.
Where the formal Weil-character computation still lacks
a proof, Draft~6 labels it as conditional rather than
asserting the former formulas. In Draft~7 the author asked how this text
is related to the companion text \cite{JC}. Claude compared the two,
wrote Chapter~\ref{ch:cover} and \S\ref{sec:running}, and added the
references to \cite{JC} in the earlier chapters. Chapter~\ref{ch:cover}
has not been refereed.

AI review, including the referee reports described above, is not an
independent certification of the results. Further details of the AI usage,
the list of changes, and the sources that were not consulted directly
during the revision are recorded in the revision log
(\texttt{revision\_log.md} in the directory \texttt{mult-one} of the
repository \texttt{github.com/wdjoyner/trace-formulas}).

%% ===== begin aabookrf.tex =====
\begin{thebibliography}{99}

\bibitem[Ad1]{Ad1} J. Adams, ``Character of the oscillator representation'',
Israel J. Math. \textbf{98} (1997) 229--252

\bibitem[Ad2]{Ad2} -------,
``Lifting of characters on orthogonal and metaplectic groups'',
Duke Math. J. \textbf{92} (1998) 129--178

%\bibitem[Ad3]{Ad3} -------, Notes on the Flicker correspondence, 1995

\bibitem[Ar]{Ar} Ariturk, ``On the composition series of principal series
representations of a three-fold covering group of $SL(2,K)$'', Nagoya Math
J. \textbf{77}(1980)177-196

%\bibitem[A1]{A1} J. Arthur,
%``The Selberg trace formula for $F$-rank one groups'', Ann.
%Math. \textbf{100}(1974)326-385

\bibitem[A1]{A2} J. Arthur,
``A trace formula for reductive groups, I : terms associated to
classes in $G(\mathbb{Q})$'', Duke Math. J. \textbf{45} (1978)911-952

\bibitem[A2]{A3}  ------,
``A trace formula for reductive groups, II : applications of a
truncation operator'', Comp. Math. \textbf{40} (1980)87-121

%\bibitem[A4]{A4} ------,``Eisenstein series and the trace formula'', in
%\emph{Automorphic forms, }
%\emph{Representations, } \emph{and
%L-functions, vol. I}, 1979 Corvallis Conference, ed. A. Borel, W. Casselman,
%Proc. Symp. Pure Math. vol 33, AMS, 1979

%\bibitem[A5]{A5} ------,
%``The trace formula for non-compact quotient'', Proc. Inter.
%Cong. Math., Warsaw, Poland, 1983

\bibitem[A3]{A6} ------, ``The trace formula in invariant form'', Ann. Math.
\textbf{114} (1981), 1--74

\bibitem[A4]{A7} ------, ``The local behaviour of weighted orbital integrals'',
Duke Math. J. \textbf{56}(1988)223-293

%\bibitem[A8]{A8} ------, ``A measure on the unipotent variety'', Can. J. Math.
%\textbf{37}(1985)1237-1274

%\bibitem[A9]{A9} ------, ``The invariant trace formula, I. Local theory'',
%J.A.M.S.\textbf{1}(1988)323-383

%\bibitem[A10]{A10} ------, ``The invariant trace formula, II. Global theory'',
%J.A.M.S.\textbf{1}(1988)501-554

%\bibitem[A11]{A11} ------, ``Intertwining operators and residues, I: weighted
%characters'', J. Func. Anal.\textbf{84}(1989)19-84

%\bibitem[A12]{A12} ------, ``A family of distributions obtained from
%orbits'', Can. J. Math. \textbf{38}(1986)179-214

%\bibitem[A13]{A13} ------, ``A family of distributions obtained from
%Eisenstein series, I'', preprint

%\bibitem[A14]{A14} ------, ``A family of distributions obtained from
%Eisenstein series, II,'' preprint

\bibitem[A5]{A15} ------,``The trace formula for reductive groups,'' in
\emph{Journ\'ees Automorphes}, Publ. Math. Univ. Paris VII,
U.E.R. Math, 1983

%\bibitem[A16]{A16} ------,``On the Fourier transform of weighted
%orbital integrals,'' preprint

\bibitem[APW]{APW} J. Arthur, ``A Paley--Wiener theorem for real reductive
groups'', Acta Math. \textbf{150} (1983) 1--89

\bibitem[AC]{AC} ----- and L. Clozel, \emph{Simple algebras,
base change, and the advanced} \emph{theory of the trace formula}
Annals of Math. Studies, \textbf{120}, Princeton Univ. Press

%\bibitem[As]{As} M. Assem,
%``The Fourier transform and some character formulae for p-adic
%$SL(\ell )$, $\ell $ a prime'', preprint

%\bibitem[Ba]{Ba} G. Bachman,
%\emph{Elements} \emph{of} \emph{abstract}
%\emph{harmonic} \emph{analysis}, Academic Press, New York, 1965

%\bibitem[B]{B} J. Bernstein, ``All reductive p-adic groups are tame'',
%Funct. Anal. Appl. \textbf{8}(1974)91-93

%\bibitem[BD]{BD} ----, ``Le centre de Bernstein'', (notes by P. Deligne),
%in \emph{Repr\'esentations} \emph{des}
%\emph{groupes} \emph{r\'eductifs}
%\emph{sur} \emph{un} \emph{corp} \emph{local},
%Hermann, Paris, 1984

%\bibitem[BDK]{BDK} J. Bernstein, P. Deligne, D. Kazhdan, ``Trace
%Paley-Wiener theorem for
%reductive p-adic groups'', J. d'Anal. Math. \textbf{47}(1986)180-192

%\bibitem[BZ1]{BZ1} J. Bernstein, A. Zelevinski, ``Representations of
%the group $GL(n,F)$,
%where $F$ is a non-archimedean local field'', Russian Math Surveys
%\textbf{31}(1976)1-68

%\bibitem[BZ2]{BZ2} ----, ``Induced representations of reductive p-adic
%groups, I'', Ann. scient. \'Ec. norm. Sup. \textbf{10}(1977)441-472

\bibitem[Blas]{Blas} D. Blasius, ``On multiplicities for $SL(n)$'',
Israel J. Math. \textbf{88}(1994)237-251

\bibitem[Bl]{Bl} C. Blondel, ``Les repr\'esentations supercuspidales
des groupes
m\'etaplectiques sur GL(2) et leurs caract\'eres'',
Bull de la Soc. Math. France, Memoir 18, 1985

%\bibitem[Bou]{Bou} A. Bouaziz, ``Int\'egrales orbitales sur les groupes
%de Lie r\'eductifs'', Ann. Sc. \'Ec. Norm. Sup. \textbf{27}(1994)573-609

\bibitem[Bo]{Bo} A. Borel, et al, \emph{Seminar in algebraic groups and
related finite groups}, Springer LNM \textbf{131}, 1970

\bibitem[BJ]{BJ} A. Borel and H. Jacquet, ``Automorphic forms and automorphic
representations'', in \emph{Automorphic} \emph{forms},
\emph{representations}, \emph{and} \emph{L-functions},
vol I, ed. A. Borel, W. Casselman, Proc. Symp. Pure Math. 33, AMS,
Providence, RI, 1979

\bibitem[CZ]{CZ} J. Chai and Q. Zhang, ``A strong multiplicity one theorem
for $SL_2$'', Pacific J. Math. \textbf{285} (2016) 345--374


\bibitem[Car]{C} P. Cartier, ``Representations of p-adic groups: a survey'',
in \emph{Automorphic} \emph{forms},
\emph{representations}, \emph{and} \emph{L-functions},
vol I, ed. A. Borel, W. Casselman, AMS, Providence, RI, 1979

%\bibitem[Ca]{Ca} W. Casselman, ``Introduction to the theory of admissible
%representations on reductive p-adic groups'', preprint

%\bibitem[Cl]{Cl} L. Clozel, ``Invariant harmonic analysis
%on the Schwartz space of a
%reductive p-adic group'', in \emph{Harmonic} \emph{analysis} %
%\emph{on} \emph{reductive} \emph{groups}, ed. W. Barker, P.
%Sally, Birkh\"auser, Boston, 1991

%\bibitem[CD]{CD} ------ and P. Delorme, ``Le th\'eor\`eme de Paley-Wiener
%invariant pour les groupes de Lie reductif, II'',
%Ann. Scient. Ec. Norm. Sup., \textbf{23}(1990)193-228

%\bibitem[D]{D} J. Dixmier, \emph{$C^{*}$-algebras}, North-Holland,
%Amsterdam, 1977

%\bibitem[Fe]{Fe} J. Fell, ``$C^{*}$-algebras with smooth dual'', Ill. J.
%Math. \textbf{4}(1960)221-230

%\bibitem[Fla]{F} D. Flath, ``Decomposition of representations into tensor
%products'', in \emph{Automorphic forms, }
%\emph{Representations, } \emph{and
%L-functions}, 1979 Corvallis Conference, ed. A. Borel, W. Casselman,
%Proc. Symp. Pure Math. vol 33, AMS, 1979

\bibitem[Fl]{Fl} Y. Flicker, ``Automorphic forms on covering groups of
$GL(2)$'', Inv
Math \textbf{57}(1980)119-182

%\bibitem[Fr]{Fr} S. Franklin, ``The reducible principal series of
%SL(2) over a p-adic field'', Thesis, Univ. Chicago, 1971

\bibitem[GS2]{GS12} W. T. Gan and G. Savin, ``Representations of metaplectic
groups II: Hecke algebra correspondences'', Represent. Theory
\textbf{16} (2012) 513--539

\bibitem[Ge1]{Ge1} S. Gelbart, \emph{Automorphic forms on adele groups},
Annals of
Math. Studies, Princeton Univ. Press, Princeton, 1975

\bibitem[Ge2]{Ge2} S. Gelbart, \emph{Weil's representation and the
spectrum of the metaplectic} \emph{group},
SLN Math \textbf{530}, 1976

\bibitem[GJ]{GJ} S. Gelbart and H. Jacquet,
``Forms of $GL(2)$ from the analytic
point of view'', in \emph{Automorphic forms},
\emph{representations}, and \emph{L-functions},
(ed. A. Borel, W. Casselman), Proc. Symp. Pure Math, vol 33, AMS, 1979

%\bibitem[GeKn]{GeKn} ------ and T. Knapp, ``L-indistinguishability and R
%groups for the special linear group'', Adv. Math. 43(1982)101-121

\bibitem[GJ2]{GJ78} S. Gelbart and H. Jacquet, ``A relation between
automorphic representations of $GL(2)$ and $GL(3)$'', Ann. Sci. \'Ec.
Norm. Sup. \textbf{11} (1978) 471--542

\bibitem[GePS]{GePS} S. Gelbart and I. Piatetski-Shapiro, ``Distinguished
representations and modular forms of half-integral weight'', Invent.
Math. \textbf{59} (1980) 145--188

\bibitem[GPS]{GPS} S. Gelbart and I. Piatetski-Shapiro, ``Some remarks
on metaplectic cusp
forms and the correspondences of Shimura and Waldspurger'', Israel
J. Math.
\textbf{44}(1983)97-126

\bibitem[GS]{GS} S. Gelbart and P. Sally, ``Intertwining operators and
automorphic forms on the metaplectic group'', Proc. Nat. Acad. Sci.
\textbf{72}(1975)1406-1410

%\bibitem[GG]{GG} I. M. Gelfand, M. I. Graev,
%``Representations of a group of the second
%order with elements from a locally compact field'', Russian Math. Surveys,
%18(1963)29-100

%\bibitem[GGPS]{GGPS} ------ and I. Piatetski-Shapiro,
%\emph{Representation
%theory} \emph{and} \emph{automorphic functions}, W. B. Saunders,
%Philadelphia, 1969

%\bibitem[Go1]{Go1} R. Godement, ``Domaines fondamentaux des groupes
%arithm\'etiques'', Sem. Bourbaki, no. 257, 1962/63

%\bibitem[Go2]{Go2} ------, ``The spectral decomposition of cusp forms'', in
%\emph{Algebraic} \emph{groups}
%\emph{and} \emph{discontinuous}
%\emph{subgroups}, Proc Symp Pure Math, vol. 9, AMS, Providence, 1966

%\bibitem[Go3]{Go3} ------, ``Analyse spectrale des functions modulaires'',
%Sem Bourbaki, no 278, 1964/65

\bibitem[HC1]{HC} Harish-Chandra, \emph{Harmonic analysis on
reductive p-adic
groups}, (notes by G. van Dijk), SLN Math \textbf{162}, 1970

\bibitem[HC2]{HC2} -------, \emph{Collected Works of
Harish-Chandra, vol III}, (ed. V. Varadarajan et al), Springer-Verlag

\bibitem[HC3]{HC68} -------, \emph{Automorphic forms on semisimple
Lie groups}, (notes by J. G. M. Mars), SLN Math \textbf{62}, 1968

%\bibitem[H1]{H1} L. H\"ormander,
%\emph{An} \emph{introduction} \emph{to}
%\emph{complex} \emph{analysis} \emph{in}
%\emph{several} \emph{complex} \emph{variables},
%2nd edition, North-Holland, Amsterdam, 1973

%\bibitem[H2]{H2} -----, \emph{The} \emph{analysis}
%\emph{of} \emph{linear} \emph{partial}
%\emph{differential} \emph{operators}, \emph{I},
%Springer-Verlag, New York, 1983

%\bibitem[HT]{HT} R. Howe, E. C. Tan,
%\emph{Non-abelian harmonic analysis}, Springer-Verlag, 1992

\bibitem[H]{H} R. Howe, ``Character of the oscillator representation'',
(appendix to ``Elements of Oscillator Representations'',
unpublished notes)

\bibitem[HM]{HM} R. Howe (with the collaboration of A. Moy),
\emph{Harish-Chandra}
 
\emph{homomorphisms} \emph{for}
\emph{p-adic} \emph{groups}, CBMS \textbf{59}, AMS (1985).

%\bibitem[Hu]{Hu} J. E. Humphreys, \emph{Arithmetic groups}, Springer
%Lecture Notes in Math., vol. \textbf{789} (1980)

\bibitem[JL]{JL} H. Jacquet and R. Langlands,
\emph{Automorphic} \emph{forms}
\emph{on} \emph{$GL(2)$}, Springer Lecture Notes in Math.,
vol 114, 1970

\bibitem[JS]{JS} H. Jacquet and J. Shalika, ``On Euler products and the
classification of automorphic representations I, II'', Amer. J. Math.
\textbf{103} (1981) 499--558, 777--815

\bibitem[J1]{J} D. Joyner, ``The metaplectic analog of Kazhdan's
endoscopic lift'', Israel J. Math. \textbf{66}(1988)113-154

\bibitem[J2]{J2} ------, ``An application of Arthur's invariant
trace formula to $SL(2)$'', preprint, Feb/1996

\bibitem[JC]{JC} D. Joyner and Claude, \emph{A hitch-hiker's guide to
the metaplectic trace formula for $SL(2)$: basic theory}, Draft~3,
November 2026. Corrected and expanded version of notes of the first
author from 1995. Directory \texttt{sl2bar-trace} of the repository
\texttt{github.com/wdjoyner/trace-formulas}.

\bibitem[J3]{J3} ------, ``A correspondence for the generalized Hecke
algebra of the metaplectic cover $\overline{SL(2,F)}$, $F$ p-adic'',
New York J. Math. \textbf{4} (1998); revised version in preparation (2026)

%\bibitem[K]{K} D. Kazhdan,
%``Cuspidal geometry of p-adic groups'', J. d'Anal. Math.
%47(1986)1-36

%\bibitem[KP]{KP} ------ and  S. Patterson, ``Metaplectic forms'',
%Publ. Math. IHES \textbf{59}(1984)35-142

%\bibitem[KP2]{KP2} ------, ``A Shimura correspondence'',
%Adv. Math. \emph{}(1986)

%\bibitem[Ko1]{Ko1} R. Kottwitz, ``Stable trace formula: elliptic singular
%terms'', Math. Ann. \textbf{275}(1986)365-399

%\bibitem[Ko2]{Ko2} ------, ``Stable trace formula: cuspidal tempered terms'',
%Duke Math. J. \textbf{51}(1984)611-650

\bibitem[K1]{K1} T. Kubota, ``Topological covering of $SL(2)$ over a
local field'', J.
Math. Soc. Japan \textbf{19} (1967)114-121

\bibitem[K2]{K2} ------, \emph{Automorphic functions and the reciprocity
law in a number field}, Kyoto U., 1969

%\bibitem[K3]{K3} ------, \emph{Eisenstein series for $SL(2,\mathbb{R})$},
%John Wiley and Sons, 1973

\bibitem[Kut]{Kut} P. Kutzko,
``Character formulas for supercuspidal representations of
$GL(\ell)$, $\ell$ a prime'', Amer J Math \textbf{109}(1987)201-222

\bibitem[Lab]{Lab} J.-P. Labesse, ``L-indistinguishable representations
and the trace formula for $SL(2)$'', in
\emph{Lie groups and their representations}, John Wiley, 1975

\bibitem[LL]{LL} ------- and R. P. Langlands,
``L-indistinguishability for
$SL(2)$'', Can. J. Math. \textbf{31}(1979) 726--785

%\bibitem[La]{La} S. Lang, \emph{$SL(2,{\mathbb{R}})$}, Addison-Wesley, 1975

\bibitem[L1]{L1} R. Langlands, ``Eisenstein series'', Proc Symp Pure Math,
vol. 9, AMS, Providence, 1966

\bibitem[L2]{L2} ------, \emph{On} \emph{the} \emph{functional}
\emph{equations} \emph{satisfied}
\emph{by} \emph{Eisenstein}
\emph{series}, Springer Lecture Notes, \textbf{544} (1976)

%\bibitem[L3]{L3} ------, \emph{Les d\'ebuts d'une formule des traces
%stable}, Publ. Math. de L'Univ. Paris, VII, 1983

\bibitem[L3]{L4} ------, \emph{Base Change for $GL(2)$},
Annals of Mathematics Studies \textbf{96},
Princeton University Press, 1980

%\bibitem[L5]{L5} ------, ``On the notion of an automorphic
%representation'', in \emph{Automorphic} \emph{forms},
%\emph{representations}, \emph{and} \emph{L-functions},
%vol I, ed. A. Borel, W. Casselman, AMS, Providence, RI, 1979

%\bibitem[Li]{Li}
%R. Lipsman, \emph{Group representations: A survey of some current
%topics}, Springer LNM \textbf{388}, 1974

%\bibitem[Ma]{Ma} G. Mackey, ``Borel structure in groups and their duals'',
%Trans. AMS \textbf{85}(1957)134-165

%\bibitem[Mac]{Mac} G. Mackey, ``Unitary representations of group extensions, I'',
%Acta Math. \textbf{99}(1958)265-311

\bibitem[Li1]{Li12} W.-W. Li, ``La formule des traces pour les rev\^etements
de groupes r\'eductifs connexes. II. Analyse harmonique locale'', Ann. Sci.
\'Ec. Norm. Sup\'er. \textbf{45} (2012) 787--859

\bibitem[Li2]{Li19} W.-W. Li, ``Spectral transfer for metaplectic groups. I.
Local character relations'', J. Inst. Math. Jussieu \textbf{18} (2019)
25--123

\bibitem[Li3]{Li21} W.-W. Li, ``Stabilization of the trace formula for
metaplectic groups'', Ast\'erisque \textbf{464} (2026)

\bibitem[M1]{M1} D. Manderscheid, ``On the supercuspidal representations of
$SL(2)$ and its 2-fold cover, I'', Math. Ann. \textbf{266}(1984)287-295

\bibitem[M2]{M2} ------, ``On the supercuspidal representations of
$SL(2)$ and its 2-fold cover, II'', Math. Ann.
\textbf{266}(1984)297-305

\bibitem[M3]{M3} ------, ``Supercuspidal representations and the
theta correspondence'', preprint

\bibitem[MW]{MW} C. Moeglin, J.-L. Waldspurger,
\emph{Spectral decomposition and Eisenstein series},
Cambridge Univ Press, 1995

\bibitem[Mo1]{Mo1} C. Moen, ``Intertwining operators for covering groups
of $SL(2)$ over
a local field'', preprint

\bibitem[Mo2]{Mo2} ------, ``Irreducibility of unitary principal series
representations
for covering groups of $SL(2)$'', Pac. J. Math. 135(1988)89-110

\bibitem[Mo3]{Mo3} -----, ``The dual pair $(U(1),U(1))$ over a p-adic field'',
Pac. J. Math. \textbf{158}(1993)

\bibitem[M]{M} C. Moore, ``Group extensions of p-adic linear groups'',
Publ. Math. IHES \textbf{35} (1968)

%\bibitem[O]{O} T. Ono, ``On Tamagawa numbers'', in Proc. Symp. Pure Math.,
%vol. 9, AMS, Providence, 1966

\bibitem[PS]{PS} I. Piatetski-Shapiro, ``Work of Waldspurger'', in Lie groups
representations, II (ed. R. Herb, et al), Springer LNM
\textbf{1041},1983

%\bibitem[PR]{PR} G. Prasad, M. Raganathan,
%``Topological central extensions of
%semi-simple groups over local fields'', Annals of Math.
%\textbf{119}(1984)143-268

%\bibitem[RS]{RS} C. Rader, A. Silberger, ``Some consequences of
%Harish-Chandra's submersion principle'', Proc. AMS \textbf{118}
%(1993)1271-1279

\bibitem[Rao]{Rao} R. Rao, ``On some explicit formulas in the theory
of Weil representation'', Pac. J. Math. \textbf{157}(1993)335-371

%\bibitem[Ro]{Ro} A. Robert, \emph{Introduction} \emph{to}
%\emph{the} \emph{representation} \emph{theory} \emph{of}
%\emph{compact} \emph{and} \emph{locally}
%\emph{compact} \emph{groups},
%London Math. Series, Cambridge Univ. Press, Cambridge, 1983

\bibitem[Rog]{R} J. Rogawski,
\emph{Automorphic representations of unitary groups
in three} \emph{variables}
Annals of Math. Studies, \textbf{123}, Princeton Univ. Press, 1990

%\bibitem[Ru]{Ru} W. Rudin, \emph{Functional} \emph{Analysis},
%McGraw-Hill, New York, 1973

\bibitem[Ra]{Ra} D. Ramakrishnan, ``Modularity of the Rankin-Selberg
$L$-series, and multiplicity one for $SL(2)$'', Ann. of Math.
\textbf{152} (2000) 45--111

\bibitem[Sa]{Sa} P. Sally, ``Analytic continuation of the irreducible unitary
representations of the covering group of $SL(2,{\mathbb{R}})$'', Memoirs AMS
\textbf{69}, 1967

%\bibitem[SS1]{SS1} ------- and J. Shalika,
%``Characters of the discrete series of
%representations of SL(2) over a p-adic field'', Proc. Acad. Sci. USA
%61(1968)1231-1237

%\bibitem[SS2]{SS2} ------, ``The Fourier transform of orbital integrals
%on SL(2) over a p-adic field'', in \emph{Lie} \emph{Group}
%\emph{Representations}, \emph{II},
%ed., R. Herb et al, Lecture Notes in Math. vol 1041,
%Springer-Verlag, New York, 1984

\bibitem[Sav]{Sav} G. Savin, ``A local Shimura correspondence'',
Math Ann. \textbf{280}(1988)185-190
\bibitem[Sch]{Sch} L. Schwartz, \emph{Theorie} \emph{des}
\emph{distributions}, \emph{vol. I}, Hermann, Paris, 1950

%\bibitem[Sen]{Sen} S. Sen, ``Explicit reciprocity laws, II', Crelle's J.,
%\textbf{323}(1981)

%\bibitem[Se]{Se} J.-P. Serre,  \emph{Representations of finite groups},
%Springer-Verlag

%\bibitem[Sh]{Sh} F. Shahidi, ``Local coefficients as Artin factors
%for real groups'', Duke Math J. \textbf{52}(1985)973-1007

\bibitem[Sha]{Sha} F. Shahidi, ``On certain $L$-functions'', Amer. J. Math.
\textbf{103} (1981) 297--355

\bibitem[Shal]{Shal} J. Shalika, ``Representations of the two by two
unimodular group over local fields'', IAS

%\bibitem[Sha]{Sha} S. Shatz, \emph{Profinite groups, arithmetic, and
%geometry}, Princeton Univ. Press, Princeton, 1972

%\bibitem[She]{She} D. Shelstad, ``L-indistinguishability for real groups'',
%Math Ann \textbf{259}(1982)385-430

\bibitem[Shi]{Shi} G. Shimura, ``On modular forms of half integral weight'',
Ann. of Math. \textbf{97} (1973) 440--481

\bibitem[Sho]{Sho} S. Shokranian, \emph{The Selberg-Arthur trace formula},
Springer Lecture Notes in Math., vol. \textbf{1503} (1992)

\bibitem[S]{S} A. Silberger, \emph{Introduction} \emph{to}
\emph{harmonic} \emph{analysis} \emph{on} \emph{reductive}
\emph{p-adic} \emph{groups}, Princeton Univ. Press, Princeton, 1979

%\bibitem[SpSt]{SpSt} T. Springer and R. Steinberg, ``Conjugacy classes'',
%in \emph{Seminar on Algebraic}
%\emph{Groups ...} (by A. Borel, et al),
%Springer LNM \textbf{131}, 2nd printing 1986

%\bibitem[T]{T} M. Tadi\'c, ``The topology of the dual space of a
%reductive group over a
%local field'', Glasnik Matem., \textbf{18}(1983)259-279

%\bibitem[Ta]{Ta} M. Taibleson, \emph{Fourier} \emph{analysis}
%\emph{on} \emph{local} \emph{fields}, Princeton Univ.
%Press, Princeton, 1975

%\bibitem[Va]{Va} V. Varadarajan,
%\emph{An introduction to harmonic analysis on semisimple}
%\emph{Lie groups}
%Cambridge Studies in Adv. in Math \textbf{16}, Cambridge, 1986

%\bibitem[Vig]{V} M.-F. Vign\'eras, ``Caract\'erisation des integrales
%orbitales sur un
%groupe r\'eductif p-adique'', J. Fac. Sci. Univ. Tokyo
%\textbf{28}(1981)945-961

\bibitem[TW]{TW} S. Takeda and A. Wood, ``Hecke algebra correspondences
for the metaplectic group'', Trans. Amer. Math. Soc. \textbf{370}
(2018) 1101--1121

\bibitem[Wa1]{Wa1} J.-L. Waldspurger, ``Correspondance de Shimura'',
J. Math. Pure Appl. \textbf{59}(1980)1-133

\bibitem[Wa2]{Wa2} -------, ``Correspondance de Shimura et quaternions'',
Forum Math. \textbf{3} (1991), 219--307

\bibitem[We]{W1} A. Weil, \emph{Basic Number Theory},
Springer-Verlag, 1975

\bibitem[We3]{We3} A. Weil, ``Sur certains groupes d'op\'erateurs unitaires'',
Acta Math. \textbf{111} (1964) 143--211

\bibitem[We2]{W2} -----, \emph{Adeles and algebraic groups},
Birkhauser, 1982


\bibitem[Zh]{Zh} Q. Zhang, ``Addendum to `A strong multiplicity one theorem for $SL_2$'\,'',
Pacific J. Math. \textbf{292} (2018) 283ff.

\end{thebibliography}
%% ===== end aabookrf.tex =====



\end{document}
